\documentclass[../../main.tex]{subfiles}
\graphicspath{{\subfix{../../image/}}}

\begin{document}

\section{差商与扩张定理}\label{section:ch05-延拓}

\subsection{差商与 $W^{1,p}$}\label{subsubsection:ch05-差商W1p}

\begin{definition}\label{definition:ch05-差商}
设$U\subseteq \mathbb{R}^n$是开集,$u:U\to\mathbb{R}$ 为局部可和函数,即$u\in L_{\text{loc}}^1(U)$,且$V\subset \subset U$。
步长为 $h$ 的第 $i$ 差商定义为
\begin{align*}
D_{i}^{h}u(x)\triangleq \frac{u(x+he_{i})-u(x)}{h}\qquad (i=1,\cdots,n),
\end{align*}
其中 $x\in V$，$h\in\mathbb{R}$，$0<|h|<\operatorname{dist}(V,\partial U)$。
并且$D^{h}u\triangleq(D_{1}^{h}u,\cdots,D_{n}^{h}u)$。
\end{definition}
\begin{remark}
差商来自导数定义中的有限增量比. 对于通常的偏导数, 我们考察
\begin{align*}
\frac{\partial u}{\partial x_i}(x)
=
\lim\limits_{h\to0}
\frac{u(x+he_i)-u(x)}{h},
\end{align*}
前提是这个极限存在. 差商则保留极限之前的表达式, 即先固定一个非零步长 $h$, 定义
\begin{align*}
D_i^h u(x)
=
\frac{u(x+he_i)-u(x)}{h}.
\end{align*}
因此, 差商可以理解为沿第 $i$ 个坐标方向、在尺度 $|h|$ 上的平均变化率.
定义差商时不需要预先知道 $u$ 可微, 这使它能够用于研究尚未确定具有弱导数的函数.

分母中的 $h$ 起到按位移大小归一化的作用. 若 $u$ 在 $x$ 附近可微, 则
\begin{align*}
u(x+he_i)-u(x)
=
h\frac{\partial u}{\partial x_i}(x)+o(|h|),
\end{align*}
所以除以 $h$ 后, 主项恰好是偏导数. 如果只考虑函数值之差而不除以 $h$, 则对于连续函数, 这个差通常直接趋于零, 无法反映导数的大小.

这里 $e_i$ 是第 $i$ 个标准单位向量, 因而 $x+he_i$ 只改变 $x$ 的第 $i$ 个坐标. 将所有坐标方向的差商放在一起, 得到
\begin{align*}
D^h u
\triangleq
(D_1^h u,\cdots,D_n^h u),
\end{align*}
它是梯度的离散近似. 上标 $h$ 表示步长参数, 并不表示求 $h$ 阶导数.
允许 $h$ 取正值或负值, 是为了同时考察两个方向的增量. 当 $h<0$ 时, 分母也带有负号, 所以两种方向的差商在极限存在时都趋于同一个偏导数.

条件 $U$ 为开集, 保证每个内点附近都有一小段可以移动而不离开定义域的空间.
但是, 如果让 $x$ 遍历整个 $U$, 则靠近边界时, 一般无法找到一个对所有 $x$ 都适用的非零步长.
因此, 我们先在一个远离边界的局部区域 $V$ 上讨论差商.

由\subcref{Real Analysis-proposition:点集间的距离的性质}{Real Analysis-proposition:点集间的距离的性质-3}, 有
\begin{align*}
d
\triangleq
\operatorname{dist}(V,\partial U)
=
\operatorname{dist}(\overline{V},\mathbb{R}^n\setminus U)
>0.
\end{align*}
因此, 当 $0<|h|<d$ 时, 对任意 $x\in\overline{V}$ 和 $t\in[0,1]$, 都有
\begin{align*}
|x+the_i-x|
=
t|h|
<d,
\end{align*}
从而
\begin{align*}
x+the_i\in U.
\end{align*}
这不仅保证差商中的两个取值点 $x$ 和 $x+he_i$ 都属于 $U$, 还保证连接它们的整条线段都在 $U$ 内.
严格不等式 $|h|<d$ 留出了正的边界距离, 使平移后的闭包仍然紧含于 $U$.
条件 $h\neq0$ 则是因为差商的分母不能为零.
如果 $U=\mathbb{R}^n$, 约定到空集的距离为 $+\infty$, 此时没有来自边界的步长限制.

要求 $u\in L_{\mathrm{loc}}^1(U)$, 是为了使差商适用于弱导数的积分框架.
差商的点值公式本身对一般函数也可以写出, 但后续需要将差商与测试函数相乘并积分, 因而需要局部可积性.
对于每个固定的允许步长 $h$, 由于 $V$ 及其平移 $V+he_i$ 都相对紧于 $U$, 有
\begin{align*}
\int_V |D_i^h u(x)|\,\mathrm{d}x
&\leqslant
\frac{1}{|h|}
\left(
\int_V |u(x+he_i)|\,\mathrm{d}x
+
\int_V |u(x)|\,\mathrm{d}x
\right)\\
&=
\frac{1}{|h|}
\left(
\int_{V+he_i}|u(y)|\,\mathrm{d}y
+
\int_V |u(x)|\,\mathrm{d}x
\right)
<\infty.
\end{align*}
因此 $D_i^h u\in L^1(V)$.
这里不必要求 $u$ 在整个 $U$ 上可积, 因为计算只涉及 $U$ 内部的相对紧区域.

差商与导数之间最直接的联系来自微积分基本定理.
若 $u\in C^1(U)$, 则由上述线段包含关系, 对任意允许的 $x$ 和 $h$, 有
\begin{align*}
u(x+he_i)-u(x)
=
\int_0^1
\frac{\mathrm{d}}{\mathrm{d}t}u(x+the_i)\,\mathrm{d}t=
h\int_0^1
\frac{\partial u}{\partial x_i}(x+the_i)\,\mathrm{d}t.
\end{align*}
所以
\begin{align*}
D_i^h u(x)
=
\int_0^1
\frac{\partial u}{\partial x_i}(x+the_i)\,\mathrm{d}t.
\end{align*}
这说明, 对光滑函数而言, 差商正是偏导数沿一小段线段的平均值.
当 $h\to0$ 时, 这段线段收缩到 $x$, 差商便趋于该点的偏导数.
这也是差商能够近似导数、并能够由导数估计控制的原因.

对于局部可积函数, 函数值及差商都应按几乎处处的意义理解.
如果改变 $u$ 在某个零测集上的取值, 则对于每个固定的 $h$, 差商只可能在该零测集及其平移上发生变化.
由于平移保持零测性, 差商所确定的可积函数等价类不受代表元选择的影响.现在严格证明差商不受代表元选择的影响.
\begin{proof}
设 $u$ 和 $\widetilde{u}$ 是两个代表元, 并且存在零测集 $N\subseteq U$, 使得
\begin{align*}
u(y)=\widetilde{u}(y),
\quad y\in U\setminus N.
\end{align*}
固定一个允许的非零步长 $h$满足:$h\in\mathbb{R}$，$0<|h|<\operatorname{dist}(V,\partial U)$.
比较两个差商时, 要同时比较两组函数值:
\begin{align*}
D_i^h u(x)
=
\frac{u(x+he_i)-u(x)}{h},\quad 
D_i^h\widetilde{u}(x)
=
\frac{\widetilde{u}(x+he_i)-\widetilde{u}(x)}{h}.
\end{align*}
只有当 $x$ 或 $x+he_i$ 落入 $N$ 时, 这两个差商才可能不同.
其中, 条件 $x+he_i\in N$ 等价于
\begin{align*}
x\in N-he_i.
\end{align*}
因此, 两个差商可能不同的点只会出现在集合
\begin{align*}
E_h\triangleq V\cap\bigl(N\cup(N-he_i)\bigr)
\end{align*}
中.
由\hyperref[Real Analysis-theorem:外测度的平移不变性]{测度的平移不变性}以及有限个零测集之并仍为零测集, 有
\begin{align*}
|E_h|
\leqslant
|N|+|N-he_i|
=0.
\end{align*}
对于每个 $x\in V\setminus E_h$, 两个取值点都不属于 $N$, 所以
\begin{align*}
u(x)=\widetilde{u}(x),\quad
u(x+he_i)=\widetilde{u}(x+he_i),
\end{align*}
从而
\begin{align*}
D_i^h u(x)=D_i^h\widetilde{u}(x).
\end{align*}
这就证明了两个差商在 $V$ 上几乎处处相等, 因而代表同一个可积函数等价类.
\end{proof}
\end{remark}

\begin{proposition}[紧集平移的紧性]\label{proposition:紧集平移与紧参数平移集}
设 $K,A\subseteq\mathbb{R}^n$ 都是紧集, 则
\begin{align*}
K+A
\triangleq
\{x+a:x\in K,\, a\in A\}
=
\bigcup\limits_{a\in A}(K+a)
\end{align*}
也是紧集.
特别地, 有以下结论:
\begin{enumerate}[(1)]
\item 对任意固定的向量 $a\in\mathbb{R}^n$, 平移集
\begin{align*}
K+a\triangleq\{x+a:x\in K\}
\end{align*}
是紧集.
\item 对任意有限个向量 $a_1,\cdots,a_m\in\mathbb{R}^n$, 集合
\begin{align*}
\bigcup\limits_{j=1}^m
\{x+ta_j:x\in K,\, t\in[0,1]\}
\end{align*}
是紧集.
\end{enumerate}
\end{proposition}
\begin{proof}
由于 $K$ 和 $A$ 都是紧集, 它们的乘积 $K\times A$ 是 $\mathbb{R}^{2n}$ 中的紧集.
定义映射
\begin{align*}
F:K\times A\longrightarrow\mathbb{R}^n,\quad
(x,a)\longmapsto x+a.
\end{align*}
对任意 $(x,a),(y,b)\in K\times A$, 有
\begin{align*}
|F(x,a)-F(y,b)|
\leqslant |x-y|+|a-b|\leqslant
\sqrt{2}\sqrt{|x-y|^2+|a-b|^2}.
\end{align*}
因此 $F$ 连续.
\hyperref[Topology-theorem:topology-紧集的连续像是紧集]{紧集在连续映射下的像仍然是紧集}, 故
\begin{align*}
K+A=F(K\times A)
\end{align*}
是紧集.

对于结论(1), 取紧集 $A=\{a\}$, 即可得到 $K+a$ 是紧集.

对于结论(2), 令
\begin{align*}
A_j\triangleq\{ta_j:t\in[0,1]\},
\quad j=1,\cdots,m.
\end{align*}
映射 $t\mapsto ta_j$ 连续, 所以 $A_j$ 是紧区间 $[0,1]$ 的连续像, 从而由\hyperref[Topology-theorem:topology-紧集的连续像是紧集]{紧集在连续映射下的像仍然是紧集}知$A_j$都是紧集.
因此, 有限并
\begin{align*}
A\triangleq\bigcup\limits_{j=1}^m A_j
\end{align*}
也是紧集.
应用已经证明的结论, 可知
\begin{align*}
K+A
=
\bigcup\limits_{j=1}^m
\{x+ta_j:x\in K,\, t\in[0,1]\}
\end{align*}
是紧集.
\end{proof}

\begin{proposition}[开集内部小幅平移后仍在原本开集中]\label{proposition:开集内部小幅平移后仍在原本开集中}
设 $U\subseteq\mathbb{R}^n$ 是开集, $\varnothing\neq V\subset\subset U$, 记
\begin{align*}
\delta\triangleq\operatorname{dist}(V,\partial U)>0.
\end{align*}
当 $U=\mathbb{R}^n$ 时, 约定 $\delta=+\infty$.
则对任意满足 $|a|<\delta$ 的向量 $a\in\mathbb{R}^n$, 都有
\begin{align*}
\{x+ta:x\in\overline{V},\, t\in[0,1]\}
\subset\subset U.
\end{align*}
特别地, $\overline{V}+a\subset\subset U$.

更一般地, 若有限个向量 $a_1,\cdots,a_m$ 都满足 $|a_j|<\delta$, 则
\begin{align*}
\bigcup\limits_{j=1}^m
\{x+ta_j:x\in\overline{V},\, t\in[0,1]\}
\subset\subset U.
\end{align*}
\end{proposition}
\begin{proof}
当 $U=\mathbb{R}^n$ 时, 结论直接由\cref{proposition:紧集平移与紧参数平移集}得到.
以下设 $U\neq\mathbb{R}^n$.

\noindent$\mathbf{Step}\bm{1}$.{\heiti 确定 $\overline{V}$ 到边界的统一正距离.}

定义
\begin{align*}
d(x)\triangleq\operatorname{dist}(x,\partial U),
\quad x\in\mathbb{R}^n.
\end{align*}
由\cref{Real Analysis-theorem:距离函数都是一致连续的}知$d$ 连续.
又因为 $U$ 是开集, $\partial U\cap U=\varnothing$,由\subcref{Real Analysis-proposition:点集间的距离的性质}{Real Analysis-proposition:点集间的距离的性质-3},对每个 $x\in U$ 都有 $d(x)>0$.
因此, $d$ 在紧集 $\overline{V}\subseteq U$ 上取得正的最小值.

由于 $V$ 在 $\overline{V}$ 中稠密, 由 $d$ 的连续性有
\begin{align*}
\delta
=\inf\limits_{x\in V}d(x)
=\min\limits_{x\in\overline{V}}d(x)
>0.
\end{align*}
具体而言, 若 $x_0\in\overline{V}$ 是最小值点, 取 $x_k\in V$ 使 $x_k\to x_0$, 则
\begin{align*}
\inf\limits_{x\in V}d(x)
\leqslant
\lim\limits_{k\to\infty}d(x_k)
=d(x_0)
=\min\limits_{x\in\overline{V}}d(x),
\end{align*}
而反向不等式由$V\subseteq\overline{V}$ 立即得到.
特别地,
\begin{align*}
|x-z|\geqslant\delta,
\quad \forall x\in\overline{V},\, z\in\partial U.
\end{align*}

\noindent$\mathbf{Step}\bm{2}$.{\heiti 证明移动距离小于 $\delta$ 时不会离开 $U$.}

任取 $x\in\overline{V}$ 及满足 $|y-x|<\delta$ 的点 $y\in\mathbb{R}^n$.
下面证明 $y\in U$.
反设 $y\notin U$, 定义
\begin{align*}
\gamma(s)\triangleq x+s(y-x),
\quad s\in[0,1],
\end{align*}
以及
\begin{align*}
A\triangleq\{s\in[0,1]:\gamma(s)\notin U\}.
\end{align*}
由于 $\gamma$ 连续且 $\mathbb{R}^n\setminus U$ 是闭集,再由\cref{Topology-proposition:连续映射下闭集的原像是闭集}知$A=\gamma^{-1}(\mathbb{R}^n\setminus U)$ 是 $[0,1]$ 的闭子集.
又因为 $1\in A$, 所以 $A$ 是非空紧集, 从而存在最小值
\begin{align*}
s_0\triangleq\min A.
\end{align*}
由 $\gamma(0)=x\in U$, 有 $0\notin A$, 因此 $0<s_0\leqslant1$.
根据 $s_0$ 的最小性, 对所有 $0\leqslant s<s_0$ 都有 $\gamma(s)\in U$, 而 $\gamma(s_0)\notin U$.
由$\gamma$的连续性和\subcref{Real Analysis-theorem:边界和开集基本性质}{Real Analysis-theorem:边界和开集基本性质-5}知
\begin{align*}
z\triangleq\gamma(s_0)
=\lim\limits_{s\to s_0}\gamma(s)
\in\overline{U}\setminus U
=\partial U.
\end{align*}
于是
\begin{align*}
\delta
\leqslant|x-z|
=s_0|y-x|
\leqslant|y-x|
<\delta,
\end{align*}
矛盾!
因此 $y\in U$.

\noindent$\mathbf{Step}\bm{3}$.{\heiti 应用于线段平移, 并验证紧包含关系.}

若 $|a|<\delta$, 则对任意 $x\in\overline{V}$ 及 $t\in[0,1]$,
\begin{align*}
|(x+ta)-x|=t|a|\leqslant|a|<\delta.
\end{align*}
由$\mathbf{Step}\bm{2}$可知 $x+ta\in U$, 因而
\begin{align*}
K_a\triangleq
\{x+ta:x\in\overline{V},\, t\in[0,1]\}
\subseteq U.
\end{align*}
由\cref{proposition:紧集平移与紧参数平移集}, $K_a$ 是紧集.
所以 $K_a\subset\subset U$.
同理, $\overline{V}+a$ 是包含于 $U$ 的紧集, 故 $\overline{V}+a\subset\subset U$.

最后, 对有限个满足 $|a_j|<\delta$ 的向量分别应用上述结论, 可知每个 $K_{a_j}$ 都是包含于 $U$ 的紧集.
它们的有限并仍然是紧集且包含于 $U$, 因此
\begin{align*}
\bigcup\limits_{j=1}^m K_{a_j}
=
\bigcup\limits_{j=1}^m
\{x+ta_j:x\in\overline{V},\, t\in[0,1]\}
\subset\subset U.
\end{align*}
\end{proof}

\begin{theorem}[差商的分部积分公式]\label{theorem:ch05-差商的分部积分公式}
设 $U\subseteq\mathbb{R}^n$ 为开集, $i\in\{1,\cdots,n\}$, $h\neq 0$. 设 $u\in L^1_{\mathrm{loc}}(U)$, $v\in C_c^\infty(U)$, 且
\begin{align*}
0<|h|<\frac{1}{2}\operatorname{dist}(\operatorname{spt}v,\partial U),
\end{align*}
则
\begin{align*}
\int_U u(x)D_i^h v(x)\,\mathrm{d}x
=-\int_U D_i^{-h}u(x)v(x)\,\mathrm{d}x.
\end{align*}
\end{theorem}
\begin{proof}
由 $v\in C_c^\infty(U)$ 及 $0<|h|<\frac{1}{2}\operatorname{dist}(\operatorname{spt}v,\partial U)$,再结合\cref{proposition:开集内部小幅平移后仍在原本开集中}可知对任意 $x\in\operatorname{spt}v$, 有 $x+h e_i\in U$ 且 $x-h e_i\in U$. 因此下列积分均有意义. 直接计算得
\begin{align*}
\int_U u(x)D_i^h v(x)\,\mathrm{d}x
=\frac{1}{h}\int_U u(x)\left(v(x+h e_i)-v(x)\right)\mathrm{d}x=\frac{1}{h}\int_U u(x)v(x+h e_i)\,\mathrm{d}x
-\frac{1}{h}\int_U u(x)v(x)\,\mathrm{d}x.
\end{align*}
对第一个积分作变量替换 $y=x+h e_i$, 则 $x=y-h e_i$, $\mathrm{d}x=\mathrm{d}y$. 由于 $v$ 的支集在 $U$ 中, 当 $y\notin U$ 时 $v(y)=0$, 故积分区域仍可取为 $U$, 从而
\begin{align*}
\int_U u(x)v(x+h e_i)\,\mathrm{d}x
=\int_U u(y-h e_i)v(y)\,\mathrm{d}y.
\end{align*}
代回上式, 得
\begin{align*}
\int_U u(x)D_i^h v(x)\,\mathrm{d}x
&=\frac{1}{h}\int_U u(x-h e_i)v(x)\,\mathrm{d}x
-\frac{1}{h}\int_U u(x)v(x)\,\mathrm{d}x\\
&=-\int_U \frac{u(x)-u(x-h e_i)}{h}v(x)\,\mathrm{d}x\\
&=-\int_U D_i^{-h}u(x)v(x)\,\mathrm{d}x.
\end{align*}
\end{proof}

\begin{theorem}[差商与弱导数]\label{theorem:ch05-差商弱导数}
设$U$是$\mathbb{R}^n$中的开集,对每个 $V\subset \subset U$.
\begin{enumerate}[(i)]
\item 设 $1\leqslant p<\infty$，$u\in W^{1,p}(U)$。则对一切 $0<|h|<\frac{1}{2}\operatorname{dist}(V,\partial U)$,都有
\begin{align}
\|D^{h}u\|_{L^{p}(V)}\leqslant n\|Du\|_{L^{p}(U)}
\end{align}

\item 设 $1<p<\infty$，$u\in L^{p}(V)$，且存在常数 $C$ 使
\begin{align}\label{eq::jijiowjiwiw23323rwe}
\|D^{h}u\|_{L^{p}(V)}\leqslant C
\end{align}
对一切 $0<|h|<\frac{1}{2}\operatorname{dist}(V,\partial U)$ 成立。则
\begin{align*}
u\in W^{1,p}(V),\quad
\|Du\|_{L^p(V)}\leqslant C.
\end{align*}
\end{enumerate}
\end{theorem}
\begin{remark}
$p=\infty$的情形见下文\cref{theorem:ch05-W1infty刻画}.断言(ii)对 $p=1$ 不成立。
下面的一维反例说明, 差商在 $L^1(V)$ 中一致有界, 并不足以推出 $u\in W^{1,1}(V)$.

取
\begin{align*}
U=(-2,2),\quad V=(-1,1),
\end{align*}
并定义
\begin{align*}
u(x)\triangleq
\begin{cases}
0,&-2<x\leqslant0,\\
1,&0<x<2.
\end{cases}
\end{align*}
显然 $V\subset \subset U$, 且 $u\in L^1(U)$. 又因为
\begin{align*}
\partial U=\{-2,2\},\quad
\operatorname{dist}(V,\partial U)=1,
\end{align*}
所以题设中的步长条件就是
$0<|h|<\frac{1}{2}.$

先计算正步长的差商. 当 $0<h<\frac{1}{2}$ 时, 只有取值点 $x$ 和 $x+h$ 分别位于跳跃点 $0$ 的两侧, 差商才可能非零. 具体而言, 对几乎处处的 $x\in V$, 有
\begin{align*}
D^h u(x)
=
\frac{u(x+h)-u(x)}{h}
=
\begin{cases}
\dfrac{1}{h},&-h<x<0,\\
0,&x\in V\setminus(-h,0).
\end{cases}
\end{align*}
因此
\begin{align*}
\|D^h u\|_{L^1(V)}
=
\int_{-h}^{0}\frac{1}{h}\,\mathrm{d}x
=1.
\end{align*}

再计算负步长的差商. 当 $-\frac{1}{2}<h<0$ 时,对几乎处处的 $x\in V$, 有
\begin{align*}
D^h u(x)
=
\frac{u(x+h)-u(x)}{h}
=
\frac{u(x)-u(x-(-h))}{-h}=
\begin{cases}
-\frac{1}{h},&0<x<-h,\\
0,&x\in V\setminus(0,-h).
\end{cases}
\end{align*}
从而
\begin{align*}
\|D^h u\|_{L^1(V)}
=
\int_0^{-h}-\frac{1}{h}\,\mathrm{d}x
=1.
\end{align*}
综合正、负步长的计算, 对一切
$0<|h|<\frac{1}{2}\operatorname{dist}(V,\partial U)$, 都有
\begin{align*}
\|D^h u\|_{L^1(V)}=1.
\end{align*}
因此, 题设中的一致有界条件成立, 可以取 $C=1$.

下面证明 $u\notin W^{1,1}(V)$.
反设 $u\in W^{1,1}(V)$, 则存在弱导数 $g\in L^1(V)$, 使得对任意 $\varphi\in C_c^\infty(V)$, 都有
\begin{align*}
\int_V u(x)\varphi'(x)\,\mathrm{d}x
=
-\int_V g(x)\varphi(x)\,\mathrm{d}x.
\end{align*}
另一方面, 由$\text{spt}\,\varphi\subseteq V=(-1,1)$知$\varphi$ 在端点 $1$ 附近恒为零,再根据 $u$ 的定义,有
\begin{align*}
\int_V u(x)\varphi'(x)\,\mathrm{d}x
=
\int_0^1\varphi'(x)\,\mathrm{d}x
=
-\varphi(0).
\end{align*}
因此, 假定存在的弱导数 $g$ 必须满足
\begin{align}\label{NMZNINWhiuhfjknxlzzvz}
\int_V g(x)\varphi(x)\,\mathrm{d}x
=
\varphi(0),
\quad
\forall\varphi\in C_c^\infty(V).
\end{align}

取标准磨光子 $\eta\in C_c^\infty(\mathbb{R}^n)$, 满足
\begin{align*}
\eta\geqslant0,\qquad
\operatorname{spt}\eta\subseteq\overline{B^{0}(0,1)},\qquad
\int_{\mathbb{R}^n}\eta(x)\,\mathrm{d}x=1.
\end{align*}
对$\varepsilon>0$, 定义
\begin{align*}
\eta_\varepsilon(x)
\triangleq\frac{1}{\varepsilon^n}
\eta\left(\frac{x}{\varepsilon}\right).
\end{align*}
于是
\begin{align*}
\eta(0)>0,\quad
0\leqslant\eta(x)\leqslant\eta(0),
\quad x\in\mathbb{R}.
\end{align*}
对 $0<r<1$, 直接取测试函数
\begin{align*}
\varphi_r(x)
\triangleq
\frac{\eta_{\frac{r}{2}}(x)}{\eta_{\frac{r}{2}}(0)}
=
\frac{\eta\left(\frac{2x}{r}\right)}{\eta(0)}.
\end{align*}
由于 $\eta\in C_c^\infty(\mathbb{R})$, 且
$\operatorname{spt}\eta=[-1,1]$, 有
\begin{align*}
\varphi_r\in C^\infty(\mathbb{R}),\quad
\operatorname{spt}\varphi_r
=
\left[-\frac{r}{2},\frac{r}{2}\right]
\subset \subset (-r,r)\subset \subset V.
\end{align*}
因此
\begin{align*}
\varphi_r\in C_c^\infty(V),\quad
0\leqslant\varphi_r\leqslant1,\quad
\varphi_r(0)=1.
\end{align*}
将 $\varphi_r$ 代入\eqref{NMZNINWhiuhfjknxlzzvz}式, 得到
\begin{align}\label{eqNVMZnnjczvnms}
1
=
\left|
\int_V g(x)\varphi_r(x)\,\mathrm{d}x
\right|
\leqslant
\int_{-\frac{r}{2}}^{\frac{r}{2}}|g(x)|\,\mathrm{d}x
\leqslant
\int_{-r}^{r}|g(x)|\,\mathrm{d}x.
\end{align}
但是, 由于 $g\in L^1(V)$, 由\hyperref[Real Analysis-theorem:积分的绝对连续性]{Lebesgue积分的绝对连续性}, 有
\begin{align*}
\lim\limits_{r\to0^+}
\int_{-r}^{r}|g(x)|\,\mathrm{d}x
=0,
\end{align*}
这与上面的不等式\eqref{eqNVMZnnjczvnms}矛盾! 因此不存在这样的 $g\in L^1(V)$, 即
\begin{align*}
u\notin W^{1,1}(V).
\end{align*}

这个反例中, 差商的非零部分位于长度为 $|h|$ 的小区间上, 高度为 $\frac{1}{|h|}$, 因而其 $L^1$ 范数始终等于 $1$.
但是当 $h\to0$ 时, 这些差商的积分质量越来越集中于跳跃点 $0$.
这正是本例中差商虽然在 $L^1(V)$ 中一致有界, 却不能产生 $L^1$ 弱导数的原因.

直观地看, 每个差商 $D^h u$ 都是一个紧贴跳跃点 $0$ 的矩形脉冲: 它的宽度为 $|h|$, 高度为 $\frac{1}{|h|}$, 所以矩形的面积始终为 $1$, 即
\begin{align*}
\int_V|D^h u(x)|\,\mathrm{d}x
=|h|\cdot\frac{1}{|h|}
=1.
\end{align*}
当 $h\to0$ 时, 这个矩形越来越窄, 同时越来越高, 但面积并没有减少. 因此, 对每个固定的 $x\neq0$, 当 $|h|$ 足够小时都有 $D^h u(x)=0$, 可是差商的总积分仍然等于 $1$. 这说明积分质量并没有消失, 而是被压缩到越来越靠近 $0$ 的小区间中.

所谓积分质量集中于 $0$, 具体是指: 无论取多小的邻域 $(-r,r)$, 只要 $0<|h|<\min\{r,\frac{1}{2}\}$, 差商的全部积分质量就已经落在这个邻域内, 即
\begin{align*}
\int_{-r}^{r}|D^h u(x)|\,\mathrm{d}x=1.
\end{align*}
如果$u$存在弱导数$g$,那么$g$的积分性质也应该与$D^hu(x)$类似.
而任何固定的 $L^1$ 函数 $g$ 都不可能具有这种性质. 因为 Lebesgue 积分的绝对连续性保证
\begin{align*}
\lim\limits_{r\to 0^+}\int_{-r}^{r}|g(x)|\,\mathrm{d}x=0.
\end{align*}
也就是说, 对于一个固定的可积函数, 区间缩得足够小时, 其中的积分质量必然趋于零; 本例中的差商却通过不断增大高度, 在不断缩小的区间内始终保留总量为 $1$ 的积分质量.

函数 $u$ 在 $0$ 两侧都是常数, 它的变化全部来自 $0$ 处大小为 $1$ 的跳跃. 前面的测试函数计算说明, 若要用一个 $L^1$ 弱导数表示这个跳跃, 就必须让它在任意小的零点邻域内保留至少为 $1$ 的绝对积分, 这与积分的绝对连续性矛盾！ 因此, 差商的 $L^1$ 范数一致有界只能控制积分质量的总量, 不能阻止这些质量集中到一个点上, 也就不足以保证存在 $L^1$ 弱导数.
\end{remark}
\begin{proof}
\noindent$\mathbf{Step}\bm{1}$.{\heiti 选择包含所有平移线段的中间区域, 并建立光滑函数的局部差商估计.}

先证明断言(i).固定任意一个满足 $0<|h|<\frac{\operatorname{dist}(V,\partial U)}{2}$ 的步长 $h$, 定义
\begin{align*}
K_h\triangleq
\bigcup\limits_{i=1}^n
\left\{
x+the_i:
x\in\overline{V},\, t\in[0,1]
\right\}.
\end{align*}
由\cref{proposition:紧集平移与紧参数平移集}知$K_h$也是紧集.
又由\cref{proposition:开集内部小幅平移后仍在原本开集中}知$K_h\subseteq U$.

由\cref{proposition:紧集的相对紧开邻域}, 可以选择开集 $W$, 使得
\begin{align*}
K_h\subseteq W\subset\subset U.
\end{align*}
特别地,有
\begin{align}\label{eqNJKNDnxjklndklszjsjwzzcxvds}
V+the_i\subseteq W,
\quad t\in[0,1],\quad i=1,\cdots,n.
\end{align}

设 $u\in C^\infty(W)\cap W^{1,p}(W)$.
对任意 $x\in V$ 及 $i=1,\cdots,n$, 线段 $\{x+the_i:t\in[0,1]\}$ 完全位于 $W$ 中, 因而由微积分基本定理有
\begin{align*}
u\left( x+he_i \right) -u\left( x \right) =\int_0^1{\frac{\mathrm{d}}{\mathrm{d}t}u\left( x+the_i \right) \,\mathrm{d}t}=h\int_0^1{u_{x_i}\left( x+the_i \right) \,\mathrm{d}t}\Longrightarrow D^h_i u(x)=\int_0^1{u_{x_i}\left( x+the_i \right) \,\mathrm{d}t}.
\end{align*}
从而再由\hyperref[Basis of Analytics-theorem:离散的Hölder不等式]{离散的Hölder不等式}得
\begin{align*}
|D^h u(x)|^p
&\leqslant \left( \sum_{i=1}^n |D_i^h u(x)| \right)^p
\leqslant n^{p-1} \sum_{i=1}^n |D_i^h u(x)|^p \\
&\leqslant n^{p-1} \sum_{i=1}^n \int_0^1 |u_{x_i}(x + t h e_i)|^p \,\mathrm{d}t \\
&\leqslant n^{p-1} \sum_{i=1}^n \int_0^1 |Du(x + t h e_i)|^p \,\mathrm{d}t.
\end{align*}
两边在 $V$ 上积分并利用\hyperref[Real Analysis-theorem:实变函数-Fubini定理]{Fubini定理}和\eqref{eqNJKNDnxjklndklszjsjwzzcxvds}式,即得
\begin{align*}
\int_V{|D^hu|^p\,\mathrm{d}x}&\leqslant n^{p-1}\sum_{i=1}^n{\int_0^1{\int_V{\left| Du\left( x+the_i \right) \right|^p\,\mathrm{d}x\,\mathrm{d}t}}}=n^{p-1}\sum_{i=1}^n{\int_0^1{\int_{V+the_i}{\left| Du\left( x \right) \right|^p\,\mathrm{d}x\,\mathrm{d}t}}}
\\
&\leqslant n^{p-1}\sum_{i=1}^n{\int_0^1{\int_W{\left| Du\left( x \right) \right|^p\,\mathrm{d}x\,\mathrm{d}t}}}=n^p\int_W{\left| Du\left( x \right) \right|^p\,\mathrm{d}x}.
\end{align*}
故
\begin{align}\label{mzvjkhshtuwwaa}
\|D^h u_m\|_{L^p(V)}\leqslant n\|Du_m\|_{L^p(W)}.
\end{align}

\noindent$\mathbf{Step}\bm{2}$.{\heiti 利用局部光滑逼近将估计推广到 Sobolev 函数.}

仍固定上述 $h$ 和相应的 $W$.
由于 $W\subset\subset U$, 由\hyperref[theorem:ch05-光滑函数的局部逼近定理]{光滑函数的局部逼近定理}, 可以选择一列磨光函数 $u_m$, 使得
\begin{align*}
u_m\in C^\infty(W)\cap W^{1,p}(W),\quad
\lim\limits_{m\to\infty}
\|u_m-u\|_{W^{1,p}(W)}=0.
\end{align*}
于是 $u_m\to u$ 于 $L^p(W)$，且 $Du_m\to Du$ 于 $L^p(W)$。
对每个 $i=1,\cdots,n$, 由 $V\subseteq W$ 及 $V+he_i\subseteq W$, 有
\begin{align*}
\|D_i^h u_m-D_i^h u\|_{L^p(V)}
&\leqslant
\frac{1}{|h|}
\left(
\|(u_m-u)(\,\cdot+he_i)\|_{L^p(V)}
+\|u_m-u\|_{L^p(V)}
\right)\\
&=
\frac{1}{|h|}
\left(
\|u_m-u\|_{L^p(V+he_i)}
+\|u_m-u\|_{L^p(V)}
\right)\\
&\leqslant
\frac{2}{|h|}\|u_m-u\|_{L^p(W)}.
\end{align*}
因此, 对整个向量差商,
\begin{align*}
\|D^h u_m-D^h u\|_{L^p(V)}
\leqslant
\sum\limits_{i=1}^n
\|D_i^h u_m-D_i^h u\|_{L^p(V)}
\leqslant
\frac{2n}{|h|}\|u_m-u\|_{L^p(W)}
\longrightarrow0,
\quad m\to\infty.
\end{align*}
故 $D^h u_m\to D^h u$ 于 $L^p(V)$。
由于 $u_m$ 光滑，利用$\mathbf{Step}\bm{1}$中的\eqref{mzvjkhshtuwwaa}式可得
\begin{align*}
\|D^h u_m\|_{L^p(V)}\leqslant n\|Du_m\|_{L^p(W)}.
\end{align*}
利用范数的连续性, 令上式$m\to\infty$, 可得
\begin{align*}
\|D^h u\|_{L^p(V)}
\leqslant n\|Du\|_{L^p(W)}\leqslant n\|Du\|_{L^p(U)}.
\end{align*}
因而断言(i)成立.

\noindent$\mathbf{Step}\bm{3}$.{\heiti 利用一致有界性, 对向量差商提取弱收敛子列.}

下面证明断言(ii).假设\eqref{eq::jijiowjiwiw23323rwe}蕴含
\begin{align*}
\sup_{0<h<\frac{\operatorname{dist}(V,\partial U)}{2}}\|D_{i}^{-h}u\|_{L^{p}(V)}<\infty;
\end{align*}
又由$1<p<\infty$和\cref{Functional Analysis-proposition:Lp对偶与自反性}知$L^p(V)$是自反空间，再根据\cref{Functional Analysis-proposition:自反空间中弱列紧性等价与有界},存在函数 $v_{i}\in L^{p}(V)$ 及子列 $h_{k}\to 0$ 使
\begin{align}\label{eqr::JIjiwjijjJFEW8334435JIOFJA-2}
D_{i}^{-h_{k}}u\rightharpoonup v_{i}\quad\text{弱收敛于}\ L^{p}(V).
\end{align}

\noindent$\mathbf{Step}\bm{4}$.{\heiti 利用差商分部积分识别弱导数并得到范数估计.}

固定 $i\in\{1,\cdots,n\}$, 任取 $\phi\in C_c^\infty(V)$, 并将 $\phi$ 在 $V$ 外延拓为零.
由于 $\phi$ 的支集紧包含于 $V$, 对充分小的 $h$, 由\hyperref[theorem:ch05-差商的分部积分公式]{差商的分部积分公式}, 有
\begin{align}
\int_VuD_i^h\phi\,\mathrm{d}x
=
-\int_V(D_i^{-h}u)\phi\,\mathrm{d}x.
\label{eqr::JIjiwjijjJFEW8334435JIOFJA-1}
\end{align}
由于 $\phi$ 光滑且具有紧支集,再结合偏导数定义知$D_i^{h_k}\phi\to\phi_{x_i}$ 一致收敛, 从而再结合\eqref{eqr::JIjiwjijjJFEW8334435JIOFJA-2}\eqref{eqr::JIjiwjijjJFEW8334435JIOFJA-1}式,得到
\begin{align*}
\int_Vu\phi_{x_i}\,\mathrm{d}x
=
\lim\limits_{k\to\infty}
\int_VuD_i^{h_k}\phi\,\mathrm{d}x
=
-\lim\limits_{k\to\infty}
\int_V(D_i^{-h_k}u)\phi\,\mathrm{d}x
=
-\int_Vv_i\phi\,\mathrm{d}x.
\end{align*}
因此 $v_i=u_{x_i}$ 是 $u$ 在 $V$ 中的第 $i$ 个弱偏导数（$i=1,\cdots,n$）.
由于每个 $v_i\in L^p(V)$, 且 $u\in L^p(V)$, 可得
\begin{align*}
u\in W^{1,p}(V),\quad Du=(v_1,\cdots,v_n).
\end{align*}
由\hyperref[Functional Analysis-proposition:范数关于弱收敛的下半连续性]{范数关于弱收敛的下半连续性}得
\begin{align*}
\|Du\|_{L^p(V)}
\leqslant
\varliminf\limits_{k\to\infty}
\|D^{-h_k}u\|_{L^p(V)}
\leqslant C.
\end{align*}
故断言(ii)成立.
\end{proof}

\begin{proposition}[平直边界附近的切向差商估计]\label{proposition:ch05-平直边界附近的切向差商估计}
设 $n\geqslant2$, $0<r<R<\infty$, 并记
\begin{align*}
U\triangleq B^{0}(0,R)\cap\{x_n>0\},\quad
V\triangleq B^{0}(0,r)\cap\{x_n>0\}.
\end{align*}
设 $1\leqslant p<\infty$, $u\in W^{1,p}(U)$. 对每个切向方向 $i=1,\cdots,n-1$ 以及每个满足 $0<|h|<R-r$ 的步长 $h$, 差商
\begin{align*}
D_i^h u(x)\triangleq\frac{u(x+he_i)-u(x)}{h}
\end{align*}
在 $V$ 上几乎处处有定义, 且满足
\begin{align}\label{eq:ch05-TanDQ-r7K2m9P4}
\int_V|D_i^h u(x)|^p\,\mathrm{d}x
\leqslant
\int_U|u_{x_i}(x)|^p\,\mathrm{d}x.
\end{align}
这里 $e_i$ 是第 $i$ 个标准基向量.

特别地, 当 $R=1$, $r=\frac{1}{2}$ 时, 尽管 $V\subset\subset U$ 不成立, 估计 \eqref{eq:ch05-TanDQ-r7K2m9P4} 对所有 $i=1,\cdots,n-1$ 和 $0<|h|<\frac{1}{2}$ 仍然成立.
\end{proposition}

\begin{proof}
固定 $i\in\{1,\cdots,n-1\}$ 以及满足 $0<|h|<R-r$ 的步长 $h$.

\noindent$\mathbf{Step}\bm{1}$.{\heiti 确认切向平移后的线段始终位于 $U$ 内.}

对任意 $x\in V$ 和 $t\in[0,1]$, 由于 $i\leqslant n-1$, 切向平移不改变第 $n$ 个坐标, 因而
\begin{align*}
(x+the_i)_n=x_n>0.
\end{align*}
另一方面,
\begin{align*}
|x+the_i|
\leqslant |x|+t|h|
<r+|h|
<R.
\end{align*}
所以
\begin{align*}
x+the_i\in U,\quad
x\in V,\quad t\in[0,1].
\end{align*}
特别地, 对所有 $t\in[0,1]$, 都有
\begin{align*}
V+the_i\subseteq U.
\end{align*}
因此, 无论 $h$ 为正还是为负, 差商中涉及的取值点都位于 $U$ 内.

\noindent$\mathbf{Step}\bm{2}$.{\heiti 估计光滑函数的差商.}

先设 $u\in C^\infty(U)\cap W^{1,p}(U)$. 对任意 $x\in V$, 由微分学基本定理得
\begin{align*}
u(x+he_i)-u(x)
=\int_0^1\frac{\mathrm{d}}{\mathrm{d}t}u(x+the_i)\,\mathrm{d}t
=h\int_0^1u_{x_i}(x+the_i)\,\mathrm{d}t.
\end{align*}
除以 $h$, 得
\begin{align*}
D_i^h u(x)
=\int_0^1u_{x_i}(x+the_i)\,\mathrm{d}t.
\end{align*}
当 $p>1$ 时, 对长度为 $1$ 的区间 $[0,1]$ 使用\hyperref[Real Analysis-theorem:实变函数--Hölder不等式]{Hölder不等式}; 当 $p=1$ 时, 直接使用积分的绝对值不等式. 两种情形均给出
\begin{align*}
|D_i^h u(x)|^p
\leqslant
\left(\int_0^1|u_{x_i}(x+the_i)|\,\mathrm{d}t\right)^p
\leqslant
\int_0^1|u_{x_i}(x+the_i)|^p\,\mathrm{d}t.
\end{align*}
两边在 $V$ 上积分, 利用\hyperref[Real Analysis-theorem:实变函数-Fubini定理]{Fubini定理}和平移换元, 得
\begin{align}
\int_V|D_i^h u(x)|^p\,\mathrm{d}x
&\leqslant
\int_0^1\int_V|u_{x_i}(x+the_i)|^p\,\mathrm{d}x\,\mathrm{d}t=
\int_0^1\int_{V+the_i}|u_{x_i}(y)|^p\,\mathrm{d}y\,\mathrm{d}t\nonumber\\
&\leqslant
\int_0^1\int_U|u_{x_i}(y)|^p\,\mathrm{d}y\,\mathrm{d}t
=
\int_U|u_{x_i}(y)|^p\,\mathrm{d}y.\label{eq::NKZNJIOIJaaaaa}
\end{align}
这就证明了光滑情形下的结论.

\noindent$\mathbf{Step}\bm{3}$.{\heiti 截去边界薄层,利用局部光滑逼近.}

现在设 $u\in W^{1,p}(U)$. 为使用局部光滑逼近, 对每个 $0<\varepsilon<r$, 定义
\begin{align*}
V_\varepsilon
\triangleq V\cap\{x_n>\varepsilon\}.
\end{align*}
考虑从 $\overline{V_\varepsilon}$ 出发的所有平移线段组成的集合
\begin{align*}
K_\varepsilon
\triangleq
\left\{
x+the_i:
x\in\overline{V_\varepsilon},\, t\in[0,1]
\right\}.
\end{align*}
由\cref{proposition:紧集平移与紧参数平移集}知$K_\varepsilon$ 是紧集. 对任意 $y=x+the_i\in K_\varepsilon$, 有
\begin{align*}
y_n=x_n\geqslant\varepsilon>0,\quad
|y|\leqslant r+|h|<R.
\end{align*}
因此 $K_\varepsilon$ 是 $U$ 的紧子集. 由\cref{proposition:紧集的相对紧开邻域},可以选取开集 $W_\varepsilon$, 使得
\begin{align*}
K_\varepsilon\subseteq W_\varepsilon\subset\subset U.
\end{align*}
于是
\begin{align*}
V_\varepsilon+the_i\subseteq W_\varepsilon,\quad t\in[0,1].
\end{align*}

由\hyperref[theorem:ch05-光滑函数的局部逼近定理]{光滑函数的局部逼近定理}, 存在
$u_m\in C^\infty(W_\varepsilon)\cap W^{1,p}(W_\varepsilon)$, 使得
\begin{align*}
\|u_m-u\|_{W^{1,p}(W_\varepsilon)}
\longrightarrow0,\quad m\to\infty.
\end{align*}
由于从 $V_\varepsilon$ 出发的上述平移线段全部包含于 $W_\varepsilon$, 将\eqref{eq::NKZNJIOIJaaaaa}式的计算应用于 $u_m$, 得
\begin{align}\label{NJKHiwiiwwww}
\int_{V_\varepsilon}|D_i^h u_m(x)|^p\,\mathrm{d}x
\leqslant
\int_0^1\int_{V_\varepsilon+the_i}
|(u_m)_{x_i}(y)|^p\,\mathrm{d}y\,\mathrm{d}t
\leqslant
\int_{W_\varepsilon}|(u_m)_{x_i}(y)|^p\,\mathrm{d}y.
\end{align}

下面说明可以在此估计中令 $m\to\infty$. 因为 $h$ 已经固定, 且
$V_\varepsilon,V_\varepsilon+he_i\subseteq W_\varepsilon$, 所以
\begin{align*}
\|D_i^h u_m-D_i^h u\|_{L^p(V_\varepsilon)}
&\leqslant
\frac{1}{|h|}
\left(
\|u_m-u\|_{L^p(V_\varepsilon+he_i)}
+
\|u_m-u\|_{L^p(V_\varepsilon)}
\right)\\
&\leqslant
\frac{2}{|h|}
\|u_m-u\|_{L^p(W_\varepsilon)}
\longrightarrow0,\quad m\to\infty.
\end{align*}
同时,还有
\begin{align*}
\|(u_m)_{x_i}-u_{x_i}\|_{L^p(W_\varepsilon)}\leqslant \|u_m-u\|_{L^p(W_\varepsilon)}
\longrightarrow0,\quad m\to\infty.
\end{align*}
因此, 由 $L^p$ 范数的连续性, 在\eqref{NJKHiwiiwwww}式中令 $m\to\infty$, 得
\begin{align*}
\int_{V_\varepsilon}|D_i^h u(x)|^p\,\mathrm{d}x
\leqslant
\int_{W_\varepsilon}|u_{x_i}(x)|^p\,\mathrm{d}x\leqslant
\int_U|u_{x_i}(x)|^p\,\mathrm{d}x.
\end{align*}

\noindent$\mathbf{Step}\bm{4}$.{\heiti 令边界薄层的厚度趋于零.}

当 $\varepsilon\to0^+$ 时, 集合 $V_\varepsilon$ 单调增大趋于$V$. 由\hyperref[Real Analysis-theorem:Beppo Levi非负渐升列积分定理]{Levi定理},得
\begin{align*}
\int_V{|D_{i}^{h}u(x)|^p\,\mathrm{d}x}&=\int_{\mathbb{R} ^n}{|D_{i}^{h}u(x)|^p\chi _V\left( x \right) \,\mathrm{d}x}=\lim_{\varepsilon \rightarrow 0^+} \int_{\mathbb{R} ^n}{|D_{i}^{h}u(x)|^p\chi _{V_{\varepsilon}}\left( x \right) \,\mathrm{d}x}
\\
&=\lim_{\varepsilon \rightarrow 0^+} \int_{V_{\varepsilon}}{|D_{i}^{h}u(x)|^p\,\mathrm{d}x}\leqslant \int_U{|u_{x_i}(x)|^p\,\mathrm{d}x}.
\end{align*}
由于方向 $i\in\{1,\cdots,n-1\}$ 和步长 $0<|h|<R-r$ 是任意选取的, \eqref{eq:ch05-TanDQ-r7K2m9P4}式对所有这些方向和步长均成立.
\end{proof}


\begin{theorem}[$W^{1,\infty}_{\mathrm{loc}}$ 的刻画]\label{theorem:ch05-W1infty局部刻画}
设 $U \subseteq \mathbb{R}^n$ 为开集，$f: U \to \mathbb{R}$。则
$f\ \text{在}\ U\ \text{中局部 Lipschitz 连续当且仅当}f \in W^{1,\infty}_{\mathrm{loc}}(U).$
特别地,$f$的Lipschitz 常数不超过 \(\|Df\|_{L^\infty(W)}\)
\end{theorem}
\begin{proof}

{\heiti 必要性:}
设 $f$ 在 $U$ 中局部 Lipschitz 连续 . 固定 $i\in\{1,\cdots,n\}$ . 任取开集 $V \subset\subset W \subset\subset U$ , 取 $0<|h|<\operatorname{dist}(V,\partial W)$ , 定义差商
\begin{align}\label{eq:局部差商定义}
D_{i}^{-h}f(x)\triangleq\frac{f(x)-f(x-he_i)}{h},\qquad x\in V.
\end{align}
由 $f$ 在 $W$ 上的 Lipschitz 条件 , 有
\begin{align*}
\sup_{0<|h|<\operatorname{dist}(V,\partial W)}\|D_{i}^{-h}f\|_{L^\infty(V)}
\leqslant \operatorname{Lip}(f|_W)<\infty.
\end{align*}
因此差商族 $\{D_{i}^{-h}f\}_{h\neq 0}$ 在 $L^\infty(V)$ 中一致有界 , 从而在 $L^2(V)$ 中也一致有界 . 由\cref{Functional Analysis-proposition:Lp对偶与自反性}知 $L^2(V)$ 是自反的 , 再根据\cref{Functional Analysis-proposition:自反空间中弱列紧性等价与有界} , 存在子列 $h_k\to 0$ 及函数 $g_i\in L^2(V)$ , 使得
\begin{align*}
D_{i}^{-h_k}f\rightharpoonup g_i\quad\text{在}\ L^2(V)\ \text{中弱收敛}.
\end{align*}
另一方面，由 $\|D_i^{-h}f\|_{L^\infty(\mathbb{R}^n)}\leqslant \operatorname{Lip}(f|_W)$ 对一切 $h\neq 0$ 成立，得
\begin{align*}
\left|\int_{\mathbb{R}^n}\phi D_i^{-h_k}f\,\mathrm{d}x\right|
\leqslant \|D_i^{-h_k}f\|_{L^\infty(\mathbb{R}^n)}\int_{\mathbb{R}^n}|\phi|\,\mathrm{d}x
\leqslant \operatorname{Lip}(f|_W)\int_{\mathbb{R}^n}|\phi|\,\mathrm{d}x.
\end{align*}
令 $k\to\infty$，即得
\begin{align*}
\left|\int_{\mathbb{R}^n}g_i\,\phi\,\mathrm{d}x\right|
\leqslant \operatorname{Lip}(u)\int_{\mathbb{R}^n}|\phi|\,\mathrm{d}x.
\end{align*}
由于\cref{Real Analysis-theorem:L无穷空间范数对偶刻画}以及$C_c^\infty(\mathbb{R}^n)$ 在 $L^1(\mathbb{R}^n)$ 中稠密，有
\begin{align*}
\|g_i\|_{L^\infty(\mathbb{R}^n)}
=\sup_{\substack{\phi\in C_c^\infty(\mathbb{R}^n)\\ \|\phi\|_{L^1(\mathbb{R}^n)}= 1}}
\left|\int_{\mathbb{R}^n}g_i\,\phi\,\mathrm{d}x\right|.
\end{align*}
因此综上可得
\begin{align}\label{eq:局部弱极限L无穷界}
\|g_i\|_{L^\infty(V)}\leqslant \operatorname{Lip}(f|_W).
\end{align}

任取 $\phi\in C_c^\infty(V)$ . 由\hyperref[theorem:ch05-差商的分部积分公式]{差商的分部积分公式} , 当$h$充分小时,有
\begin{align}\label{eq:局部差商分部积分}
\int_{V}f(x)D_{i}^{h}\phi(x)\,\mathrm{d}x
=-\int_{V}D_{i}^{-h}f(x)\phi(x)\,\mathrm{d}x.
\end{align}
其中 $D_{i}^{h}\phi(x)=\frac{\phi(x+he_i)-\phi(x)}{h}$ . 因为 $\phi$ 光滑且具有紧支集 , 所以当 $h_k\to 0$ 时 , $D_{i}^{h_k}\phi\to \phi_{x_i}$ 在 $L^\infty(V)$ 中成立 . 对 \eqref{eq:局部差商分部积分} 式令 $h=h_k\to 0$ , 左端由\hyperref[Real Analysis-theorem:控制收敛定理]{控制收敛定理}收敛到 $\int_V f\phi_{x_i}\,\mathrm{d}x$ ; 右端由弱收敛定义收敛到 $-\int_V g_i\phi\,\mathrm{d}x$ . 于是
\begin{align*}
\int_{V}f\phi_{x_i}\,\mathrm{d}x=-\int_{V}g_i\phi\,\mathrm{d}x.
\end{align*}
上式对一切 $\phi\in C_c^\infty(V)$ 成立 , 故 $g_i$ 是 $f$ 关于 $x_i$ 的弱导数 , 即 $g_i=f_{x_i}$ 在弱意义下成立 . 由于 $g_i\in L^\infty(V)$ , 且这对任意 $V\subset\subset U$ 均成立 , 所以 $f\in W^{1,\infty}_{\mathrm{loc}}(U)$ .

{\heiti 充分性：}
设 \(f\in W^{1,\infty}_{\mathrm{loc}}(U)\)。任取闭球 \(B\subset\subset U\)，取开集 \(W\) 使得 \(B\subset\subset W\subset\subset U\)。令 \(f^\varepsilon\triangleq\eta_\varepsilon*f\) 为通常的磨光函数，其中 \(\varepsilon<\operatorname{dist}(B,\partial W)\)。
由\hyperref[theorem:appC-光滑算子的性质]{光滑算子的性质}知 \(f^\varepsilon\in C^\infty(W)\)，且对任意 \(i=1,\cdots,n\)，有
\begin{align*}
D_if^\varepsilon(x)=(\eta_\varepsilon*D_if)(x).
\end{align*}
根据\hyperref[theorem:appC-光滑算子的性质]{光滑算子的性质}，有
\begin{align}\label{eq:磨光导数一致界}
\|Df^\varepsilon\|_{L^\infty(B)}\leqslant \|Df\|_{L^\infty(W)}<\infty.
\end{align}
因为 \(f^\varepsilon\) 是 \(C^\infty\) 的，对任意 \(x,y\in B\)，由微积分基本定理有
\begin{align*}
f^\varepsilon(x)-f^\varepsilon(y)=\int_0^1 Df^\varepsilon(y+t(x-y))\,\mathrm{d}t\cdot(x-y).
\end{align*}
两边取绝对值，并利用\eqref{eq:磨光导数一致界}式，得到
\begin{align*}
|f^\varepsilon(x)-f^\varepsilon(y)|
&\leqslant \|Df^\varepsilon\|_{L^\infty(B)}|x-y|
\leqslant \|Df\|_{L^\infty(W)}|x-y|.
\end{align*}
因此对于固定的 \(x,y\in B\)，令 \(\varepsilon\to 0^+\)，由磨光函数的性质 \(f^\varepsilon\to f\) 在 \(L^1_{\mathrm{loc}}(U)\) 中成立，故由\cref{Real Analysis-corollary:非负可测函数积分为0则几乎处处为0}知,存在子列几乎处处收敛，从而
\begin{align*}
|f(x)-f(y)|\leqslant \|Df\|_{L^\infty(W)}|x-y|.
\end{align*}
这说明 \(f|_B\) 是 Lipschitz 连续的，Lipschitz 常数不超过 \(\|Df\|_{L^\infty(W)}\)。由于闭球 \(B\subset\subset U\) 是任意的，故 \(f\) 在 \(U\) 中局部Lipschitz连续。
\end{proof}



\subsection{扩张定理和$W^{1,\infty}$的刻画}

下面我们的目标是把Sobolev空间 $W^{1,p}(U)$ 中的函数延拓为Sobolev空间 $W^{1,p}(\mathbb{R}^{n})$ 中的函数。这可能是微妙的。例如注意到，把 $u\in W^{1,p}(U)$ 在 $\mathbb{R}^{n}-U$ 中取零来延拓一般行不通，因为我们可能因此沿 $\partial U$ 造成如此恶劣的间断，以致延拓后的函数不再具有弱一阶偏导数。我们必须另想办法延拓 $u$，使之“在穿过 $\partial U$ 时保持弱导数”。

\begin{theorem}[Kirszbraun扩张定理, 实值情形]\label{theorem:实值-Kirszbraun扩张定理}
设 $A\subseteq\mathbb{R}^n$ 是任意非空子集, $f:A\to\mathbb{R}$ 是 $L$-Lipschitz 函数, 即对任意 $x,y\in A$,
\begin{align*}
|f(x)-f(y)|\leqslant L|x-y|.
\end{align*}
则存在 $L$-Lipschitz 函数 $F:\mathbb{R}^n\to\mathbb{R}$, 使得 $F|_A=f$, 且对任意 $x,y\in\mathbb{R}^n$,
\begin{align*}
|F(x)-F(y)|\leqslant L|x-y|.
\end{align*}
\end{theorem}
\begin{remark}
一般情形下, Kirszbraun 定理对向量值函数 $f:A\to\mathbb{R}^m$ 也成立, 即存在 $L$-Lipschitz 扩张 $F:\mathbb{R}^n\to\mathbb{R}^m$. 其证明需要更复杂的几何或泛函分析方法.
\end{remark}
\begin{proof}

定义
\begin{align*}
F(x)\triangleq\inf\limits_{a\in A}\left\{f(a)+L|x-a|\right\},\quad x\in\mathbb{R}^n.
\end{align*}
首先证明 $F$ 是 $L$-Lipschitz 的. 任取 $x,y\in\mathbb{R}^n$. 对任意 $a\in A$, 由三角不等式得
\begin{align*}
f(a)+L|x-a|\leqslant f(a)+L|y-a|+L|x-y|.
\end{align*}
对 $a\in A$ 取下确界, 得
\begin{align*}
F(x)\leqslant F(y)+L|x-y|.
\end{align*}
交换 $x$ 与 $y$ 的位置, 同理可得
\begin{align*}
F(y)\leqslant F(x)+L|x-y|.
\end{align*}
因此
\begin{align*}
|F(x)-F(y)|\leqslant L|x-y|.
\end{align*}

其次证明 $F|_A=f$. 任取 $b\in A$. 一方面, 对任意 $a\in A$, 由 $f$ 的 $L$-Lipschitz性得
\begin{align*}
f(a)+L|b-a|\geqslant f(b).
\end{align*}
对 $a\in A$ 取下确界, 得
\begin{align*}
F(b)=\inf\limits_{a\in A}\left\{f(a)+L|b-a|\right\}\geqslant f(b).
\end{align*}
另一方面, 在上式中取 $a=b$, 得
\begin{align*}
F(b)\leqslant f(b)+L|b-b|=f(b).
\end{align*}
故 $F(b)=f(b)$. 由 $b\in A$ 的任意性, $F|_A=f$. 证毕.
\end{proof}

\begin{lemma}[片状函数的弱导数]\label{lemma:ch05-片状Wkp}
设 $B=B(0,r)\subseteq\mathbb{R}^{n}$，$k\geqslant1$，$1\leqslant p\leqslant\infty$。记
\begin{align*}
B^{+}\triangleq B\cap\{x_n>0\},\quad
B^{-}\triangleq B\cap\{x_n<0\},\quad
B^{0}\triangleq B\cap\{x_n=0\},\quad
x'\triangleq(x_1,\cdots,x_{n-1}).
\end{align*}
若 $f\in C^{k-1}(B)$ 且 $f|_{B^{\pm}}\in W^{k,p}(B^{\pm})$，则 $f\in W^{k,p}(B)$，且对每个满足 $|\alpha|\leqslant k$ 的多重指标 $\alpha$，有
\begin{align*}
D^{\alpha}f=
\begin{cases}
D^{\alpha}(f|_{B^{+}}), & \text{a.e. 于 }B^{+},\\
D^{\alpha}(f|_{B^{-}}), & \text{a.e. 于 }B^{-}.
\end{cases}
\end{align*}
\end{lemma}
\begin{remark}
注意$B^0$是零测集,因此$D^\alpha f$可以在$B^0$上取任意值.
\end{remark}
\begin{proof}

以下总设 $n\geqslant2$；当 $n=1$ 时所有含 $x'$ 的式子按 $\mathbb{R}^{0}$ 为单点集理解，论证完全相同。
    
\noindent$\mathbf{Step}\bm{1}$.{\heiti 构造截断函数并验证乘积属于 Sobolev 空间.}

先设 $k=1$。此时 $f\in C^{0}(B)$，$f|_{B^{\pm}}\in W^{1,p}(B^{\pm})$。由于 $W^{1,\infty}\subseteq W^{1,1}$，且弱导数不依赖于 $p$，以下对 $1\leqslant p\leqslant\infty$ 统一论证。

取 $\theta\in C^{\infty}(\mathbb{R})$，使得
\begin{align*}
0\leqslant\theta\leqslant1,\quad
\theta(t)=0\ (t\leqslant1),\quad
\theta(t)=1\ (t\geqslant2),
\end{align*}
并对 $\varepsilon>0$ 令 $\theta_{\varepsilon}(t)\triangleq\theta\left(\frac{t}{\varepsilon}\right)$。则
\begin{align}
\theta_{\varepsilon}(t)=0\ (t\leqslant\varepsilon),\qquad
\theta_{\varepsilon}(t)=1\ (t\geqslant2\varepsilon).
\label{eq:ch05-theta性质}
\end{align}

任取 $\varphi\in C_{c}^{\infty}(B)$，取 $r'\in(0,r)$ 使 $\operatorname{spt}\varphi\subseteq B(0,r')$。令
\begin{align*}
K_{\varepsilon}\triangleq\operatorname{spt}\varphi\cap\{x_n\geqslant\varepsilon\}.
\end{align*}
由于 $\operatorname{spt}\varphi$ 紧且 $K_{\varepsilon}\subseteq\{x_n\geqslant\varepsilon\}$，$K_{\varepsilon}$ 是紧集；又 $K_{\varepsilon}\subseteq B^{+}$，故 $K_{\varepsilon}$ 紧含于开集 $B^{+}$。由\cref{proposition:紧集的相对紧开邻域}，存在开集 $O_{\varepsilon}$ 使
\begin{align*}
K_{\varepsilon}\subseteq O_{\varepsilon}\subseteq\overline{O_{\varepsilon}}\subseteq B^{+},
\end{align*}
且 $\overline{O_{\varepsilon}}$ 紧。再取 $\rho\in C_{c}^{\infty}(B^{+})$ 使
\begin{align}
0\leqslant\rho\leqslant1,\qquad \rho=1\ \text{于}\ \overline{O_{\varepsilon}}.
\label{eq:ch05-rho性质}
\end{align}
令
\begin{align}
\Theta(x)\triangleq\theta_{\varepsilon}(x_n)\rho(x),\qquad x\in B.
\label{eq:ch05-Theta定义}
\end{align}
则 $\Theta\in C_{c}^{\infty}(B^{+})$。

由于 $f|_{B^{+}}\in W^{1,p}(B^{+})$ 且 $\Theta\in C_{c}^{\infty}(B^{+})$，由\hyperref[theorem:ch05-弱导数性质-4]{Leibniz 公式}，$\Theta f\in W^{1,p}(B^{+})$，且对每个 $i=1,\cdots,n$，
\begin{align}
D_i(\Theta f)=\Theta D_if+fD_i\Theta.
\label{eq:ch05-片状乘积导数}
\end{align}

\noindent$\mathbf{Step}\bm{2}$.{\heiti 在 $B^{+}$ 上验证切向弱导数.}

固定 $i\in\{1,\cdots,n-1\}$。由于$\operatorname{spt}\Theta\subset \subset B^+$,类似$\mathbf{Step}\bm{1}$中$\rho$的取法,同理可取 $\rho_1\in C_{c}^{\infty}(B^{+})$ 使
\begin{align}
0\leqslant\rho_1\leqslant1,\qquad \rho_1=1\ \text{于}\ \operatorname{spt}\Theta.
\label{eq:ch05-rho1性质}
\end{align}
则 $\rho_1\varphi\in C_{c}^{\infty}(B^{+})$。由 $\Theta f\in W^{1,p}(B^{+})$ 及弱导数定义，
\begin{align}
\int_{B^{+}}\Theta fD_i(\rho_1\varphi)\,\mathrm{d}x
=-\int_{B^{+}}D_i(\Theta f)\rho_1\varphi\,\mathrm{d}x.
\label{eq:ch05-片状弱导数}
\end{align}
因为 $\rho_1=1$ 于 $\operatorname{spt}\Theta$，所以
\begin{align*}
\Theta fD_i(\rho_1\varphi)=\Theta fD_i\varphi,\qquad
D_i(\Theta f)\rho_1\varphi=D_i(\Theta f)\varphi.
\end{align*}
代入 \eqref{eq:ch05-片状弱导数} 得
\begin{align}
\int_{B^{+}}\Theta fD_i\varphi\,\mathrm{d}x
=-\int_{B^{+}}D_i(\Theta f)\varphi\,\mathrm{d}x.
\label{eq:ch05-片状Theta}
\end{align}
又因 $i<n$，有 $D_i\theta_{\varepsilon}(x_n)=0$，故
\begin{align*}
D_i\Theta=\rho D_i\theta_{\varepsilon}+\theta_{\varepsilon}D_i\rho=\theta_{\varepsilon}D_i\rho,\,\text{a.e.}\,x\in B^+.
\end{align*}
又由$\mathbf{Step}\bm{1}$知
\begin{align*}
\mathrm{spt}(\theta _{\varepsilon}\varphi )\subset \mathrm{spt}\varphi \cap \mathrm{spt}\theta _{\varepsilon}\subset \mathrm{spt}\varphi \cap \{x_n\geqslant \varepsilon \}=K_{\varepsilon},
\end{align*}
而在 $K_{\varepsilon}$ 上 $\rho=1$，故 $D_i\rho(x)=0,\,\text{a.e.}\,x\in \mathrm{spt}(\theta _{\varepsilon}\varphi ).$于是
\begin{align}
fD_i\Theta\,\varphi=f\theta_{\varepsilon}D_i\rho\,\varphi=0,\,\text{a.e.}\,x\in B^+.
\label{eq:ch05-片状切向消失}
\end{align}

下面把 \eqref{eq:ch05-片状Theta} 的左端换成 $\theta_{\varepsilon}fD_i\varphi$。将 $B^{+}$ 分成两块
\begin{align}
\int_{B^{+}}\Theta fD_i\varphi
&=\int_{B^{+}\cap\{x_n\geqslant\varepsilon\}}\Theta fD_i\varphi
+\int_{B^{+}\cap\{x_n<\varepsilon\}}\Theta fD_i\varphi.
\label{eq:ch05-拆分}
\end{align}
在 $B^{+}\cap\{x_n<\varepsilon\}$ 上，由 \eqref{eq:ch05-theta性质}知$\theta_{\varepsilon}=0$，故由 \eqref{eq:ch05-Theta定义}知$\Theta=0$，于是第二个积分为 $0$。在 $B^{+}\cap\{x_n\geqslant\varepsilon\}$ 上，$x\in\operatorname{spt}\varphi$ 时，因为 $\operatorname{spt}\varphi\cap\{x_n\geqslant\varepsilon\}=K_{\varepsilon}\subseteq\{\rho=1\}$,所以 $\rho(x)=1$，故 $\Theta=\theta_{\varepsilon}$；而 $x\notin\operatorname{spt}\varphi$ 时 $\varphi(x)=0$。因此
\begin{align}
\int_{B^{+}\cap\{x_n\geqslant\varepsilon\}}\Theta fD_i\varphi
=\int_{B^{+}\cap\{x_n\geqslant\varepsilon\}}\theta_{\varepsilon}fD_i\varphi.
\label{eq:ch05-第一块}
\end{align}
又因为 $\theta_{\varepsilon}=0$ 于$B^{+}\cap\{x_n<\varepsilon\}$，所以
\begin{align}
\int_{B^{+}\cap\{x_n<\varepsilon\}}\theta_{\varepsilon}fD_i\varphi=0.
\label{eq:ch05-第二块}
\end{align}
由 \eqref{eq:ch05-拆分}\eqref{eq:ch05-第二块}\eqref{eq:ch05-第一块}，
\begin{align}
\int_{B^{+}}\Theta fD_i\varphi\,\mathrm{d}x
=\int_{B^{+}}\theta_{\varepsilon}fD_i\varphi\,\mathrm{d}x.
\label{eq:ch05-替换Theta}
\end{align}
将 \eqref{eq:ch05-替换Theta} 与 \eqref{eq:ch05-片状乘积导数}\eqref{eq:ch05-片状切向消失} 代入 \eqref{eq:ch05-片状Theta}，得到
\begin{align}
\int_{B^{+}}\theta_{\varepsilon}fD_i\varphi\,\mathrm{d}x
=-\int_{B^{+}}\theta_{\varepsilon}D_if\,\varphi\,\mathrm{d}x.
\label{eq:ch05-片状eps}
\end{align}

\noindent$\mathbf{Step}\bm{3}$.{\heiti 令 $\varepsilon\to0^{+}$，得到 $B^{+}$ 上的切向弱导数.}

在 \eqref{eq:ch05-片状eps} 中，对每个固定的 $x\in B^{+}$，由 \eqref{eq:ch05-theta性质}，当 $\varepsilon<\frac{x_n}{2}$ 时 $\theta_{\varepsilon}(x_n)=1$，故 $\theta_{\varepsilon}(x_n)\to1$（$\varepsilon\to0^{+}$）。又
\begin{align*}
|\theta_{\varepsilon}(x_n)f(x)D_i\varphi(x)|\leqslant|f(x)|\,|D_i\varphi(x)|,\qquad
|\theta_{\varepsilon}(x_n)D_if(x)\varphi(x)|\leqslant|D_if(x)|\,|\varphi(x)|.
\end{align*}
由$f|_{B^{\pm}}\in W^{k,p}(B^{\pm})$知 $f,D_if\in L^{1}(B^{+})$，而 $\varphi,D_i\varphi$ 有界且支集紧含于 $B^{+}$，右端两个函数都属于 $L^{1}(B^{+})$。故由\hyperref[Real Analysis-theorem:控制收敛定理]{控制收敛定理}，在 \eqref{eq:ch05-片状eps} 中令 $\varepsilon\to0^{+}$，得
\begin{align}
\int_{B^{+}}fD_i\varphi\,\mathrm{d}x
=-\int_{B^{+}}D_if\,\varphi\,\mathrm{d}x.
\label{eq:ch05-片状切向}
\end{align}

\noindent$\mathbf{Step}\bm{4}$.{\heiti 计算 $B^{+}$ 上的法向弱导数.}

仍取 $\Theta=\theta_{\varepsilon}\rho$ 如 \eqref{eq:ch05-Theta定义}。在 \eqref{eq:ch05-片状Theta} 中把 $i$ 换成 $n$，再代入 \eqref{eq:ch05-片状乘积导数}，得
\begin{align}
\int_{B^{+}}\Theta fD_n\varphi\,\mathrm{d}x
=-\int_{B^{+}}\Theta D_nf\,\varphi\,\mathrm{d}x
-\int_{B^{+}}fD_n\Theta\,\varphi\,\mathrm{d}x.
\label{eq:ch05-片状法向1}
\end{align}
又由\hyperref[theorem:ch05-弱导数性质-4]{Leibniz公式}，得
\begin{align}\label{eqJKNHJNMZZZ}
D_n\Theta=\rho D_n\theta_{\varepsilon}+\theta_{\varepsilon}D_n\rho
=\theta_{\varepsilon}'\rho+\theta_{\varepsilon}D_n\rho.
\end{align}
由$\mathbf{Step}\bm{1}$知
\begin{align}\label{eqrnjNJKHJV}
\mathrm{spt}(\theta _{\varepsilon}\varphi )\subset \mathrm{spt}\varphi \cap \mathrm{spt}\theta _{\varepsilon}\subset \mathrm{spt}\varphi \cap \{x_n\geqslant \varepsilon \}=K_{\varepsilon}\subset B^+,
\end{align}
而在 $K_{\varepsilon}$ 上 $\rho=1$，故
\begin{align}
\Theta fD_n\varphi=\theta_\varepsilon\rho fD_n\varphi=\theta_\varepsilon fD_n\varphi,\,\,f\theta' _{\varepsilon}\rho \varphi =f\theta'_{\varepsilon}\varphi ,\,\,D_n\rho (x)=0,\quad \mathrm{a}.\mathrm{e}.\,x\in B^+.\label{JKNVZN2342hnj}
\end{align}
再结合\eqref{eqJKNHJNMZZZ}式得
\begin{align}
\int_{B^{+}}fD_n\Theta\,\varphi\,\mathrm{d}x
=\int_{B^{+}}f\theta_{\varepsilon}'\varphi\,\mathrm{d}x.
\label{eq:ch05-片状法向2}
\end{align}

下面计算 \eqref{eq:ch05-片状法向2} 的极限。由\eqref{eqrnjNJKHJV}式知$\mathrm{spt}(\theta'_{\varepsilon}\varphi )\subset B^+$，从而再由\hyperref[Real Analysis-theorem:实变函数-Fubini定理]{Fubini 定理}，得
\begin{align}
\int_{B^+}{f\theta _{\varepsilon}^{\prime}\varphi \,\mathrm{d}x}&=\int_{B^0}{\left( \int_0^{\sqrt{r^2-\left| x^{\prime} \right|^2}}{f(x^{\prime},x_n)\varphi (x^{\prime},x_n)\theta _{\varepsilon}^{\prime}(x_n)\,\mathrm{d}x_n} \right) \mathrm{d}x^{\prime}}\nonumber
\\
&=\int_{B^0}{\left( \int_0^{\infty}{f(x^{\prime} ,x_n)\varphi (x^{\prime},x_n)\theta _{\varepsilon}^{\prime}(x_n)\,\mathrm{d}x_n} \right) \mathrm{d}x^{\prime}}.
\label{eq:ch05-法向内层}
\end{align}
对固定的 $x'$，令 $x_n=\varepsilon s$，则 $\mathrm{d}x_n=\varepsilon\,\mathrm{d}s$，且
\begin{align*}
\theta_{\varepsilon}'(x_n)=\frac{\mathrm{d}}{\mathrm{d}x_n}\theta\left(\frac{x_n}{\varepsilon}\right)=\frac{1}{\varepsilon}\theta'(s).
\end{align*}
故内层积分为
\begin{align}
\int_{0}^{\infty}f(x',\varepsilon s)\varphi(x',\varepsilon s)\theta'(s)\,\mathrm{d}s.
\label{eq:ch05-法向内层2}
\end{align}
由于 $f\in C^{0}(B)$ 且 $\varphi\in C_{c}^{\infty}(B(0,r'))$，函数 $(x',s)\mapsto f(x',\varepsilon s)\varphi(x',\varepsilon s)$ 在 $\{|x'|\leqslant r'\}\times[0,\infty)$ 上有界；又 $\theta'\in C_{c}^{\infty}(\mathbb{R})$ 有界且支集紧。故由\hyperref[Real Analysis-theorem:控制收敛定理]{控制收敛定理}，得
\begin{align*}
\int_{0}^{\infty}f(x',\varepsilon s)\varphi(x',\varepsilon s)\theta'(s)\,\mathrm{d}s
\longrightarrow f(x',0)\varphi(x',0)\int_{0}^{\infty}\theta'(s)\,\mathrm{d}s=f(x',0)\varphi(x',0)\left[\theta(+\infty)-\theta(0)\right]=f(x',0)\varphi(x',0),\, \varepsilon \to 0^+.
\end{align*}
于是 \eqref{eq:ch05-法向内层2} 的极限为 $f(x',0)\varphi(x',0)$。而$f(x',0)\varphi(x',0)$关于$x'$在$\{|x'|\leqslant r'\}$上有界,从而对充分小的$\varepsilon$,由上式知
\begin{align*}
\int_{0}^{\infty}f(x',\varepsilon s)\varphi(x',\varepsilon s)\theta'(s)\,\mathrm{d}s
\end{align*}
也关于$x'$在$B^0=\{|x'|\leqslant r'\}$上有界.于是再对\eqref{eq:ch05-法向内层}式中 $x'$用\hyperref[Real Analysis-theorem:控制收敛定理]{控制收敛定理},得
\begin{align}
\int_{B^{+}}f\theta_{\varepsilon}'\varphi\,\mathrm{d}x
\longrightarrow\int_{B^{0}}f(x',0)\varphi(x',0)\,\mathrm{d}x',\, \varepsilon \to 0^+.
\label{eq:ch05-法向极限}
\end{align}
在 \eqref{eq:ch05-片状法向1}中,结合\eqref{JKNVZN2342hnj}式知,左端等于 $\int_{B^{+}}\theta_{\varepsilon}fD_n\varphi\,\mathrm{d}x$,右端第一项等于$\int_{B^+}{\theta _{\varepsilon}\varphi D_nf\,\mathrm{d}x}$.由\hyperref[Real Analysis-theorem:控制收敛定理]{控制收敛定理}可得
\begin{gather*}
\int_{B^+}{\theta _{\varepsilon}fD_n\varphi \,\mathrm{d}x}\rightarrow \int_{B^+}{fD_n\varphi \,\mathrm{d}x},\,\varepsilon \rightarrow 0^+;
\\
\int_{B^+}{\theta _{\varepsilon}\varphi D_nf\,\mathrm{d}x}\rightarrow -\int_{B^+}{\varphi D_nf},\,\varepsilon \rightarrow 0^+.
\end{gather*}
右端第二项由\eqref{eq:ch05-法向极限}式趋于$-\int_{B^{0}}f\varphi\,\mathrm{d}x'$。故
\begin{align}
\int_{B^{+}}fD_n\varphi\,\mathrm{d}x
=-\int_{B^{+}}D_nf\,\varphi\,\mathrm{d}x
-\int_{B^{0}}f\varphi\,\mathrm{d}x'.
\label{eq:ch05-片状法向+}
\end{align}

\noindent$\mathbf{Step}\bm{5}$.{\heiti 同理得到 $B^{-}$ 上的切向弱导数.}

固定 $i\in\{1,\cdots,n-1\}$。在 $B^{-}$ 上取
\begin{align*}
\theta_{\varepsilon}^{-}(x_n)\triangleq\theta\left(-\frac{x_n}{\varepsilon}\right),\qquad x_n<0.
\end{align*}
由 $\theta$ 的取法，
\begin{align}
\theta_{\varepsilon}^{-}(t)=0\ (t\geqslant-\varepsilon),\qquad
\theta_{\varepsilon}^{-}(t)=1\ (t\leqslant-2\varepsilon),
\label{eq:ch05-片状负切向theta性质}
\end{align}
且对每个固定的 $t<0$，当 $0<\varepsilon<-\frac{t}{2}$ 时 $\theta_{\varepsilon}^{-}(t)=1$，故
\begin{align}
\theta_{\varepsilon}^{-}(t)\longrightarrow1\quad(\varepsilon\to0^{+}).
\label{eq:ch05-片状负切向theta极限}
\end{align}

任取 $\varphi\in C_{c}^{\infty}(B)$，取 $r'\in(0,r)$ 使 $\operatorname{spt}\varphi\subseteq B(0,r')$。令
\begin{align*}
K_{\varepsilon}^{\mathrm{t}}\triangleq\operatorname{spt}\varphi\cap\{x_n\leqslant-\varepsilon\}.
\end{align*}
由于 $\operatorname{spt}\varphi$ 紧且 $K_{\varepsilon}^{\mathrm{t}}\subseteq\{x_n\leqslant-\varepsilon\}$，$K_{\varepsilon}^{\mathrm{t}}$ 是紧集；又 $K_{\varepsilon}^{\mathrm{t}}\subseteq B^{-}$，故 $K_{\varepsilon}^{\mathrm{t}}$ 紧含于开集 $B^{-}$。由\cref{proposition:紧集的相对紧开邻域}，存在开集 $O_{\varepsilon}^{\mathrm{t}}$ 使
\begin{align*}
K_{\varepsilon}^{\mathrm{t}}\subseteq O_{\varepsilon}^{\mathrm{t}}\subseteq\overline{O_{\varepsilon}^{\mathrm{t}}}\subseteq B^{-},
\end{align*}
且 $\overline{O_{\varepsilon}^{\mathrm{t}}}$ 紧。取 $\rho_{\mathrm{t}}\in C_{c}^{\infty}(B^{-})$ 使
\begin{align}
0\leqslant\rho_{\mathrm{t}}\leqslant1,\qquad \rho_{\mathrm{t}}=1\ \text{于}\ \overline{O_{\varepsilon}^{\mathrm{t}}}.
\label{eq:ch05-片状负切向rho}
\end{align}
令
\begin{align}
\Theta_{\mathrm{t}}(x)\triangleq\theta_{\varepsilon}^{-}(x_n)\rho_{\mathrm{t}}(x),\qquad x\in B^{-}.
\label{eq:ch05-片状负切向Theta}
\end{align}
则 $\Theta_{\mathrm{t}}\in C_{c}^{\infty}(B^{-})$。

由于 $f|_{B^{-}}\in W^{1,p}(B^{-})$ 且 $\Theta_{\mathrm{t}}\in C_{c}^{\infty}(B^{-})$，由\hyperref[theorem:ch05-弱导数性质-4]{Leibniz 公式}，$\Theta_{\mathrm{t}}f\in W^{1,p}(B^{-})$，且对每个 $i=1,\cdots,n$，
\begin{align}
D_i(\Theta_{\mathrm{t}}f)=\Theta_{\mathrm{t}}D_if+fD_i\Theta_{\mathrm{t}}.
\label{eq:ch05-片状负切向乘积导数}
\end{align}
类似$\rho_t$的取法，取 $\rho_{1,\mathrm{t}}\in C_{c}^{\infty}(B^{-})$ 使
\begin{align}
0\leqslant\rho_{1,\mathrm{t}}\leqslant1,\qquad \rho_{1,\mathrm{t}}=1\ \text{于}\ \operatorname{spt}\Theta_{\mathrm{t}}.
\label{eq:ch05-片状负切向rho1}
\end{align}
则 $\rho_{1,\mathrm{t}}\varphi\in C_{c}^{\infty}(B^{-})$。由 $\Theta_{\mathrm{t}}f\in W^{1,p}(B^{-})$ 及弱导数定义得
\begin{align}
\int_{B^{-}}\Theta_{\mathrm{t}}fD_i(\rho_{1,\mathrm{t}}\varphi)\,\mathrm{d}x
=-\int_{B^{-}}D_i(\Theta_{\mathrm{t}}f)\rho_{1,\mathrm{t}}\varphi\,\mathrm{d}x.
\label{eq:ch05-片状负切向弱导数}
\end{align}
因为 $\rho_{1,\mathrm{t}}=1$ 于 $\operatorname{spt}\Theta_{\mathrm{t}}$，所以
\begin{align*}
\Theta_{\mathrm{t}}fD_i(\rho_{1,\mathrm{t}}\varphi)=\Theta_{\mathrm{t}}fD_i\varphi,\qquad
D_i(\Theta_{\mathrm{t}}f)\rho_{1,\mathrm{t}}\varphi=D_i(\Theta_{\mathrm{t}}f)\varphi.
\end{align*}
代入 \eqref{eq:ch05-片状负切向弱导数} 得
\begin{align}
\int_{B^{-}}\Theta_{\mathrm{t}}fD_i\varphi\,\mathrm{d}x
=-\int_{B^{-}}D_i(\Theta_{\mathrm{t}}f)\varphi\,\mathrm{d}x.
\label{eq:ch05-片状负切向Theta}
\end{align}
又因 $i<n$，所以$D_i\theta_{\varepsilon}^{-}(x_n)=0$，故
\begin{align*}
D_i\Theta_{\mathrm{t}}
=\rho_{\mathrm{t}}D_i\theta_{\varepsilon}^{-}
+\theta_{\varepsilon}^{-}D_i\rho_{\mathrm{t}}
=\theta_{\varepsilon}^{-}D_i\rho_{\mathrm{t}}.
\end{align*}
在 $\operatorname{spt}(\theta_{\varepsilon}^{-}\varphi)$ 上 $\rho_{\mathrm{t}}=1$，故 $D_i\rho_{\mathrm{t}}=0$；于是
\begin{align}
fD_i\Theta_{\mathrm{t}}\,\varphi
=f\theta_{\varepsilon}^{-}D_i\rho_{\mathrm{t}}\,\varphi=0
\quad\text{a.e. }x\in B^{-}.
\label{eq:ch05-片状负切向消失}
\end{align}

下面把 \eqref{eq:ch05-片状负切向Theta} 的左端换成 $\theta_{\varepsilon}^{-}fD_i\varphi$。将 $B^{-}$ 分成两块：
\begin{align}
\int_{B^{-}}\Theta_{\mathrm{t}}fD_i\varphi
&=\int_{B^{-}\cap\{x_n\leqslant-\varepsilon\}}\Theta_{\mathrm{t}}fD_i\varphi
+\int_{B^{-}\cap\{x_n>-\varepsilon\}}\Theta_{\mathrm{t}}fD_i\varphi.
\label{eq:ch05-片状负切向拆分}
\end{align}
在 $B^{-}\cap\{x_n>-\varepsilon\}$ 上，由 \eqref{eq:ch05-片状负切向theta性质}，$\theta_{\varepsilon}^{-}=0$，故由 \eqref{eq:ch05-片状负切向Theta}，$\Theta_{\mathrm{t}}=0$，于是第二个积分为 $0$。在 $B^{-}\cap\{x_n\leqslant-\varepsilon\}$ 上，若 $x\in\operatorname{spt}\varphi$，则
\begin{align*}
x\in\operatorname{spt}\varphi\cap\{x_n\leqslant-\varepsilon\}=K_{\varepsilon}^{\mathrm{t}}\subseteq\{\rho_{\mathrm{t}}=1\},
\end{align*}
故 $\rho_{\mathrm{t}}(x)=1$，从而 $\Theta_{\mathrm{t}}(x)=\theta_{\varepsilon}^{-}(x_n)$；若 $x\notin\operatorname{spt}\varphi$，则 $\varphi(x)=0$。因此
\begin{align}
\int_{B^{-}\cap\{x_n\leqslant-\varepsilon\}}\Theta_{\mathrm{t}}fD_i\varphi
=\int_{B^{-}\cap\{x_n\leqslant-\varepsilon\}}\theta_{\varepsilon}^{-}fD_i\varphi.
\label{eq:ch05-片状负切向第一块}
\end{align}
又因为 $\theta_{\varepsilon}^{-}=0$ 于 $B^{-}\cap\{x_n>-\varepsilon\}$，
\begin{align}
\int_{B^{-}\cap\{x_n>-\varepsilon\}}\theta_{\varepsilon}^{-}fD_i\varphi=0.
\label{eq:ch05-片状负切向第二块}
\end{align}
由 \eqref{eq:ch05-片状负切向拆分}\eqref{eq:ch05-片状负切向第一块}\eqref{eq:ch05-片状负切向第二块}，
\begin{align}
\int_{B^{-}}\Theta_{\mathrm{t}}fD_i\varphi\,\mathrm{d}x
=\int_{B^{-}}\theta_{\varepsilon}^{-}fD_i\varphi\,\mathrm{d}x.
\label{eq:ch05-片状负切向替换}
\end{align}
将 \eqref{eq:ch05-片状负切向乘积导数}\eqref{eq:ch05-片状负切向替换}\eqref{eq:ch05-片状负切向消失} 代入 \eqref{eq:ch05-片状负切向Theta}，得
\begin{align}
\int_{B^{-}}\theta_{\varepsilon}^{-}fD_i\varphi\,\mathrm{d}x
=-\int_{B^{-}}\theta_{\varepsilon}^{-}D_if\,\varphi\,\mathrm{d}x.
\label{eq:ch05-片状负eps}
\end{align}

最后令 $\varepsilon\to0^{+}$。由 \eqref{eq:ch05-片状负切向theta极限}，对每个 $x\in B^{-}$，$\theta_{\varepsilon}^{-}(x_n)\to1$。又由 $B^{-}$ 有界及 $1\leqslant p\leqslant\infty$，有 $W^{1,p}(B^{-})\subseteq W^{1,1}(B^{-})$，故 $f,D_if\in L^{1}(B^{-})$；而 $\varphi,D_i\varphi$ 有界且支集紧含于 $B$，故
\begin{align*}
|fD_i\varphi|\leqslant\|D_i\varphi\|_{L^{\infty}}|f|\in L^{1}(B^{-}),\qquad
|D_if\,\varphi|\leqslant\|\varphi\|_{L^{\infty}}|D_if|\in L^{1}(B^{-}).
\end{align*}
由\hyperref[Real Analysis-theorem:控制收敛定理]{控制收敛定理}，在 \eqref{eq:ch05-片状负eps} 中令 $\varepsilon\to0^{+}$，得
\begin{align}
\int_{B^{-}}fD_i\varphi\,\mathrm{d}x
=-\int_{B^{-}}D_if\,\varphi\,\mathrm{d}x,\qquad i=1,\cdots,n-1.
\label{eq:ch05-片状切向-}
\end{align}

\noindent$\mathbf{Step}\bm{6}$.{\heiti 计算 $B^{-}$ 上的法向弱导数并合并.}

沿用 $\mathbf{Step}\bm{4}$ 的 $\theta_{\varepsilon}^{-}$，其性质见 \eqref{eq:ch05-片状负切向theta性质}\eqref{eq:ch05-片状负切向theta极限}，并注意
\begin{align}
\theta_{\varepsilon}^{-\prime}(t)=-\frac{1}{\varepsilon}\theta'\left(-\frac{t}{\varepsilon}\right),\qquad t<0.
\label{eq:ch05-片状法向theta导数}
\end{align}
与 $\mathbf{Step}\bm{4}$ 相同地取 $K_{\varepsilon}^{\mathrm{n}},O_{\varepsilon}^{\mathrm{n}},\rho_{\mathrm{n}},\Theta_{\mathrm{n}},\rho_{1,\mathrm{n}}$，则 $\Theta_{\mathrm{n}}\in C_{c}^{\infty}(B^{-})$，
\begin{align}
\Theta_{\mathrm{n}}=\theta_{\varepsilon}^{-}\quad\text{于}\ \operatorname{spt}\varphi\cap\{x_n\leqslant-\varepsilon\},
\label{eq:ch05-片状法向Theta定义}
\end{align}
且同 $\mathbf{Step}\bm{4}$ 由弱导数定义与\hyperref[theorem:ch05-弱导数性质-4]{Leibniz 公式}得
\begin{align}
\int_{B^{-}}\Theta_{\mathrm{n}}fD_i\varphi\,\mathrm{d}x
=-\int_{B^{-}}D_i(\Theta_{\mathrm{n}}f)\varphi\,\mathrm{d}x.
\label{eq:ch05-片状法向Theta}
\end{align}
取 $i=n$，代入 $D_n(\Theta_{\mathrm{n}}f)=\Theta_{\mathrm{n}}D_nf+fD_n\Theta_{\mathrm{n}}$，得
\begin{align}
\int_{B^{-}}\Theta_{\mathrm{n}}fD_n\varphi\,\mathrm{d}x
=-\int_{B^{-}}\Theta_{\mathrm{n}}D_nf\,\varphi\,\mathrm{d}x
-\int_{B^{-}}fD_n\Theta_{\mathrm{n}}\,\varphi\,\mathrm{d}x.
\label{eq:ch05-片状法向1}
\end{align}
又
\begin{align}
D_n\Theta_{\mathrm{n}}
=\rho_{\mathrm{n}}D_n\theta_{\varepsilon}^{-}+\theta_{\varepsilon}^{-}D_n\rho_{\mathrm{n}}
=\theta_{\varepsilon}^{-\prime}\rho_{\mathrm{n}}+\theta_{\varepsilon}^{-}D_n\rho_{\mathrm{n}},
\label{eq:ch05-片状法向导数}
\end{align}
而在 $\operatorname{spt}(\theta_{\varepsilon}^{-}\varphi)$ 上 $\rho_{\mathrm{n}}=1$，故 $D_n\rho_{\mathrm{n}}=0$，且 $\Theta_{\mathrm{n}}=\theta_{\varepsilon}^{-}$，于是
\begin{align}
fD_n\Theta_{\mathrm{n}}\,\varphi=f\theta_{\varepsilon}^{-\prime}\varphi,
\qquad
\int_{B^{-}}fD_n\Theta_{\mathrm{n}}\,\varphi\,\mathrm{d}x
=\int_{B^{-}}f\theta_{\varepsilon}^{-\prime}\varphi\,\mathrm{d}x.
\label{eq:ch05-片状法向2}
\end{align}
同 $\mathbf{Step}\bm{4}$ 把 $B^{-}$ 拆成 $\{x_n\leqslant-\varepsilon\}$ 与 $\{x_n>-\varepsilon\}$：在前者上 $\Theta_{\mathrm{n}}=\theta_{\varepsilon}^{-}$，在后者上 $\theta_{\varepsilon}^{-}=0$，故
\begin{align}
\int_{B^{-}}\Theta_{\mathrm{n}}fD_n\varphi\,\mathrm{d}x
=\int_{B^{-}}\theta_{\varepsilon}^{-}fD_n\varphi\,\mathrm{d}x.
\label{eq:ch05-片状法向替换}
\end{align}

下面计算 \eqref{eq:ch05-片状法向2} 的极限。由 Fubini 定理，
\begin{align}
\int_{B^{-}}f\theta_{\varepsilon}^{-\prime}\varphi\,\mathrm{d}x
&=\int_{B^{0}}\left(\int_{-\sqrt{r^{2}-|x'|^{2}}}^{0}f(x',x_n)\varphi(x',x_n)\theta_{\varepsilon}^{-\prime}(x_n)\,\mathrm{d}x_n\right)\mathrm{d}x'
\notag\\
&=\int_{B^{0}}\left(\int_{-\infty}^{0}f(x',x_n)\varphi(x',x_n)\theta_{\varepsilon}^{-\prime}(x_n)\,\mathrm{d}x_n\right)\mathrm{d}x'.
\label{eq:ch05-法向负内层}
\end{align}
令 $x_n=-\varepsilon s$，则
\begin{align*}
\theta_{\varepsilon}^{-\prime}(x_n)=-\frac{1}{\varepsilon}\theta'(s),\qquad \mathrm{d}x_n=-\varepsilon\,\mathrm{d}s,
\end{align*}
且 $x_n:-\infty\to0^{-}$ 对应 $s:+\infty\to0^{+}$，故内层积分为
\begin{align}
\int_{-\infty}^{0}f\varphi\theta_{\varepsilon}^{-\prime}\,\mathrm{d}x_n
&=\int_{+\infty}^{0}f(x',-\varepsilon s)\varphi(x',-\varepsilon s)\left(-\frac{1}{\varepsilon}\theta'(s)\right)(-\varepsilon\,\mathrm{d}s)
\notag\\
&=-\int_{0}^{\infty}f(x',-\varepsilon s)\varphi(x',-\varepsilon s)\theta'(s)\,\mathrm{d}s.
\label{eq:ch05-法向负内层2}
\end{align}
由 $f\in C^{0}(B)$、$\varphi\in C_{c}^{\infty}(B)$ 及 $\theta'\in C_{c}^{\infty}(\mathbb{R})$，对 $s$ 用\hyperref[Real Analysis-theorem:控制收敛定理]{控制收敛定理}得
\begin{align*}
-\int_{0}^{\infty}f(x',-\varepsilon s)\varphi(x',-\varepsilon s)\theta'(s)\,\mathrm{d}s
\longrightarrow-f(x',0)\varphi(x',0)\int_{0}^{\infty}\theta'(s)\,\mathrm{d}s
=-f(x',0)\varphi(x',0),
\end{align*}
再对 $x'$ 用\hyperref[Real Analysis-theorem:控制收敛定理]{控制收敛定理}，得
\begin{align}
\int_{B^{-}}f\theta_{\varepsilon}^{-\prime}\varphi\,\mathrm{d}x
\longrightarrow-\int_{B^{0}}f(x',0)\varphi(x',0)\,\mathrm{d}x',\quad\varepsilon\to0^{+}.
\label{eq:ch05-法向负极限}
\end{align}

在 \eqref{eq:ch05-片状法向1} 中令 $\varepsilon\to0^{+}$。由 \eqref{eq:ch05-片状法向2}\eqref{eq:ch05-片状法向替换}\eqref{eq:ch05-法向负极限} 及\hyperref[Real Analysis-theorem:控制收敛定理]{控制收敛定理}，得
\begin{align}
\int_{B^{-}}fD_n\varphi\,\mathrm{d}x
=-\int_{B^{-}}D_nf\,\varphi\,\mathrm{d}x
-\left(-\int_{B^{0}}f\varphi\,\mathrm{d}x'\right)
=-\int_{B^{-}}D_nf\,\varphi\,\mathrm{d}x
+\int_{B^{0}}f\varphi\,\mathrm{d}x'.
\label{eq:ch05-片状法向-}
\end{align}
与 \eqref{eq:ch05-片状法向+} 相加，沿 $B^{0}$ 的边界项相消，得
\begin{align}
\int_{B}fD_n\varphi\,\mathrm{d}x
=-\int_{B}D_nf\,\varphi\,\mathrm{d}x.
\label{eq:ch05-片状法向}
\end{align}
结合 \eqref{eq:ch05-片状切向}\eqref{eq:ch05-片状切向-}\eqref{eq:ch05-片状法向}，并注意 $\varphi\in C_{c}^{\infty}(B)$ 任意，即知 $f$ 的一阶弱偏导数存在，且
\begin{align*}
Df=
\begin{cases}
D(f|_{B^{+}}), & \text{a.e. 于 }B^{+},\\
D(f|_{B^{-}}), & \text{a.e. 于 }B^{-}.
\end{cases}
\end{align*}
故 $f\in W^{1,p}(B)$，且公式对 $|\alpha|\leqslant1$ 成立。

\noindent$\mathbf{Step}\bm{7}$.{\heiti 归纳证明 $k>1$ 的情形.}

设 $k\geqslant2$，$f\in C^{k-1}(B)$，$f|_{B^{\pm}}\in W^{k,p}(B^{\pm})$。固定 $|\beta|=k-1$，令
\begin{align*}
g\triangleq D^{\beta}f.
\end{align*}
由于 $f\in C^{k-1}(B)$，其直到 $k-1$ 阶的常义导数连续，故
\begin{align*}
g\in C^{0}(B).
\end{align*}
又因为 $f|_{B^{\pm}}\in W^{k,p}(B^{\pm})$ 且 $|\beta|=k-1$，$D^{\beta}(f|_{B^{\pm}})$ 是 $f|_{B^{\pm}}$ 的 $k-1$ 阶弱导数，而它的各一阶弱导数正是 $f|_{B^{\pm}}$ 的 $k$ 阶弱导数，后者属于 $L^{p}(B^{\pm})$，故
\begin{align*}
D^{\beta}(f|_{B^{\pm}})\in W^{1,p}(B^{\pm}).
\end{align*}
另一方面，$f|_{B^{\pm}}\in C^{k-1}(B^{\pm})$,对任意 $\psi\in C_{c}^{\infty}(B^{\pm})$，反复利用
\hyperref[Basis of Analytics-theorem:Green恒等式-3]{Green恒等式\ref{Basis of Analytics-theorem:Green恒等式-3}}得
\begin{align*}
\int_{B^{\pm}}f|_{B^{\pm}}D^{\beta}\psi\,\mathrm{d}x
=(-1)^{|\beta|}\int_{B^{\pm}}D^{\beta}(f|_{B^{\pm}})\psi\,\mathrm{d}x,
\end{align*}
其中边界项因 $\psi$ 在 $B^{\pm}$ 内有紧支集而消失。故$f$的常义导数 $D^{\beta}(f|_{B^{\pm}})$ 就是其弱导数.从而
\begin{align}
g|_{B^{\pm}}=D^{\beta}(f|_{B^{\pm}})\in W^{1,p}(B^{\pm}).
\label{eq:ch05-片状归纳-g}
\end{align}
于是 $g$ 满足 $k=1$ 时的全部假设：$g\in C^{0}(B)$，且由 \eqref{eq:ch05-片状归纳-g}，$g|_{B^{\pm}}\in W^{1,p}(B^{\pm})$。由 $\mathbf{Step}\bm{1}$--$\bm{6}$，即 $k=1$ 的情形，有
\begin{align}
g\in W^{1,p}(B),
\label{eq:ch05-片状归纳-gW}
\end{align}
且对每个 $i\in\{1,\cdots,n\}$，
\begin{align}
D_{i}g=
\begin{cases}
D_{i}D^{\beta}(f|_{B^{+}}), & \text{a.e. 于 }B^{+},\\
D_{i}D^{\beta}(f|_{B^{-}}), & \text{a.e. 于 }B^{-}.
\end{cases}
\label{eq:ch05-片状归纳-Dig}
\end{align}

另一方面，$f\in C^{k-1}(B)$ 蕴含 $g=D^{\beta}f$ 是 $f$ 的弱导数，因为对任意 $\varphi\in C_{c}^{\infty}(B)$，同样反复利用
\hyperref[Basis of Analytics-theorem:Green恒等式-3]{Green恒等式\ref{Basis of Analytics-theorem:Green恒等式-3}}得
\begin{align*}
\int_{B}fD^{\beta}\varphi\,\mathrm{d}x
=(-1)^{|\beta|}\int_{B}g\varphi\,\mathrm{d}x.
\end{align*}
因此对每个 $i\in\{1,\cdots,n\}$，$D_{i}g=D_{i}D^{\beta}f=D^{\beta+e_i}f$ 是 $f$ 的 $k$ 阶弱导数，且由 \eqref{eq:ch05-片状归纳-Dig}式得
\begin{align}
D^{\beta+e_i}f
=D_{i}g
=
\begin{cases}
D_{i}D^{\beta}(f|_{B^{+}})=D^{\beta+e_i}(f|_{B^{+}}), & \text{a.e. 于 }B^{+},\\
D_{i}D^{\beta}(f|_{B^{-}})=D^{\beta+e_i}(f|_{B^{-}}), & \text{a.e. 于 }B^{-}.
\end{cases}
\label{eq:ch05-片状归纳-k阶}
\end{align}
由于任意 $|\alpha|=k$ 都可写成 $\alpha=\beta+e_i$，其中 $|\beta|=k-1$，$i\in\{1,\cdots,n\}$。由 \eqref{eq:ch05-片状归纳-k阶} 知 $f$ 的所有 $k$ 阶弱导数都存在，且属于 $L^{p}(B)$。

对 $|\alpha|\leqslant k-1$，$D^{\alpha}f$ 连续.反复利用\hyperref[Basis of Analytics-theorem:Green恒等式-3]{Green恒等式\ref{Basis of Analytics-theorem:Green恒等式-3}}知，$f$ 的常义导数$D^{\alpha}f$ 就是 $f$ 的弱导数，且属于 $L^{p}(B)$。于是 $f$ 的所有不超过 $k$ 阶的弱导数都存在且属于 $L^{p}(B)$，即
\begin{align*}
f\in W^{k,p}(B).
\end{align*}

最后验证逐块公式。若 $|\alpha|\leqslant k-1$，则 $D^{\alpha}f$ 连续，且 $f|_{B^{\pm}}\in C^{k-1}(B^{\pm})$，故
\begin{align*}
D^{\alpha}f=D^{\alpha}(f|_{B^{\pm}})\quad\text{于}\ B^{\pm}.
\end{align*}
若 $|\alpha|=k$，写 $\alpha=\beta+e_i$，$|\beta|=k-1$，则由 \eqref{eq:ch05-片状归纳-k阶}可知
\begin{align*}
D^{\alpha}f=
\begin{cases}
D^{\alpha}(f|_{B^{+}}), & \text{a.e. 于 }B^{+},\\
D^{\alpha}(f|_{B^{-}}), & \text{a.e. 于 }B^{-}.
\end{cases}
\end{align*}
故结论对一切 $|\alpha|\leqslant k$ 成立。
\end{proof}

\begin{lemma}[$W^{k,\infty}$ 的 $C^{k-1}$ 代表元]\label{lemma:ch05-Wkinfty代表元}
设 $B=B(0,r)\subseteq\mathbb{R}^{n}$，$k\geqslant1$。记
\begin{align*}
B^{+}\triangleq B\cap\{x_n>0\},\quad
B^{-}\triangleq B\cap\{x_n<0\},\quad
B^{0}\triangleq B\cap\{x_n=0\},\quad
x'\triangleq(x_1,\cdots,x_{n-1}).
\end{align*}
设 $u\in W^{k,\infty}(B^{+})$。则 $u$ 有代表元 $\widetilde{u}\in C^{k-1}(\overline{B^{+}})$，使得对每个满足 $|\alpha|\leqslant k-1$ 的多重指标 $\alpha$，$D^{\alpha}\widetilde{u}$ 在 $B^{+}$ 上 Lipschitz 连续，并且存在只依赖于 $k,n$ 的常数 $C$，使得
\begin{align}
\|D^{\alpha}\widetilde{u}\|_{L^{\infty}(B^{+})}
+\operatorname{Lip}(D^{\alpha}\widetilde{u})
\leqslant C\|u\|_{W^{k,\infty}(B^{+})},\quad |\alpha|\leqslant k-1.
\label{eq:ch05-Wkinfty代表元估计}
\end{align}
\end{lemma}
\begin{proof}
\noindent$\mathbf{Step}\bm{1}$.{\heiti 磨光并建立各阶导数的 $L^{\infty}$ 估计.}

取标准磨光子$\eta\in C_{c}^{\infty}(\mathbb{R}^{n})$，$\eta\geqslant0$，$\operatorname{spt}\eta\subseteq\overline{B(0,1)}$，$\int_{\mathbb{R}^{n}}\eta=1$，令 $\eta_{\varepsilon}(x)\triangleq\varepsilon^{-n}\eta\left( \frac{x}{\varepsilon} \right)$，并令
\begin{align*}
u_{\varepsilon}\triangleq\eta_{\varepsilon}*u.
\end{align*}
由\hyperref[theorem:appC-光滑算子的性质]{磨光子的性质}，$u_{\varepsilon}\in C^{\infty}(B_{\varepsilon}^{+})$，其中
\begin{align*}
B_{\varepsilon}^{+}\triangleq\{x\in B^{+}:\operatorname{dist}(x,\partial B^{+})>\varepsilon\},
\end{align*}
由\hyperref[theorem:ch05-光滑函数的局部逼近定理]{光滑函数的局部逼近定理},对每个 $|\alpha|\leqslant k$，有
\begin{align}
D^{\alpha}u_{\varepsilon}\in C^{\infty}(B_{\varepsilon}^{+}),\quad D^{\alpha}u_{\varepsilon}=\eta_{\varepsilon}*D^{\alpha}u
\quad\text{于 }B_{\varepsilon}^{+}.
\label{eq:ch05-磨光导数}
\end{align}
于是对 $x\in B_{\varepsilon}^{+}$，
\begin{align*}
|D^{\alpha}u_{\varepsilon}(x)|
\leqslant\int_{\mathbb{R}^{n}}\eta_{\varepsilon}(x-y)|D^{\alpha}u(y)|\,\mathrm{d}y
\leqslant\|D^{\alpha}u\|_{L^{\infty}(B^{+})}\int_{\mathbb{R}^{n}}\eta_{\varepsilon}(x-y)\,\mathrm{d}y=\|D^{\alpha}u\|_{L^{\infty}(B^{+})}.
\end{align*}
故
\begin{align}
\|D^{\alpha}u_{\varepsilon}\|_{L^{\infty}(B_{\varepsilon}^{+})}
\leqslant\|D^{\alpha}u\|_{L^{\infty}(B^{+})}.
\label{eq:ch05-磨光Linfty}
\end{align}

\noindent$\mathbf{Step}\bm{2}$.{\heiti 证明最高阶磨光导数的 Lipschitz 估计.}

对 $|\alpha|=k-1$，由 \eqref{eq:ch05-磨光导数}，$D^{\alpha}u_{\varepsilon}\in C^{1}(B_{\varepsilon}^{+})$。对 $x,y\in B_{\varepsilon}^{+}$，由微分中值定理，
\begin{align*}
|D^{\alpha}u_{\varepsilon}(x)-D^{\alpha}u_{\varepsilon}(y)|
\leqslant\sum_{i=1}^{n}\|D_iD^{\alpha}u_{\varepsilon}\|_{L^{\infty}(B_{\varepsilon}^{+})}|x_i-y_i|
\leqslant n\|u\|_{W^{k,\infty}(B^{+})}|x-y|,
\end{align*}
其中最后一个不等号用到\eqref{eq:ch05-磨光Linfty}式。于是
\begin{align}
|D^{\alpha}u_{\varepsilon}(x)-D^{\alpha}u_{\varepsilon}(y)|
\leqslant n\|u\|_{W^{k,\infty}(B^{+})}|x-y|,\quad x,y\in B_{\varepsilon}^{+}.
\label{eq:ch05-磨光Lipschitz}
\end{align}

\noindent$\mathbf{Step}\bm{3}$.{\heiti 在紧集上使用 Arzel\`a--Ascoli 定理并取对角线子列.}

取
\begin{align*}
K_m\triangleq\left\{x\in B^{+}:\operatorname{dist}(x,\partial B^{+})\geqslant\frac{1}{m},\ |x|\leqslant r-\frac{1}{m}\right\},\qquad m\geqslant3,
\end{align*}
则 $K_m$ 紧且 $K_m\uparrow B^{+}$，故 $u_{\varepsilon}\in C^{k-1}(K_m)$。由 \eqref{eq:ch05-磨光Linfty}\eqref{eq:ch05-磨光Lipschitz}式，对每个 $|\alpha|\leqslant k-1$，
\begin{align*}
\|D^{\alpha}u_{\varepsilon}\|_{L^{\infty}(K_m)}\leqslant\|u\|_{W^{k,\infty}(B^{+})},\qquad
\operatorname{Lip}(D^{\alpha}u_{\varepsilon}|_{K_m})\leqslant n\|u\|_{W^{k,\infty}(B^{+})}.
\end{align*}
因此族 $\{D^{\alpha}u_{\varepsilon}\}_{0<\varepsilon<\varepsilon_0}$ 在 $C(K_m)$ 中一致有界且等度连续。由\hyperref[Functional Analysis-theorem:泛函分析--Arzelà-Ascoli定理]{Arzel\`a--Ascoli 定理}，$\{D^{\alpha}u_{\varepsilon}\}_{0<\varepsilon<\varepsilon_0}$ 在 $C(K_m)$ 中列紧。
由于满足 $|\alpha|\leqslant k-1$ 的 $\alpha$ 只有有限个，先对每个 $\alpha$ 取子列使其在 $K_m$ 上一致收敛，再其取有限次子列，故可不妨设对一切 $|\alpha|\leqslant k-1$，$\{D^{\alpha}u_{\varepsilon}\}$ 在 $K_m$ 上一致收敛。

于是依次对 $m=1,2,\cdots$, 取子列
\begin{align*}
\{\varepsilon_{1,j}\}_{j=1}^{\infty}\supseteq\{\varepsilon_{2,j}\}_{j=1}^{\infty}\supseteq\cdots,
\end{align*}
使得对每个 $m$ 及每个 $|\alpha|\leqslant k-1$，$\{D^{\alpha}u_{\varepsilon_{m,j}}\}_{j=1}^{\infty}$ 在 $K_m$ 上一致收敛。由\hyperref[Functional Analysis-proposition:对角线法则]{对角线法则}，取对角线子列
\begin{align*}
\varepsilon_l\triangleq\varepsilon_{l,l},\qquad l=1,2,\cdots.
\end{align*}
则对每个固定的 $m$，$\{\varepsilon_l\}_{l\geqslant m}$ 是 $\{\varepsilon_{m,j}\}_{j}$ 的子列，故对每个 $|\alpha|\leqslant k-1$，$\{D^{\alpha}u_{\varepsilon_l}\}_{l}$ 在 $K_m$ 上一致收敛.

由于 $K_m\uparrow B^{+}$，故对$\forall x\in B^+$,存在$m_x>0$,使得$x\in K_{m_x}$.于是对每个 $|\alpha|\leqslant k-1$,可定义
\begin{align}
g_{\alpha}\triangleq \lim_{l\to\infty}D^{\alpha}u_{\varepsilon_l}(x),\quad g_0(x)\triangleq \widetilde{u}(x)\triangleq\lim_{l\to\infty}u_{\varepsilon_l}(x),\qquad x\in B^{+}.
\label{eq:ch05-极限函数}
\end{align}
对每个 $|\alpha|\leqslant k-1$,由$\{D^{\alpha}u_{\varepsilon_l}\}_{l}$在每个 $K_m$ 上一致收敛，且$D^{\alpha}u_{\varepsilon_l}\in C^{0}(K_m)$ 对每个 $m$ 成立，故 $g_{\alpha}\in C^{0}(B^{+})$.

再证 $\widetilde{u}\in C^{k-1}(B^{+})$ 且其常义导数在$K$上一致收敛。对 $|\alpha|\leqslant k-2$ 及 $x,y\in B^{+}$，因为$B^+$是凸集，所以线段 $[x,y]\subset B^{+}$.由微积分基本定理，得
\begin{align*}
D^{\alpha}u_{\varepsilon_l}(x)-D^{\alpha}u_{\varepsilon_l}(y)
=\int_{0}^{1}D^{\alpha+e_i}u_{\varepsilon_l}(y+t(x-y))\cdot(x-y)\,\mathrm{d}t.
\end{align*}
令 $l\to\infty$，得
\begin{align}
g_{\alpha}(x)-g_{\alpha}(y)
=\int_{0}^{1}g_{\alpha+e_i}(y+t(x-y))\cdot(x-y)\,\mathrm{d}t.
\label{eq:ch05-导数极限}
\end{align}
由于$g_{\alpha+e_i}$ 连续和\cref{Real Analysis-corollary:可积函数的Lebesgue点},故结合上式知$g_{\alpha}$ 可微且 $Dg_{\alpha}=g_{\alpha+e_i}$。从 $g_0=\widetilde{u}$ 出发，依次取 $|\alpha|=0,1,\cdots,k-2$，即得
\begin{align*}
D^{\alpha}\widetilde{u}=g_{\alpha},\qquad |\alpha|\leqslant k-1.
\end{align*}
故 $\widetilde{u}\in C^{k-1}(B^{+})$，且对每个紧集 $K\subset B^{+}$ 及每个 $|\alpha|\leqslant k-1$，
\begin{align}
D^{\alpha}u_{\varepsilon_l}\longrightarrow D^{\alpha}\widetilde{u}
\quad\text{在 }K\text{ 上一致}.
\label{eq:ch05-磨光收敛}
\end{align}

\noindent$\mathbf{Step}\bm{4}$.{\heiti 确定极限为 $u$ 的代表元.}

由\hyperref[theorem:appC-光滑算子的性质]{磨光子的性质}，$u_{\varepsilon}\to u$ 于 $L^{1}_{\mathrm{loc}}(B^{+})$。另一方面，由 \eqref{eq:ch05-磨光收敛}知$u_{\varepsilon_l}\to\widetilde{u}$ 在每个紧集上一致，从而也于 $L^{1}_{\mathrm{loc}}(B^{+})$收敛。由 $L^{1}_{\mathrm{loc}}(B^{+})$ 中极限的唯一性知，$\widetilde{u}=u$ a.e. 于 $B^{+}$。即 $\widetilde{u}$ 是 $u$ 的代表元。

\noindent$\mathbf{Step}\bm{5}$.{\heiti 把 Lipschitz 估计与 $L^{\infty}$ 估计过渡到 $\widetilde{u}$.}

在\eqref{eq:ch05-磨光Lipschitz}中令 $\varepsilon_l\to 0$，由 \eqref{eq:ch05-磨光收敛}，对每个紧集 $K\subset B^{+}$ 及 $x,y\in K$，
\begin{align}
|D^{\alpha}\widetilde{u}(x)-D^{\alpha}\widetilde{u}(y)|
\leqslant n\|u\|_{W^{k,\infty}(B^{+})}|x-y|.
\label{eq:ch05-代表元Lipschitz局部}
\end{align}
任取 $x,y\in B^{+}$。由于 $B^{+}$ 是凸集，线段 $[x,y]\subseteq B^{+}$；又 $[x,y]$ 紧且与 $\partial B^{+}$ 不交，故存在紧集 $K\subset B^{+}$ 包含 $[x,y]$。在 \eqref{eq:ch05-代表元Lipschitz局部} 中取此 $K$，即得
\begin{align}
|D^{\alpha}\widetilde{u}(x)-D^{\alpha}\widetilde{u}(y)|
\leqslant n\|u\|_{W^{k,\infty}(B^{+})}|x-y|,\quad x,y\in B^{+}.
\label{eq:ch05-代表元Lipschitz}
\end{align}
同理，在 \eqref{eq:ch05-磨光Linfty} 中令 $\varepsilon_l\to 0$，得
\begin{align}
\|D^{\alpha}\widetilde{u}\|_{L^{\infty}(B^{+})}
\leqslant\|D^{\alpha}u\|_{L^{\infty}(B^{+})}.
\label{eq:ch05-代表元Linfty}
\end{align}
将 \eqref{eq:ch05-代表元Lipschitz} 与 \eqref{eq:ch05-代表元Linfty} 合并，即得 \eqref{eq:ch05-Wkinfty代表元估计}。

\noindent$\mathbf{Step}\bm{6}$.{\heiti 把各阶导数连续延拓到闭包 $\overline{B^{+}}$.}

任取 $z\in B^{0}$ 及点列 $x_k\in B^{+}$，$x_k\to z$。由 \eqref{eq:ch05-代表元Lipschitz}，对每个 $|\alpha|\leqslant k-1$，
\begin{align*}
|D^{\alpha}\widetilde{u}(x_k)-D^{\alpha}\widetilde{u}(x_m)|
\leqslant n\|u\|_{W^{k,\infty}(B^{+})}|x_k-x_m|.
\end{align*}
由于 $\{x_k\}$ 收敛，右端当 $k,m\to\infty$ 时趋于 $0$，故 $\{D^{\alpha}\widetilde{u}(x_k)\}$ 是 Cauchy 列，从而收敛。若另有 $y_k\in B^{+}$，$y_k\to z$，则同理
\begin{align*}
|D^{\alpha}\widetilde{u}(x_k)-D^{\alpha}\widetilde{u}(y_k)|
\leqslant n\|u\|_{W^{k,\infty}(B^{+})}|x_k-y_k|\longrightarrow0,
\end{align*}
故极限不依赖于点列的选取。于是可定义
\begin{align*}
D^{\alpha}\widetilde{u}(z)\triangleq\lim_{x\to z,\ x\in B^{+}}D^{\alpha}\widetilde{u}(x).
\end{align*}
由 \eqref{eq:ch05-代表元Lipschitz} 取极限知延拓后的 $D^{\alpha}\widetilde{u}$ 在 $\overline{B^{+}}$ 上仍满足同一 Lipschitz 估计，且连续。特别地，$\widetilde{u}\in C^{k-1}(\overline{B^{+}})$。
\end{proof}

\begin{lemma}[$C^{k}$ 微分同胚下的复合]\label{lemma:ch05-复合Wkp}
设 $W,B\subseteq\mathbb{R}^{n}$ 为有界开集，$k\geqslant1$，$\Phi:W\to B$ 为 $C^{k}$ 微分同胚，$\Psi=\Phi^{-1}:B\to W$ 也为 $C^{k}$ 微分同胚，且
\begin{align*}
M\triangleq\max\left\{\|D^{\alpha}\Phi\|_{L^{\infty}(W)},\|D^{\alpha}\Psi\|_{L^{\infty}(B)}:1\leqslant|\alpha|\leqslant k\right\}<\infty.
\end{align*}
设 $1\leqslant p\leqslant\infty$，$u\in W^{k,p}(B)$。则 $u\circ\Phi\in W^{k,p}(W)$，且存在只依赖于 $n,k,p,M$ 的常数 $C$，使得
\begin{align}
\|u\circ\Phi\|_{W^{k,p}(W)}\leqslant C\|u\|_{W^{k,p}(B)}.
\label{eq:ch05-复合估计}
\end{align}
更精确地，对每个满足 $1\leqslant|\alpha|\leqslant k$ 的多重指标 $\alpha$，存在非负整数 $c_{\alpha,\beta,\gamma^{1},\cdots,\gamma^{|\beta|}}$（只依赖于 $\alpha,\beta$ 与 $\gamma^{1},\cdots,\gamma^{|\beta|}$），使得对 $u\in C^{k}(\overline{B})$ 及 $x\in W$，
\begin{align}
D^{\alpha}(u\circ\Phi)(x)
=\sum_{\substack{\beta\in\mathbb{N}^{n}\setminus\{0\}\\ 1\leqslant|\beta|\leqslant|\alpha|}}
\sum_{\substack{\gamma^{1},\cdots,\gamma^{|\beta|}\in\mathbb{N}^{n}\setminus\{0\}\\ \gamma^{1}+\cdots+\gamma^{|\beta|}=\alpha}}
c_{\alpha,\beta,\gamma^{1},\cdots,\gamma^{|\beta|}}
(D^{\beta}u)(\Phi(x))\prod_{r=1}^{|\beta|}D^{\gamma^{r}}\Phi(x).
\label{eq:ch05-复合显式}
\end{align}
特别地，
\begin{align}
D_i(u\circ\Phi)(x)&=\sum_{l=1}^{n}(D_lu)(\Phi(x))D_i\Phi_l(x),\label{eq:ch05-复合一阶}\\
D_{ij}(u\circ\Phi)(x)&=\sum_{l,m=1}^{n}(D_{lm}u)(\Phi(x))D_i\Phi_l(x)D_j\Phi_m(x)
+\sum_{l=1}^{n}(D_lu)(\Phi(x))D_{ij}\Phi_l(x).\label{eq:ch05-复合二阶}
\end{align}
且对 $u\in W^{k,p}(B)$，公式 \eqref{eq:ch05-复合显式} 几乎处处成立。
\end{lemma}
\begin{proof}
\noindent$\mathbf{Step}\bm{1}$.{\heiti 证明 $|\alpha|=1$ 的显式公式.}

设 $u\in C^{1}(\overline{B})$。由\hyperref[theorem:向量求导基本性质]{向量求导基本性质\ref{theorem:向量求导基本性质-4}}，对 $x\in W$，
\begin{align*}
D_i(u\circ\Phi)(x)=\sum_{l=1}^{n}(D_lu)(\Phi(x))D_i\Phi_l(x).
\end{align*}
这正是 \eqref{eq:ch05-复合一阶}。在 \eqref{eq:ch05-复合显式} 中取 $|\alpha|=1$，则 $\alpha=e_i$，$\beta=e_l$（$1\leqslant l\leqslant n$），$|\beta|=1$，$\gamma^{1}=e_i$，系数 $c_{e_i,e_l,e_i}=1$，故 \eqref{eq:ch05-复合显式} 与 \eqref{eq:ch05-复合一阶} 一致。

\noindent$\mathbf{Step}\bm{2}$.{\heiti 对 $|\alpha|$ 归纳证明显式公式 \eqref{eq:ch05-复合显式}.}

设公式 \eqref{eq:ch05-复合显式} 对一切 $|\tilde\alpha|\leqslant m$ 成立，其中 $m\geqslant1$。固定 $|\alpha|=m+1$，写 $\alpha=\tilde\alpha+e_i$，其中 $|\tilde\alpha|=m$。由归纳假设，
\begin{align}
D^{\tilde\alpha}(u\circ\Phi)(x)
=\sum_{\substack{\beta\in\mathbb{N}^{n}\setminus\{0\}\\ 1\leqslant|\beta|\leqslant|\tilde\alpha|}}
\sum_{\substack{\gamma^{1},\cdots,\gamma^{|\beta|}\in\mathbb{N}^{n}\setminus\{0\}\\ \gamma^{1}+\cdots+\gamma^{|\beta|}=\tilde\alpha}}
c_{\tilde\alpha,\beta,\gamma^{1},\cdots,\gamma^{|\beta|}}
(D^{\beta}u)(\Phi(x))\prod_{r=1}^{|\beta|}D^{\gamma^{r}}\Phi(x).
\label{eq:ch05-归纳假设}
\end{align}
对上式两端求 $D_i$。由乘积法则，对每一项
\begin{align*}
(D^{\beta}u)(\Phi(x))\prod_{r=1}^{|\beta|}D^{\gamma^{r}}\Phi(x)
\end{align*}
求 $D_i$ 得到两类项：
\begin{align}
&\sum_{l=1}^{n}(D^{\beta+e_l}u)(\Phi(x))D_i\Phi_l(x)\prod_{r=1}^{|\beta|}D^{\gamma^{r}}\Phi(x),\label{eq:ch05-第一类}\\
&\sum_{s=1}^{|\beta|}(D^{\beta}u)(\Phi(x))D^{\gamma^{s}+e_i}\Phi(x)\prod_{r\ne s}D^{\gamma^{r}}\Phi(x).\label{eq:ch05-第二类}
\end{align}
第一类项对应新的 $\beta'=\beta+e_l$（于是 $|\beta'|=|\beta|+1$）与新的元组
\begin{align*}
(\gamma'^{1},\cdots,\gamma'^{|\beta'|})\triangleq(\gamma^{1},\cdots,\gamma^{|\beta|},e_i),
\end{align*}
其和仍为 $\tilde\alpha+e_i=\alpha$。第二类项对应同一个 $\beta$ 与把第 $s$ 个 $\gamma^{s}$ 换成 $\gamma^{s}+e_i$ 的元组，其和仍为 $\alpha$。因此求导后的每一项都具有 \eqref{eq:ch05-复合显式} 右端的形式。将同类项合并，即得存在非负整数 $c_{\alpha,\beta,\gamma^{1},\cdots}$ 使 \eqref{eq:ch05-复合显式} 对 $\alpha$ 成立。具体地，系数满足递推式
\begin{align}
c_{\alpha,\beta,\gamma^{1},\cdots,\gamma^{|\beta|}}
&=\sum_{l=1}^{n}c_{\tilde\alpha,\beta-e_l,\gamma^{1},\cdots,\gamma^{|\beta|-1}}
+\sum_{s=1}^{|\beta|}c_{\tilde\alpha,\beta,\gamma^{1},\cdots,\gamma^{s}-e_i,\cdots,\gamma^{|\beta|}},
\label{eq:ch05-系数递推}
\end{align}
其中约定：若 $\beta-e_l\notin\mathbb{N}^{n}$、或 $\beta-e_l=0$ 而右端第一项中出现零个 $\gamma$、或 $\gamma^{s}-e_i\notin\mathbb{N}^{n}\setminus\{0\}$，则相应的 $c$ 取 $0$。递推 \eqref{eq:ch05-系数递推} 与初始条件 $c_{e_i,e_l,e_i}=1$ 唯一确定了所有系数。归纳完成。

\noindent$\mathbf{Step}\bm{3}$.{\heiti 写出 $|\alpha|=2$ 的显式公式并验证.}

取 $\alpha=(i,j)$，不妨设 $i\neq j$；$i=j$ 时同理。在 \eqref{eq:ch05-复合显式} 中，$\beta$ 与 $\gamma^{1},\cdots,\gamma^{|\beta|}$ 满足 $\gamma^{1}+\cdots+\gamma^{|\beta|}=\alpha$。可能的情形只有：
\begin{enumerate}[(1)]
\item $\beta=e_l+e_m$（$|\beta|=2$），$\gamma^{1}=e_j$，$\gamma^{2}=e_i$（或交换），系数按递推为 $1$；
\item $\beta=e_l$（$|\beta|=1$），$\gamma^{1}=\alpha$，系数为 $1$。
\end{enumerate}
于是
\begin{align*}
D_{ij}(u\circ\Phi)(x)
&=\sum_{l,m=1}^{n}(D_{lm}u)(\Phi(x))D_i\Phi_l(x)D_j\Phi_m(x)
+\sum_{l=1}^{n}(D_lu)(\Phi(x))D_{ij}\Phi_l(x),
\end{align*}
这正是 \eqref{eq:ch05-复合二阶}。当 $i=j$ 时，第一类项中 $\gamma^{1}=\gamma^{2}=e_i$，系数仍为 $1$，公式不变。直接求导验证：由 \eqref{eq:ch05-复合一阶}，
\begin{align*}
D_jD_i(u\circ\Phi)
=D_j\left[\sum_{l=1}^{n}(D_lu)(\Phi)D_i\Phi_l\right]=\sum_{l=1}^{n}\left[\sum_{m=1}^{n}(D_mD_lu)(\Phi)D_j\Phi_m\right]D_i\Phi_l
+\sum_{l=1}^{n}(D_lu)(\Phi)D_{ji}\Phi_l,
\end{align*}
与 \eqref{eq:ch05-复合二阶} 一致（注意 $D_{lm}u=D_{ml}u$，$D_{ij}\Phi_l=D_{ji}\Phi_l$）。

\noindent$\mathbf{Step}\bm{4}$.{\heiti 由显式公式得到逐点估计.}

由 \eqref{eq:ch05-复合显式} 及 $|D^{\gamma}\Phi(x)|\leqslant M$（$1\leqslant|\gamma|\leqslant k$，$x\in W$），对 $x\in W$，
\begin{align}
|D^{\alpha}(u\circ\Phi)(x)|
&\leqslant\sum_{\substack{\beta\in\mathbb{N}^{n}\setminus\{0\}\\ 1\leqslant|\beta|\leqslant|\alpha|}}
\left(\sum_{\substack{\gamma^{1},\cdots,\gamma^{|\beta|}\in\mathbb{N}^{n}\setminus\{0\}\\ \gamma^{1}+\cdots+\gamma^{|\beta|}=\alpha}}
c_{\alpha,\beta,\gamma^{1},\cdots,\gamma^{|\beta|}}M^{|\beta|}\right)
|(D^{\beta}u)(\Phi(x))|
\nonumber\\
&\leqslant C_{\alpha}\sum_{1\leqslant|\beta|\leqslant|\alpha|}|(D^{\beta}u)(\Phi(x))|,
\label{eq:ch05-复合逐点}
\end{align}
其中
\begin{align}
C_{\alpha}\triangleq\max_{1\leqslant|\beta|\leqslant|\alpha|}
\left(\sum_{\substack{\gamma^{1},\cdots,\gamma^{|\beta|}\in\mathbb{N}^{n}\setminus\{0\}\\ \gamma^{1}+\cdots+\gamma^{|\beta|}=\alpha}}
c_{\alpha,\beta,\gamma^{1},\cdots,\gamma^{|\beta|}}M^{|\beta|}\right)
\label{eq:ch05-Calpha}
\end{align}
只依赖于 $n,k,M$（因为 $\alpha$ 只取有限多个值）。

\noindent$\mathbf{Step}\bm{5}$.{\heiti 对光滑 $u$ 建立 $L^{p}$ 估计.}

设 $N$ 为满足 $1\leqslant|\beta|\leqslant k$ 的多重指标 $\beta$ 的个数（只依赖于 $n,k$）。当 $1\leqslant p<\infty$ 时，由 \eqref{eq:ch05-复合逐点} 及\hyperref[Basis of Analytics-proposition:凸性不等式]{凸性不等式}得
\begin{align*}
\left(\sum_{j=1}^{N}a_j\right)^{p}\leqslant N^{p-1}\sum_{j=1}^{N}a_j^{p}\qquad(a_j\geqslant0),
\end{align*}
得
\begin{align}
\int_{W}|D^{\alpha}(u\circ\Phi)|^{p}\,\mathrm{d}x
&\leqslant C_{\alpha}^{p}\int_{W}\left(\sum_{1\leqslant|\beta|\leqslant|\alpha|}|(D^{\beta}u)(\Phi(x))|\right)^{p}\mathrm{d}x
\nonumber\\
&\leqslant C_{\alpha}^{p}N^{p-1}\sum_{1\leqslant|\beta|\leqslant|\alpha|}\int_{W}|(D^{\beta}u)(\Phi(x))|^{p}\,\mathrm{d}x.
\label{eq:ch05-复合Lp1}
\end{align}
作变量替换 $y=\Phi(x)$，即 $x=\Psi(y)$。由于 $\Phi$ 是 $C^{1}$ 微分同胚，变量替换公式给出
\begin{align*}
\int_{W}|(D^{\beta}u)(\Phi(x))|^{p}\,\mathrm{d}x
=\int_{B}|D^{\beta}u(y)|^{p}|\det D\Psi(y)|\,\mathrm{d}y.
\end{align*}
因 $B$ 有界、$\Psi\in C^{1}$ 且 $D\Psi$ 在 $\overline{B}$ 上连续,所以可设
\begin{align*}
C_{M}\triangleq\sup_{y\in B}|\det D\Psi(y)|<\infty,
\end{align*}
则
\begin{align}
\int_{W}|D^{\alpha}(u\circ\Phi)|^{p}\,\mathrm{d}x
\leqslant C_{\alpha}^{p}N^{p-1}C_{M}\sum_{1\leqslant|\beta|\leqslant|\alpha|}\int_{B}|D^{\beta}u(y)|^{p}\,\mathrm{d}y.
\label{eq:ch05-复合Lp}
\end{align}
又
\begin{align}
\int_{W}|u\circ\Phi|^{p}\,\mathrm{d}x
=\int_{B}|u(y)|^{p}|\det D\Psi(y)|\,\mathrm{d}y
\leqslant C_{M}\int_{B}|u|^{p}\,\mathrm{d}y.
\label{eq:ch05-复合Lp0}
\end{align}
将 \eqref{eq:ch05-复合Lp} 对一切 $|\alpha|\leqslant k$ 求和，并与 \eqref{eq:ch05-复合Lp0} 合并，即得光滑情形的 \eqref{eq:ch05-复合估计}。

当 $p=\infty$ 时，由 \eqref{eq:ch05-复合逐点} 直接得到
\begin{align*}
\|D^{\alpha}(u\circ\Phi)\|_{L^{\infty}(W)}
\leqslant C_{\alpha}\sum_{1\leqslant|\beta|\leqslant|\alpha|}\|D^{\beta}u\|_{L^{\infty}(B)},
\end{align*}
再对 $|\alpha|\leqslant k$ 求和即得 \eqref{eq:ch05-复合估计}式$p=\infty$的情形。

\noindent$\mathbf{Step}\bm{6}$.{\heiti $1\leqslant p<\infty$：用 $C^{\infty}(\overline{B})$ 逼近.}

由于 $B$ 是球，$\partial B$ 为 $C^{\infty}$，由\hyperref[theorem:ch05-直到边界光滑的函数整体逼近定理]{直到边界光滑的函数整体逼近定理}，存在 $u_m\in C^{\infty}(\overline{B})$，使得
\begin{align}\label{eqNVJZNJNZNNN}
\|u_m-u\|_{W^{k,p}(B)}\longrightarrow0,\quad m\to\infty.
\end{align}
对 $u_m-u_l$ 应用已证的光滑情形估计 \eqref{eq:ch05-复合估计}，得
\begin{align*}
\|(u_m-u_l)\circ\Phi\|_{W^{k,p}(W)}
\leqslant C\|u_m-u_l\|_{W^{k,p}(B)}.
\end{align*}
故 $\{u_m\circ\Phi\}$ 是 $W^{k,p}(W)$ 中的 Cauchy 列。由于\hyperref[theorem:ch05-WkpBanach]{Sobolev空间都是Banach空间}，故存在 $v\in W^{k,p}(W)$，使得
\begin{align}\label{eqJNLKZSKLASASAS}
u_m\circ\Phi\longrightarrow v\quad\text{于 }W^{k,p}(W).
\end{align}
另一方面，把\eqref{eq:ch05-复合Lp0}式中$u$ 换成 $u_m-u$，得
\begin{align*}
\|u_m\circ\Phi-u\circ\Phi\|_{L^{p}(W)}
\leqslant C_{M}^{\frac{1}{p}}\|u_m-u\|_{L^{p}(B)}\longrightarrow0,
\end{align*}
故 $u_m\circ\Phi\to u\circ\Phi$ 于 $L^{p}(W)$。由 $L^{p}(W)$ 中极限的唯一性知, 几乎处处有$v=u\circ\Phi$。于是 $u\circ\Phi\in W^{k,p}(W)$.再3把\eqref{eq:ch05-复合估计}式$u$换成$u_m$,并令$m\to \infty$得,\eqref{eq:ch05-复合估计}式对 $u$ 成立。

\noindent$\mathbf{Step}\bm{7}$.{\heiti 把显式公式过渡到 $W^{k,p}$ 函数.}

对每个 $|\beta|\leqslant k$，由变量替换和\eqref{eqNVJZNJNZNNN}式可得
\begin{align*}
\int_{W}|(D^{\beta}u_m)(\Phi(x))-(D^{\beta}u)(\Phi(x))|^{p}\,\mathrm{d}x
\leqslant C_{M}\int_{B}|D^{\beta}u_m-D^{\beta}u|^{p}\,\mathrm{d}y\longrightarrow0,
\end{align*}
故 $D^{\beta}u_m\circ\Phi\to D^{\beta}u\circ\Phi$ 于 $L^{p}(W)$。在 \eqref{eq:ch05-复合显式}式中令 $m\to\infty$,由\eqref{eqJNLKZSKLASASAS}式知,\eqref{eq:ch05-复合显式}式左端 $D^{\alpha}(u_m\circ\Phi)\to D^{\alpha}(u\circ\Phi)$ 于 $L^{p}(W)$；\eqref{eq:ch05-复合显式}式右端每一项
\begin{align*}
(D^{\beta}u_m)(\Phi(x))\prod_{r=1}^{|\beta|}D^{\gamma^{r}}\Phi(x)\to (D^{\beta}u)(\Phi(x))\prod_{r=1}^{|\beta|}D^{\gamma^{r}}\Phi(x)\quad\text{于 }L^p(W).
\end{align*}
中，第一因子于 $L^{p}(W)$ 收敛，由$\Phi\in C^{k}$知第二因子有界且逐点收敛，故乘积于 $L^{p}(W)$ 收敛。于是 \eqref{eq:ch05-复合显式}式对一般的$u\in W^{k,p}(W)$ 几乎处处成立。特别地，取 $|\alpha|=1,2$ 即得 \eqref{eq:ch05-复合一阶}\eqref{eq:ch05-复合二阶}式。

\noindent$\mathbf{Step}\bm{8}$.{\heiti $p=\infty$ 的情形.}

设 $u\in W^{k,\infty}(B)$,则 $u\in W^{k,2}(B)$。由 $\mathbf{Step}\bm{6}$--$\mathbf{7}$（取 $p=2$），$u\circ\Phi\in W^{k,2}(W)$，且对每个 $|\alpha|\leqslant k$，$D^{\alpha}(u\circ\Phi)$ 几乎处处等于 \eqref{eq:ch05-复合显式}式的右端。由 \eqref{eq:ch05-复合逐点}式，得
\begin{align*}
|D^{\alpha}(u\circ\Phi)(x)|
\leqslant C_{\alpha}\sum_{1\leqslant|\beta|\leqslant|\alpha|}\|D^{\beta}u\|_{L^{\infty}(B)},
\quad\text{a.e. }x\in W.
\end{align*}
故 $D^{\alpha}(u\circ\Phi)\in L^{\infty}(W)$，从而 $u\circ\Phi\in W^{k,\infty}(W)$，且由上式对 $|\alpha|\leqslant k$ 求和即得 \eqref{eq:ch05-复合估计}式$p=\infty$的情形。
\end{proof}

\begin{lemma}[半空间上的高阶反射]\label{lemma:ch05-高阶反射}
设 $k\geqslant1$，$B=B(0,r)\subseteq\mathbb{R}^{n}$，记
\begin{align*}
B^{+}\triangleq B\cap\{x_n>0\},\quad
B^{-}\triangleq B\cap\{x_n<0\},\quad
B^{0}\triangleq B\cap\{x_n=0\},\quad
x'\triangleq(x_1,\cdots,x_{n-1}).
\end{align*}
设 $1\leqslant p\leqslant\infty$。则存在常数 $c_1,\cdots,c_{k+1}\in\mathbb{R}$，使得
\begin{align}
\sum_{j=1}^{k+1}c_j\left(-\frac{1}{j}\right)^{m}=1,\quad m=0,1,\cdots,k.
\label{eq:ch05-反射矩条件}
\end{align}
并且对每个 $u\in C^{k}(\overline{B^{+}})$，按下式定义的 $\overline{u}:B\to\mathbb{R}$：
\begin{align}
\overline{u}(x)\triangleq
\begin{cases}
u(x), & x\in B^{+},\\
\displaystyle\sum_{j=1}^{k+1}c_ju\left(x',-\frac{x_n}{j}\right), & x\in B^{-},\\
u(x',0), & x\in B^{0},
\end{cases}
\label{align:ch05-高阶反射}
\end{align}
当$1\leqslant p<\infty$时，有
\begin{enumerate}[(i)]
\item $\overline{u}\in C^{k-1}(B)$，且 $\overline{u}|_{B^{\pm}}\in C^{k}(B^{\pm})$；
\item $\overline{u}\in W^{k,p}(B)$，且存在只依赖于 $k,n,p$ 的常数 $C$，使得
\begin{align}
\|\overline{u}\|_{W^{k,p}(B)}\leqslant C\|u\|_{W^{k,p}(B^{+})}.
\label{eq:ch05-反射估计}
\end{align}
\end{enumerate}
当$p=\infty$时，有$u\in W^{k,\infty}(B^{+})$，且存在 $\widetilde{u}\in C^{k-1}(\overline{B^{+}})$ 是 $u$ 的等价形，以 $\widetilde{u}$ 代替 $u$ 按\eqref{align:ch05-高阶反射}式定义的 $\overline{u}$ 仍满足
\begin{enumerate}[(i)]
\item $\overline{u}\in C^{k-1}(B)$，且 $\overline{u}|_{B^{\pm}}\in C^{k}(B^{\pm})$；
\item $\overline{u}\in W^{k,\infty}(B)$，且存在只依赖于 $k,n,p$ 的常数 $C$，使得
\begin{align}
\|\overline{u}\|_{W^{k,\infty}(B)}\leqslant C\|u\|_{W^{k,\infty}(B^{+})}.
\end{align}
\end{enumerate}
\end{lemma}
\begin{proof}
\noindent$\mathbf{Step}\bm{1}$.{\heiti 证明矩条件方程组有唯一解.}

考虑关于 $c_1,\cdots,c_{k+1}$ 的线性方程组 \eqref{eq:ch05-反射矩条件}。其系数矩阵为
\begin{align*}
A=\left(t_j^{m}\right)_{0\leqslant m\leqslant k,\ 1\leqslant j\leqslant k+1},
\end{align*}
其中 $t_j=-\frac{1}{j}$。由于 $j\mapsto-\frac{1}{j}$ 在 $\{1,\cdots,k+1\}$ 上严格单调，$t_1,\cdots,t_{k+1}$ 两两不同。$A$ 是 Vandermonde 矩阵，其行列式为
\begin{align*}
\det A=\prod_{1\leqslant i<j\leqslant k+1}(t_j-t_i)\neq0.
\end{align*}
故非齐次线性方程组对任意右端常数项存在唯一解。

\noindent$\mathbf{Step}\bm{2}$.{\heiti 证明 $\overline{u}$ 在 $B$ 上连续.}

在 $B^{+}$ 上 $\overline{u}=u$ 连续；在 $B^{-}$ 上 $\overline{u}$ 是 $u$ 与 $C^{k}$ 映射的复合，连续。只需验证 $B^{0}$ 处。任取 $x_0=(x_0',0)\in B^{0}$。当 $x=(x',x_n)\in B^{-}$，$x\to x_0$ 时，对每个 $j$，$\left( x^{\prime},-\frac{x_n}{j} \right) \to(x_0',0)$，由 $u\in C^{k}(\overline{B^{+}})\subseteq C^{0}(\overline{B^{+}})$，得
\begin{align*}
u\left(x',-\frac{x_n}{j}\right)\longrightarrow u(x_0',0).
\end{align*}
于是
\begin{align*}
\overline{u}(x)=\sum_{j=1}^{k+1}c_j u\left(x',-\frac{x_n}{j}\right)
\longrightarrow\left(\sum_{j=1}^{k+1}c_j\right)u(x_0',0)=u(x_0',0),
\end{align*}
其中最后一步用到 \eqref{eq:ch05-反射矩条件} 中 $m=0$ 的等式。而 $x\in B^{+}$ 时 $\overline{u}(x)=u(x)\to u(x_0)$。故 $\overline{u}$ 在 $B$ 上连续。

\noindent$\mathbf{Step}\bm{3}$.{\heiti 归纳证明各阶偏导数存在且连续.}

下面证明：对每个 $|\alpha|\leqslant k-1$，$D^{\alpha}\overline{u}$ 在 $B$ 上存在、连续，且在 $B^{\pm}$ 上由相应的逐块公式给出。

假设对一切 $|\beta|\leqslant m-1$（$1\leqslant m\leqslant k-1$）已证。固定 $|\beta|=m-1$。由归纳假设，$D^{\beta}\overline{u}\in C^{0}(B)$，且在 $B^{\pm}$ 内 $C^{1}$。对 $x\in B^{-}$，由 \eqref{align:ch05-高阶反射}式和求导链式法则得
\begin{align}
D^{\beta}\overline{u}(x)
=\sum_{j=1}^{k+1}c_j\left(-\frac{1}{j}\right)^{\beta_n}
D^{\beta}u\left(x',-\frac{x_n}{j}\right),
\label{eq:ch05-反射导数}
\end{align}
故
\begin{align*}
D_iD^{\beta}\overline{u}(x)
&=\sum_{j=1}^{k+1}c_j\left(-\frac{1}{j}\right)^{\beta_n}
D_iD^{\beta}u\left(x',-\frac{x_n}{j}\right),\quad i<n,\\
D_nD^{\beta}\overline{u}(x)
&=\sum_{j=1}^{k+1}c_j\left(-\frac{1}{j}\right)^{\beta_n+1}
D_nD^{\beta}u\left(x',-\frac{x_n}{j}\right).
\end{align*}
由于 $|\beta|+1=m\leqslant k-1$，有 $\beta_n\leqslant k-1$ 且 $\beta_n+1\leqslant k$。由 \eqref{eq:ch05-反射矩条件}，当 $x\to x_0\in B^{0}$ 且 $x\in B^{-}$ 时，
\begin{align*}
D_iD^{\beta}\overline{u}(x)
&\longrightarrow\sum_{j=1}^{k+1}c_j\left(-\frac{1}{j}\right)^{\beta_n}D_iD^{\beta}u(x_0',0)
=D_iD^{\beta}u(x_0',0),\\
D_nD^{\beta}\overline{u}(x)
&\longrightarrow\sum_{j=1}^{k+1}c_j\left(-\frac{1}{j}\right)^{\beta_n+1}D_nD^{\beta}u(x_0',0)
=D_nD^{\beta}u(x_0',0).
\end{align*}
而从 $B^{+}$ 一侧，$D_iD^{\beta}\overline{u}=D_iD^{\beta}u\to D_iD^{\beta}u(x_0',0)$。故 $D_iD^{\beta}\overline{u}$ 在 $B^{0}$ 两侧有相同的连续极限。

下面证明 $D_iD^{\beta}\overline{u}(x_0)$ 存在且等于该公共极限。

对 $i<n$，取 $h\neq0$ 充分小，使 $x_0+he_i\in B^{0}$。由 $\overline{u}=u$ 于 $B^{0}$，故
\begin{align*}
\frac{D^{\beta}\overline{u}(x_0+he_i)-D^{\beta}\overline{u}(x_0)}{h}
=\frac{D^{\beta}u(x_0+he_i)-D^{\beta}u(x_0)}{h}.
\end{align*}
由于 $u\in C^{k}(\overline{B^{+}})$ 且 $k\geqslant1$，$D^{\beta}u$ 在 $x_0$ 处沿 $e_i$ 方向可微，故上式当 $h\to0$ 时趋于 $D_iD^{\beta}u(x_0)$。

对 $i=n$，先取 $h>0$。函数 $t\mapsto D^{\beta}\overline{u}(x_0+te_n)$ 在 $[0,h]$ 上连续，在 $(0,h)$ 内可微。由微分中值定理，存在 $\theta_h\in(0,h)$，使得
\begin{align*}
\frac{D^{\beta}\overline{u}(x_0+he_n)-D^{\beta}\overline{u}(x_0)}{h}
=D_nD^{\beta}\overline{u}(x_0+\theta_he_n).
\end{align*}
当 $h\to0^{+}$ 时，$x_0+\theta_he_n\in B^{+}$，由 $B^{+}$ 一侧的连续性，
\begin{align*}
D_nD^{\beta}\overline{u}(x_0+\theta_he_n)\longrightarrow D_nD^{\beta}u(x_0).
\end{align*}
再取 $h<0$。同理，在 $B^{-}$ 上对 $t\mapsto D^{\beta}\overline{u}(x_0+te_n)$ 用中值定理，存在 $\theta_h\in(h,0)$，使得
\begin{align*}
\frac{D^{\beta}\overline{u}(x_0+he_n)-D^{\beta}\overline{u}(x_0)}{h}
=D_nD^{\beta}\overline{u}(x_0+\theta_he_n)\longrightarrow D_nD^{\beta}u(x_0).
\end{align*}
故两个单侧差商极限相同，$D_nD^{\beta}\overline{u}(x_0)$ 存在且等于 $D_nD^{\beta}u(x_0)$。

于是 $D^{\beta}\overline{u}$ 的一阶偏导数在 $B$ 上存在且连续，即 $D^{\beta}\overline{u}\in C^{1}(B)$。由归纳法，$\overline{u}\in C^{k-1}(B)$，且 $\overline{u}|_{B^{\pm}}\in C^{k}(B^{\pm})$。

\noindent$\mathbf{Step}\bm{4}$.{\heiti 证明 $\overline{u}\in W^{k,p}(B)$ 并建立估计.}

由 $\mathbf{Step}\bm{3}$，$\overline{u}\in C^{k-1}(B)$，且 $\overline{u}|_{B^{\pm}}\in C^{k}(B^{\pm})\subseteq W^{k,p}(B^{\pm})$。由\cref{lemma:ch05-片状Wkp}，$\overline{u}\in W^{k,p}(B)$，且对每个满足 $|\alpha|\leqslant k$ 的多重指标 $\alpha$，有
\begin{align*}
D^{\alpha}\overline{u}=
\begin{cases}
D^{\alpha}(\overline{u}|_{B^{+}}), & \text{a.e. 于 }B^{+},\\
D^{\alpha}(\overline{u}|_{B^{-}}), & \text{a.e. 于 }B^{-}.
\end{cases}
\end{align*}

当 $1\leqslant p<\infty$ 时，对 $|\alpha|\leqslant k$，由 \eqref{align:ch05-高阶反射}式和求导链式法则得
\begin{align}\label{HJKLVNLKAAA}
D^{\alpha}\overline{u}(x)=\sum_{j=1}^{k+1}c_j\left(-\frac{1}{j}\right)^{\alpha_n}
D^{\alpha}u\left(x',-\frac{x_n}{j}\right),\quad x\in B^{-},
\end{align}
再由\hyperref[Real Analysis-theorem:实变函数--Minkowski不等式]{Minkowski 不等式}得
\begin{align*}
\|D^{\alpha}\overline{u}\|_{L^{p}(B^{-})}
&\leqslant\sum_{j=1}^{k+1}|c_j|\left(\int_{B^{-}}\left|D^{\alpha}u\left(x',-\frac{x_n}{j}\right)\right|^{p}\mathrm{d}x\right)^{\frac{1}{p}}.
\end{align*}
对第 $j$ 项作变量替换
\begin{align*}
y'=x',\qquad y_n=-\frac{x_n}{j},\qquad\text{即 }x_n=-jy_n,\quad \mathrm{d}x_n=j\,\mathrm{d}y_n.
\end{align*}
$B^{-}$ 的像为
\begin{align*}
E_j\triangleq\left\{y\in\mathbb{R}^{n}:y_1^2+\cdots+y_{n-1}^2+j^2y_n^2<r^2,\ y_n>0\right\}\subseteq B^{+},
\end{align*}
且 Jacobi 行列式绝对值为 $j$。故
\begin{align*}
\left(\int_{B^{-}}\left|D^{\alpha}u\left(x',-\frac{x_n}{j}\right)\right|^{p}\mathrm{d}x\right)^{\frac{1}{p}}
=j^{\frac{1}{p}}\left(\int_{E_j}|D^{\alpha}u(y)|^{p}\,\mathrm{d}y\right)^{\frac{1}{p}}\leqslant j^{\frac{1}{p}}\|D^{\alpha}u\|_{L^{p}(B^{+})}.
\end{align*}
于是
\begin{align*}
\|D^{\alpha}\overline{u}\|_{L^{p}(B^{-})}
\leqslant\left(\sum_{j=1}^{k+1}|c_j|j^{\frac{1}{p}}\right)\|D^{\alpha}u\|_{L^{p}(B^{+})}.
\end{align*}
又 $\|D^{\alpha}\overline{u}\|_{L^{p}(B^{+})}=\|D^{\alpha}u\|_{L^{p}(B^{+})}$。将对 $|\alpha|\leqslant k$ 的估计相加，即得 \eqref{eq:ch05-反射估计}，其中
\begin{align*}
C\triangleq1+\sum_{j=1}^{k+1}|c_j|j^{\frac{1}{p}}
\end{align*}
只依赖于 $k,n,p$。

\noindent$\mathbf{Step}\bm{5}$.{\heiti 处理 $W^{k,\infty}$ 与 $C^{k-1}$ 代表元的情形.}

设 $u\in W^{k,\infty}(B^{+})$，则由\cref{lemma:ch05-Wkinfty代表元}可取代表元$\widetilde{u}\in C^{k-1}(\overline{B^{+}})$。以 $\widetilde{u}$ 代替 $u$ 定义 $\overline{u}$。由于 $\mathbf{Step}\bm{3}$ 的归纳论证只用到 $\widetilde{u}$ 直到 $k-1$ 阶的连续性，故同样得到 $\overline{u}\in C^{k-1}(B)$。

又 $\overline{u}|_{B^{+}}=\widetilde{u}\in W^{k,\infty}(B^{+})$。而 $\overline{u}|_{B^{-}}$ 是 $\widetilde{u}$ 与仿射映射 $x\mapsto\left(x',-\frac{x_n}{j}\right)$ 的复合的线性组合。由于仿射映射是 $C^{\infty}$ 微分同胚，由\cref{lemma:ch05-复合Wkp}知 $\overline{u}|_{B^{-}}\in W^{k,\infty}(B^{-})$。于是由\cref{lemma:ch05-片状Wkp}，$\overline{u}\in W^{k,\infty}(B)$，且其弱 $k$ 阶导数在 $B^{-}$ 上几乎处处由\eqref{HJKLVNLKAAA}式给出。对该公式取本性上界，即得
\begin{align*}
\|D^{\alpha}\overline{u}\|_{L^{\infty}(B^{-})}
\leqslant\left(\sum_{j=1}^{k+1}|c_j|\right)\|D^{\alpha}u\|_{L^{\infty}(B^{+})},
\end{align*}
从而 \eqref{eq:ch05-反射估计}式 对 $p=\infty$ 成立。
\end{proof}

\begin{theorem}[扩张定理]\label{theorem:ch05-扩张定理}
设 $k\geqslant1$，$1\leqslant p\leqslant\infty$，$U\subseteq\mathbb{R}^{n}$ 是有界开集，$\partial U$ 为 $C^{k}$，并取有界开集 $V$ 使 $U\subset\subset V$。则存在有界线性算子
\begin{align*}
E_k:W^{k,p}(U)\to W^{k,p}(\mathbb{R}^{n}),
\end{align*}
使得对每个 $u\in W^{k,p}(U)$满足
\begin{enumerate}[(i)]
\item $E_ku=u$ a.e. 于 $U$；
\item $\operatorname{spt}E_ku\subset\subset V$；
\item $\|E_ku\|_{W^{k,p}(\mathbb{R}^{n})}\leqslant C\|u\|_{W^{k,p}(U)}$，其中常数 $C$ 只依赖于 $n,k,p,U,V$。
\end{enumerate}
\end{theorem}
\begin{remark}
\textbf{曲面 $\{z_n=\gamma(z')\}$ 的单位法向量:}\label{曲面的单位法向量选取}
设 $\gamma:\mathbb{R}^{n-1}\to\mathbb{R}$ 为 $C^1$ 函数，考虑曲面
\[
\Sigma\triangleq\{z=(z',z_n)\in\mathbb{R}^n:z_n=\gamma(z')\},
\qquad z'=(z_1,\dots,z_{n-1}).
\]
下面求 $\Sigma$ 在点 $\Phi(z')=(z',\gamma(z'))$ 处的单位法向量。

\noindent\textbf{(1) 参数化与切向量。}
$\Sigma$ 是函数 $\gamma$ 的图像，有自然参数化
\[
\Phi(z')=(z',\gamma(z')),\qquad z'\in\mathbb{R}^{n-1}.
\]
对第 $i$ 个坐标求偏导，得到切向量
\[
\partial_i\Phi
=(0,\dots,1,\dots,0,\partial_i\gamma(z')),
\qquad i=1,\dots,n-1,
\]
其中 $1$ 位于第 $i$ 个位置。这 $n-1$ 个向量张成 $\Sigma$ 在 $\Phi(z')$ 处的切空间。

\noindent\textbf{(2) 法向量与切向量正交。}
设 $N=(N_1,\dots,N_n)\in\mathbb{R}^n$。$N$ 是 $\Sigma$ 的法向量，当且仅当
\[
N\cdot\partial_i\Phi=0,\qquad i=1,\dots,n-1.
\]
即
\[
N_i+N_n\partial_i\gamma(z')=0,\qquad i=1,\dots,n-1.
\]
取 $N_n=1$，便得 $N_i=-\partial_i\gamma(z')$，于是
\[
N=(-\nabla\gamma(z'),1).
\]
也可把 $\Sigma$ 看成水平集 $\{F=0\}$，其中
\[
F(z)=z_n-\gamma(z');
\]
此时
\[
\nabla F(z)=(-\nabla\gamma(z'),1)
\]
就是 $\Sigma$ 的法向量。

\noindent\textbf{(3) 单位化。}
上述法向量的长度为
\[
|N|
=\sqrt{|-\nabla\gamma(z')|^2+1^2}
=\sqrt{1+|\nabla\gamma(z')|^2}.
\]
故 $\Sigma$ 在 $\Phi(z')$ 处的单位法向量为
\[
\nu(z')
=\frac{(-\nabla\gamma(z'),1)}{\sqrt{1+|\nabla\gamma(z')|^2}}.
\]
特别地，在 $z'=0$ 处，
\[
\nu(0)
=\frac{(-\nabla\gamma(0),1)}{\sqrt{1+|\nabla\gamma(0)|^2}}.
\]

\noindent\textbf{(4) 方向的选择。}
法向量有正负两个方向。$N=(-\nabla\gamma(z'),1)$ 的最后一个分量为 $1>0$，所以它指向 $z_n$ 增大的方向。若区域为
\[
U=\{z_n>\gamma(z')\},
\]
则$N$的方向从边界指向 $U$ 的内部，因此 $\nu(z')$ 是 $U$ 在边界处的\textbf{内法向量}；相应的外法向量为
\[
\nu_{\mathrm{out}}(z')
=\frac{(\nabla\gamma(z'),-1)}{\sqrt{1+|\nabla\gamma(z')|^2}}.
\]
在边界拉直坐标中，$Q$ 的选取使 $U$ 的内法向量指向 $z_n$ 轴正方向，即内法向量为 $e_n$。于是
\[
\frac{(-\nabla\gamma(0),1)}{\sqrt{1+|\nabla\gamma(0)|^2}}=e_n.
\]
比较前 $n-1$ 个分量，即得
\[
\nabla\gamma(0)=0.
\]
几何上，这说明曲面在 $0$ 处的切平面是水平的。

\noindent\textbf{(5) 二维情形。}
当 $n=2$ 时，$\Sigma$ 是曲线 $y=\gamma(x)$。切向量为 $(1,\gamma'(x))$，法向量 $(-\gamma'(x),1)$ 与它正交：
\[
(1,\gamma'(x))\cdot(-\gamma'(x),1)
=-\gamma'(x)+\gamma'(x)=0.
\]
单位法向量为
\[
\frac{(-\gamma'(x),1)}{\sqrt{1+\gamma'(x)^2}}.
\]
当 $\gamma\equiv0$ 时它是 $(0,1)=e_2$；当 $\gamma'(0)=0$ 时，曲线在 $0$ 处水平，法向量也正好是 $e_2$。
\end{remark}
\begin{proof}
\noindent$\mathbf{Step}\bm{1}$.{\heiti 在标准半空间上用第 $k$ 阶反射建立 $W^{k,p}$ 估计.}

以下总设 $n\geqslant2$；当 $n=1$ 时所有含 $x'$ 的式子按 $\mathbb{R}^{0}$ 为单点集理解，论证完全相同。设 $B_{r'}\triangleq B(0,r')\subseteq\mathbb{R}^{n}$，记
\begin{align*}
B_{r'}^{+}\triangleq B_{r'}\cap\{x_n>0\},\quad
B_{r'}^{-}\triangleq B_{r'}\cap\{x_n<0\},\quad
B_{r'}^{0}\triangleq B_{r'}\cap\{x_n=0\},\quad
x'\triangleq(x_1,\cdots,x_{n-1}).
\end{align*}

\begin{figure}[H]
\centering
\includegraphics[width=0.55\textwidth]{ch05-half-ball.png}
\caption{边界处的半球}
\label{figure:ch05-边界半球}
\end{figure}

由\cref{lemma:ch05-高阶反射}，存在 $c_1,\cdots,c_{k+1}\in\mathbb{R}$ 满足
\begin{align}
\sum_{j=1}^{k+1}c_j\left(-\frac{1}{j}\right)^{m}=1,\quad m=0,1,\cdots,k,
\label{eq:ch05-统一-矩条件}
\end{align}
且对每个 $u\in C^{k}(\overline{B_{r'}^{+}})$，按
\begin{align}
\overline{u}(x)\triangleq
\begin{cases}
u(x), & x\in B_{r'}^{+},\\
\displaystyle\sum_{j=1}^{k+1}c_ju\left(x',-\frac{x_n}{j}\right), & x\in B_{r'}^{-},\\
u(x',0), & x\in B_{r'}^{0}
\label{eq:ch05-统一-反射}
\end{cases}
\end{align}
定义的 $\overline{u}$ 满足 $\overline{u}\in W^{k,p}(B_{r'})$ 及
\begin{align}
\|\overline{u}\|_{W^{k,p}(B_{r'})}\leqslant C_0\|u\|_{W^{k,p}(B_{r'}^{+})},
\label{eq:ch05-统一-反射估计}
\end{align}
其中 $C_0$ 只依赖于 $k,n,p$。

\noindent$\mathbf{Step}\bm{2}$.{\heiti 对每个边界点构造 $C^{k}$ 拉直坐标邻域.}

固定 $x_0\in\partial U$。设 $\nu(x_0)$ 为 $x_0$ 处关于 $U$ 的单位外法向量。取 $n$ 阶正交矩阵 $Q$，使得 $Q^{T}$ 的最后一列为 $-\nu(x_0)$。定义
\begin{align}
z\triangleq Q(x-x_0).
\label{eq:ch05-统一-坐标}
\end{align}
又因\hyperref[lemma:正交平移变换为等距双射]{正交平移变换为等距双射},故在此坐标系下 $x_0$ 对应 $z=0$，且 $U$ 在 $x_0$ 处的内法向量指向 $z_n$ 轴正方向。由\cref{definition:appC-Ck边界}，存在 $r>0$ 与 $\gamma\in C^{k}(\mathbb{R}^{n-1})$，使得
\begin{align}
U\cap B_z(0,r)=\left\{z\in B_z(0,r):z_n>\gamma(z')\right\},
\label{eq:ch05-统一-边界}
\end{align}
其中 $B_z(0,r)$ 是 $z$ 空间中以原点为心、$r$ 为半径的开球，$z'=(z_1,\cdots,z_{n-1})$。

先证 $\gamma(0)=0$。由于 $0\in\partial U$，有 $0\notin U$。若$0>\gamma(0)$，由 \eqref{eq:ch05-统一-边界} 得 $0\in U$，矛盾！若 $\gamma(0)>0$，由 $0\in\overline{U}=U\cup \partial U$，可取 $z_l\in U\cap B_z(0,r)$，$z_l\to0$。由 \eqref{eq:ch05-统一-边界}，$(z_l)_n>\gamma(z_l')$；令 $l\to\infty$，由 $\gamma$ 连续得 $0\geqslant\gamma(0)$，矛盾！故
\begin{align}
\gamma(0)=0.
\label{eq:ch05-统一-gamma0}
\end{align}
再证 $\nabla\gamma(0)=0$(详见\hyperref[曲面的单位法向量选取]{注})。曲面 $\{z_n=\gamma(z')\}$ 在 $0$ 处的单位法向量为
\begin{align*}
\frac{(-\nabla\gamma(0),1)}{\sqrt{1+|\nabla\gamma(0)|^{2}}}.
\end{align*}
由 \eqref{eq:ch05-统一-边界}，$U$ 位于该曲面之上（$z_n>\gamma(z')$），故其内法向量就是该法向量；而内法向量为 $e_n$。于是
\begin{align*}
\frac{(-\nabla\gamma(0),1)}{\sqrt{1+|\nabla\gamma(0)|^{2}}}=e_n.
\end{align*}
比较前 $n-1$ 个分量得
\begin{align}
\nabla\gamma(0)=0.
\label{eq:ch05-统一-gradgamma0}
\end{align}

定义
\begin{align*}
\Phi_1(z)\triangleq\left(z',z_n-\gamma(z')\right),\qquad
\Psi_1(y)\triangleq\left(y',y_n+\gamma(y')\right).
\end{align*}
直接计算得
\begin{align*}
\Phi_1(\Psi_1(y))&=\left(y',y_n+\gamma(y')-\gamma(y')\right)=y,\\
\Psi_1(\Phi_1(z))&=\left(z',z_n-\gamma(z')+\gamma(z')\right)=z.
\end{align*}
故 $\Phi_1,\Psi_1$ 互逆。由于 $\gamma\in C^{k}(\mathbb{R}^{n-1})$，因此$\Phi_1,\Psi_1\in C^{k}(\mathbb{R}^{n})$，且互为 $C^{k}$ 逆映射，从而都是 $C^{k}$ 微分同胚。其 Jacobi 矩阵为
\begin{align*}
D\Phi_1(z)=
\begin{pmatrix}
I_{n-1} & 0\\
-\nabla\gamma(z')^{T} & 1
\end{pmatrix},
\qquad
D\Psi_1(y)=
\begin{pmatrix}
I_{n-1} & 0\\
\nabla\gamma(y')^{T} & 1
\end{pmatrix}.
\end{align*}
两者均为三角矩阵，对角元全为 $1$，故
\begin{align}
|\det D\Phi_1(z)|=|\det D\Psi_1(y)|=1.
\label{eq:ch05-统一-det1}
\end{align}

\noindent$\mathbf{Step}\bm{3}$.{\heiti 选取坐标球与边界邻域 $W$，并验证 $W\subset\subset V$.}

令 $y_0\triangleq\Phi_1(0)=0$。由于 $\Phi_1(B_z(0,r))$ 是含 $0$ 的开集，存在 $\rho_0>0$，使得
\begin{align*}
B(0,\rho_0)\subseteq\Phi_1(B_z(0,r)).
\end{align*}
又 $\Psi_1(0)=0$ 且 $\Psi_1$ 在 $0$ 处连续，故存在 $0<\rho_1\leqslant\rho_0$，使得
\begin{align}
|y|<\rho_1\Longrightarrow|\Psi_1(y)|<r.
\label{eq:ch05-统一-rho1}
\end{align}
由于 $x_0\in\partial U\subseteq\overline{U}\subseteq V$ 且 $V$ 是开集，存在 $\delta>0$ 使得
\begin{align}
B(x_0,\delta)\subseteq V.
\label{eq:ch05-统一-delta}
\end{align}
由 $\Psi_1$ 在 $0$ 处的连续性，存在 $0<\rho\leqslant\rho_1$，使得
\begin{align}
|y|<\rho\Longrightarrow|\Psi_1(y)|<\frac{\delta}{2}.
\label{eq:ch05-统一-rho}
\end{align}
记
\begin{gather*}
B\triangleq B(0,\rho),\qquad
W\triangleq x_0+Q^{T}\Psi_1(B),\\
B^{+}\triangleq B\cap\{y_n>0\},\quad
B^{-}\triangleq B\cap\{y_n<0\},\quad
B^{0}\triangleq B\cap\{y_n=0\}.
\end{gather*}
由于 $\Psi_1$ 是 $C^{k}$ 微分同胚，$W$ 是开集。对 $x=x_0+Q^{T}\Psi_1(y)\in W$，由 $Q^{T}$ 正交及 \eqref{eq:ch05-统一-rho}，
\begin{align*}
|x-x_0|=|Q^{T}\Psi_1(y)|=|\Psi_1(y)|<\frac{\delta}{2}<r,
\end{align*}
故
\begin{align}
W\subseteq B\left( x_0,\frac{\delta}{2} \right)\Longrightarrow \overline{W}\subseteq\overline{B\left(x_0,\frac{\delta}{2}\right)}
\subseteq B(x_0,\delta)\subseteq V.
\label{eq:ch05-统一-W紧含}
\end{align}
又 $\overline{W}$ 是闭且有界集，从而是紧集，故
\begin{align}
W\subset\subset V.
\label{eq:ch05-统一-W紧含2}
\end{align}

\noindent$\mathbf{Step}\bm{4}$.{\heiti 定义拉直映射 $\Phi,\Psi$，验证其为 $C^{k}$ 微分同胚且 Jacobi 行列式为 $1$.}

定义
\begin{align*}
\Phi:W\to B,\qquad \Phi(x)\triangleq\Phi_1(Q(x-x_0)),
\end{align*}
以及
\begin{align*}
\Psi:B\to W,\qquad \Psi(y)\triangleq x_0+Q^{T}\Psi_1(y).
\end{align*}
对 $x\in W$，由 $QQ^{T}=I$，
\begin{align*}
\Psi(\Phi(x))=x_0+Q^{T}\Psi_1(\Phi_1(Q(x-x_0)))=x_0+Q^{T}Q(x-x_0)=x.
\end{align*}
对 $y\in B$，由 $Q^{T}Q=I$，
\begin{align*}
\Phi(\Psi(y))=\Phi_1(Q(x_0+Q^{T}\Psi_1(y)-x_0))=\Phi_1(QQ^{T}\Psi_1(y))=\Phi_1(\Psi_1(y))=y.
\end{align*}
故 $\Phi,\Psi$ 互逆。由于 $Q$ 是正交矩阵、$\Phi_1,\Psi_1\in C^{k}(\mathbb{R}^{n})$，故 $\Phi\in C^{k}(W),\Psi\in C^{k}(B)$，且互为 $C^{k}$ 逆映射，从而是 $C^{k}$ 微分同胚。由\hyperref[theorem:向量求导基本性质]{向量求导基本性质\ref{theorem:向量求导基本性质-4}}，得
\begin{align}
D\Phi(x)=D\Phi_1(Q(x-x_0))Q,\qquad
D\Psi(y)=Q^{T}D\Psi_1(y).
\label{eq:ch05-统一-DPhiD-Psi}
\end{align}
由 \eqref{eq:ch05-统一-det1} 及 $|\det Q|=|\det Q^{T}|=1$，
\begin{align}
|\det D\Phi(x)|
&=|\det D\Phi_1(Q(x-x_0))|\cdot|\det Q|=1,
\label{eq:ch05-统一-detPhi}\\
|\det D\Psi(y)|
&=|\det Q^{T}|\cdot|\det D\Psi_1(y)|=1.
\label{eq:ch05-统一-detPsi}
\end{align}
又由 $\Phi,\Psi$ 互逆，
\begin{align}
\Phi(W)=B,\qquad\Psi(B)=W.
\label{eq:ch05-统一-PhiW}
\end{align}

\begin{figure}[H]
\centering
\includegraphics[width=0.8\textwidth]{8e1ffee58f4e42d8a67b7ea75f26d99d.png}
\caption{拉直边界}
\label{figure:ch05-拉直边界}
\end{figure}

\noindent$\mathbf{Step}\bm{5}$.{\heiti 验证 $\Phi$ 把 $W\cap U$ 映为 $B^{+}$，把 $W\cap\partial U$ 映为 $B^{0}$.}

先证 $\Phi(W\cap U)=B^{+}$。
若 $x\in W\cap U$，则由\eqref{eq:ch05-统一-W紧含}式知$W\subseteq B\left( x_0,\frac{\delta}{2} \right)$ 及 $\frac{\delta}{2}<r$,从而$z=Q(x-x_0)$满足$|z|=|x-x_0|<r$，故 $z\in B_z(0,r)$.,又由 \eqref{eq:ch05-统一-边界}，$z_n>\gamma(z')$。于是
\begin{align*}
\Phi_n(x)=(\Phi_1(z))_n=z_n-\gamma(z')>0,
\end{align*}
即 $\Phi(x)\in B^{+}$。故 $\Phi(W\cap U)\subseteq B^{+}$。

反之，若 $y\in B^{+}$，令 $x=\Psi(y)\in W$。则
\begin{align*}
z=Q(x-x_0)=Q(x_0+Q^{T}\Psi_1(y)-x_0)=\Psi_1(y).
\end{align*}
于是 $z'=y'$，$z_n=y_n+\gamma(y')$，从而
\begin{align*}
z_n=y_n+\gamma(y')>\gamma(y')=\gamma(z').
\end{align*}
由 \eqref{eq:ch05-统一-边界} 得,在$z$坐标系下有$z\in U\cap B_z(0,r)$,从而在$x$坐标系下有$x\in U$。故 $B^{+}\subseteq\Phi(W\cap U)$。因此
\begin{align}
\Phi(W\cap U)=B^{+}.
\label{eq:ch05-统一-PhiWU}
\end{align}

再证 $\Phi(W\cap\partial U)=B^{0}$。
若 $x\in W\cap\partial U$，则 $z=Q(x-x_0)\in B_z(0,r)$ 且 $x\notin U$。由 \eqref{eq:ch05-统一-边界}，$x\notin U$ 等价于 $z_n\leqslant\gamma(z')$；又因为 $x\in\partial U\subseteq\overline U$，所以存在点列 $x_j\in U$，使得 $x_j\to x$，则 $z_j\to z$。由于 $z\in B_z(0,r)$ 且 $B_z(0,r)$ 是开集，当 $j$ 充分大时,有 $z_j\in B_z(0,r)$。由\eqref{eq:ch05-统一-边界}知
$z_{j,n}>\gamma(z_j').$
令 $j\to\infty$，由 $\gamma$ 的连续性，得$z_n\geqslant\gamma(z').$
故 $z_n=\gamma(z')$，于是
\begin{align*}
\Phi_n(x)=z_n-\gamma(z')=0,
\end{align*}
即 $\Phi(x)\in B^{0}$。故 $\Phi(W\cap\partial U)\subseteq B^{0}$。

反之，若 $y\in B^{0}$，令 $x=\Psi(y)\in W$,则 $z=Q(x-x_0)=\Psi_1(y)$ 满足 $z_n=y_n+\gamma(y')=\gamma(y')=\gamma(z')$。由 \eqref{eq:ch05-统一-边界}知$x\notin U$.又取 $y_l\in B^{+}$，$y_l\to y$，则由\eqref{eq:ch05-统一-PhiWU}式知$x_l=\Psi(y_l)\in W\cap U$，$x_l\to x$，故 $x\in\overline{U}$。于是 $x\in\partial U$。因此 $B^{0}\subseteq\Phi(W\cap\partial U)$。所以
\begin{align}
\Phi(W\cap\partial U)=B^{0}.
\label{eq:ch05-统一-PhiWbU}
\end{align}

\noindent$\mathbf{Step}\bm{6}$.{\heiti 在拉直坐标下作局部反射延拓，并建立局部 $W^{k,p}$ 估计.}

设 $u\in C^{\infty}(\overline{U})$。定义
\begin{align*}
u'(y)\triangleq u(\Psi(y)),\qquad y\in B^{+}.
\end{align*}
由于 $\Psi\in C^{k}(\mathbb{R}^n)$,且 $\Psi(\overline{B^{+}})\subseteq\overline{U}$，故$u'\in C^{k}(\overline{B^{+}})$。

对 $u'$ 应用 $\mathbf{Step}\bm{1}$ 的高阶反射（即用 \eqref{eq:ch05-统一-反射}式，把 $u$ 换成 $u'$,把 $B_{r'}$ 换成 $B$），得到 $\overline{u}'\in W^{k,p}(B)$，且由\eqref{eq:ch05-统一-反射估计}式知
\begin{align}
\overline{u}'(y)=u'(y)\quad(y\in B^{+}),\qquad
\|\overline{u}'\|_{W^{k,p}(B)}\leqslant C_0\|u'\|_{W^{k,p}(B^{+})}.
\label{eq:ch05-统一-局部反射}
\end{align}

定义
\begin{align*}
\overline{u}(x)\triangleq\overline{u}'(\Phi(x)),\qquad x\in W.
\end{align*}
由\cref{lemma:ch05-复合Wkp}知，$\overline{u}\in W^{k,p}(W)$，且存在只依赖于 $\Phi,\Psi,n,k,p$ 的常数 $C_1$，使得
\begin{align}
\|\overline{u}\|_{W^{k,p}(W)}
&\leqslant C_1\|\overline{u}'\|_{W^{k,p}(B)}
\leqslant C_1C_0\|u'\|_{W^{k,p}(B^{+})}
\notag\\
&\leqslant C_1C_0\cdot C_1\|u\|_{W^{k,p}(W\cap U)}
=C_0'\|u\|_{W^{k,p}(W\cap U)},
\label{eq:ch05-统一-局部估计}
\end{align}
其中 $C_0'\triangleq C_0C_1^{2}$，第二、第三个不等号都用了\cref{lemma:ch05-复合Wkp} 中的\eqref{eq:ch05-复合估计}式。若 $x\in W\cap U$，则由 \eqref{eq:ch05-统一-PhiWU}式，$\Phi(x)\in B^{+}$，于是
\begin{align*}
\overline{u}(x)=\overline{u}'(\Phi(x))=u'(\Phi(x))=u(\Psi(\Phi(x)))=u(x).
\end{align*}
故
\begin{align}
\overline{u}|_{W\cap U}=u.
\label{eq:ch05-统一-局部等于u}
\end{align}

\noindent$\mathbf{Step}\bm{7}$.{\heiti 取有限覆盖与辅助开集 $W_0$.}

由 $\partial U$ 紧及 $\mathbf{Step}\bm{2}$--$\bm{6}$，存在有限个点 $x_i^{0}\in\partial U$ 与开集 $W_i\subset\subset V$（$i=1,\cdots,N$），以及 $u$ 到 $W_i$ 的延拓 $\overline{u}_i\in W^{k,p}(W_i)$，以及只依赖于 $\Phi,\Psi,n,k,p$ 的常数 $C_i'$,使得
\begin{gather}
\partial U\subseteq\bigcup_{i=1}^{N}W_i,\quad \overline{u}_i|_{W_i\cap U}=u,\nonumber\\
\|\overline{u}_i\|_{W^{k,p}(W_i)}\leqslant C'_i\|u\|_{W^{k,p}(U)}\leqslant C'\|u\|_{W^{k,p}(U)},\quad i=1,\cdots,N,\label{eq:ch05-统一-覆盖估计}
\end{gather}
其中 $C'\triangleq \underset{i=1,\cdots ,N}{\max}\{C_{i}^{\prime}\}$.

令
\begin{align}
K\triangleq U\setminus\bigcup_{i=1}^{N}W_i.
\label{eq:ch05-统一-K}
\end{align}
先证 $\overline{K}\subseteq U$。若存在 $x\in\overline{K}\cap\partial U$，则由 $\partial U\subseteq\bigcup_{i=1}^{N}W_i$ 知存在 $j\in\{1,\cdots,N\}$ 使 $x\in W_j$。由于 $W_j$ 是开集，存在 $\delta>0$ 使 $B(x,\delta)\subseteq W_j$。又 $x\in\overline{K}$，故 $B(x,\delta)\cap K\neq\varnothing$。但
\begin{align*}
K\subseteq U\setminus\bigcup_{i=1}^{N}W_i\subseteq\mathbb{R}^{n}\setminus W_j,
\end{align*}
与 $B(x,\delta)\cap K\neq\varnothing$ 矛盾！故 $\overline{K}\cap\partial U=\varnothing$。又由 \eqref{eq:ch05-统一-K} 及\subcref{Real Analysis-theorem:边界和开集基本性质}{Real Analysis-theorem:边界和开集基本性质-5}知
\begin{align*}
\overline{K}\subseteq\overline{U}=U\cup\partial U.
\end{align*}
结合 $\overline{K}\cap\partial U=\varnothing$ 得 $\overline{K}\subseteq U$。由于 $\overline{K}$ 是闭集且 $\overline{K}\subseteq\overline{U}$，而 $\overline{U}$ 紧，故 $\overline{K}$ 紧。

若 $U=\varnothing$，则定理平凡，取 $E_k=0$ 即可。以下设 $U\neq\varnothing$。

若 $K=\varnothing$，任取 $x_*\in U$ 及 $r>0$ 使 $B(x_*,r)\subset\subset U$，令
\begin{align*}
W_0\triangleq B(x_*,r).
\end{align*}
此时 $U\subseteq\bigcup_{i=1}^{N}W_i$，故
\begin{align*}
\overline{U}\subseteq\bigcup_{i=1}^{N}W_i\subseteq\bigcup_{i=0}^{N}W_i.
\end{align*}

若 $K\neq\varnothing$，由于 $\overline{K}$ 紧、$\overline{K}\subseteq U$，而 $\mathbb{R}^{n}\setminus U$ 是闭集且与 $\overline{K}$ 不交，由\cref{proposition:紧集的相对紧开邻域}知，存在开集$W_0$,使得
\begin{align*}
K\subseteq\overline{K}\subseteq W_0\subseteq\overline{W_0}\subseteq U,
\end{align*}
且 $\overline{W_0}$ 紧，故 $W_0\subset\subset U$。

最后验证 $\overline{U}\subseteq\bigcup_{i=0}^{N}W_i$。任取 $x\in\overline{U}$。若 $x\in\bigcup_{i=1}^{N}W_i$，则已属于右端。若 $x\notin\bigcup_{i=1}^{N}W_i$，则 $x\in\overline{U}\setminus\bigcup_{i=1}^{N}W_i=U\setminus\bigcup_{i=1}^{N}W_i=K\subseteq W_0$。故总有 $x\in\bigcup_{i=0}^{N}W_i$。

\noindent$\mathbf{Step}\bm{8}$.{\heiti 取单位分解并对光滑 $u$ 建立整体估计.}

由\hyperref[corollary:紧集上的光滑单位分解定理]{紧集上的光滑单位分解定理}（取 $K=\overline{U}$，$V_i=W_i$），存在 $\zeta_i\in C_c^{\infty}(W_i)$，使得
\begin{align*}
0\leqslant\zeta_i\leqslant1,\qquad
\sum_{i=0}^{N}\zeta_i(x)=1\quad(x\in\overline{U}).
\end{align*}
对每个 $i=0,1,\dots,N$，将 $\zeta_i\overline{u}_i$ 零延拓到 $\mathbb{R}^n$，仍记为 $\zeta_i\overline{u}_i$,即
\begin{align*}
(\zeta_i\overline{u}_i)(x)\triangleq
\begin{cases}
\zeta_i(x)\overline{u}_i(x), & x\in W_i,\\[2mm]
0, & x\in\mathbb{R}^n\setminus W_i,
\end{cases}
\end{align*}
其中 $\overline{u}_0\triangleq u$.令
\begin{align*}
E_0u\triangleq\sum_{i=0}^{N}\zeta_i\overline{u}_i,
\end{align*}

下面证明 $E_0u\in W^{k,p}(\mathbb{R}^n)$。对每个 $i=0,1,\cdots,N$，记
$v_i\triangleq\zeta_i\overline{u}_i$,则由\hyperref[theorem:ch05-弱导数性质]{Leibniz公式}知,$v_i\in W^{k,p}(W_i)$,且
\begin{align*}
\operatorname{spt}v_i\subseteq\operatorname{spt}\zeta_i\subset\subset W_i.
\end{align*}
取 $\chi_i\in C_c^\infty(W_i)$，使得 $\chi_i=1$ 于包含 $\operatorname{spt}v_i$的一个开集内。对任意 $\varphi\in C_c^\infty(\mathbb{R}^n)$，有 $\chi_i\varphi\in C_c^\infty(W_i)$，且由 由\hyperref[theorem:ch05-弱导数性质]{Leibniz公式}可得,在$\operatorname{spt}v_i$内都有
\begin{align*}
D^{\alpha}(\chi _i\varphi )=\sum_{\beta \leqslant \alpha}{\binom{\alpha}{\beta}D^{\beta}\chi _i\,D^{\alpha -\beta}\varphi}=\chi _iD^{\alpha}\varphi +\sum_{|\beta |\geqslant 1,\beta \leqslant \alpha}{\binom{\alpha}{\beta}D^{\beta}\chi _i\,D^{\alpha -\beta}\varphi}=D^{\alpha}\varphi .
\end{align*}
而在$\operatorname{spt}v_i$外都有$v_i=0$,故
\begin{align*}
v_iD^\alpha\varphi(x)=v_iD^\alpha(\chi_i\varphi)(x),\quad x\in \mathbb{R}^n.
\end{align*}
于是由 $v_i\in W^{k,p}(W_i)$ 的弱导数定义，得
\begin{align*}
\int_{\mathbb{R}^n}v_iD^\alpha\varphi\,\mathrm{d}x
=\int_{W_i}v_iD^\alpha(\chi_i\varphi)\,\mathrm{d}x=(-1)^{|\alpha|}\int_{W_i}D^\alpha v_i\,\chi_i\varphi\,\mathrm{d}x=(-1)^{|\alpha|}\int_{W_i}D^\alpha v_i\varphi\,\mathrm{d}x=(-1)^{|\alpha|}\int_{\mathbb{R}^n}\widetilde{D^\alpha v_i}\,\varphi\,\mathrm{d}x,
\end{align*}
其中 $\widetilde{D^\alpha v_i}$ 表示 $D^\alpha v_i$ 到 $\mathbb{R}^n$ 的零延拓。这说明 $\widetilde{D^\alpha v_i}$ 是 $v_i$ 的 $\alpha$ 阶弱导数。因此 $v_i\in W^{k,p}(\mathbb{R}^n)$，且
\begin{align*}
\|v_i\|_{W^{k,p}(\mathbb{R}^n)}=\|\zeta_i\overline{u}_i\|_{W^{k,p}(W_i)}.
\end{align*}
由于
\begin{align*}
E_0u=\sum\limits_{i=0}^{N}v_i
\end{align*}
是有限和，而 $W^{k,p}(\mathbb{R}^n)$ 是线性空间，故
\begin{align*}
E_0u\in W^{k,p}(\mathbb{R}^n).
\end{align*}

又由\hyperref[theorem:ch05-弱导数性质]{Leibniz公式}知 $\zeta_i\overline{u}_i\in W^{k,p}(W_i)$,且对 $|\alpha|\leqslant k$，有
\begin{align*}
D^{\alpha}(\zeta_i\overline{u}_i)=\sum_{\beta\leqslant\alpha}\binom{\alpha}{\beta}D^{\beta}\zeta_i\,D^{\alpha-\beta}\overline{u}_i.
\end{align*}
由于 $|\alpha-\beta|\leqslant k$，取$M_i\triangleq \underset{\left| \alpha \right|\leqslant k}{\max}\sum_{\beta \leqslant \alpha}{\binom{\alpha}{\beta}}\left\| D^{\beta}\zeta _i \right\| _{L^{\infty}\left( W_i \right)}$,则$M_i$只依赖于 $\zeta_i,n,k$，且
\begin{align*}
|D^{\alpha}(\zeta_i\overline{u}_i)(x)|\leqslant M_i\sum_{|\gamma|\leqslant k}|D^{\gamma}\overline{u}_i(x)|,\qquad x\in W_i.
\end{align*}
于是由\hyperref[Basis of Analytics-proposition:凸性不等式]{凸性不等式}得
\begin{align*}
\left\| \zeta _i\overline{u}_i \right\| _{W^{k,p}\left( W_i \right)}^{p}&=\sum_{|\alpha |\leqslant k}{\int_{W_i}{\left| D^{\alpha}\left( \zeta _i\overline{u}_i \right) \left( x \right) \right|^p\mathrm{d}x}}\leqslant M_{i}^{p}N_k\int_{W_i}{\left( \sum_{\left| \gamma \right|\leqslant k}{\left| D^{\gamma}\overline{u}_i\left( x \right) \right|} \right) ^p\mathrm{d}x}
\\
&\leqslant M_{i}^{p}N_{k}^{p}\sum_{|\gamma |\leqslant k}{\int_{W_i}{|D^{\gamma}\overline{u}_i|^p\,\mathrm{d}x}}=M_{i}^{p}N_{k}^{p}\left\| \overline{u}_i \right\|_{W^{k,p}(W_i)}^{p},
\end{align*}
其中$N_k$ 为满足 $|\alpha|\leqslant k$ 的 $\alpha$ 的个数.
记$M_{i}^{\prime}\triangleq M_iN_k$，得
\begin{align}
\|\zeta_i\overline{u}_i\|_{W^{k,p}(W_i)}\leqslant M_i'\|\overline{u}_i\|_{W^{k,p}(W_i)}.
\label{eq:ch05-统一-zeta}
\end{align}
由 \eqref{eq:ch05-统一-覆盖估计} 与 \eqref{eq:ch05-统一-zeta}得，对$u\in C^{\infty}(\overline{U})$,有
\begin{align}
\|E_0u\|_{W^{k,p}(\mathbb{R}^{n})}
&\leqslant\sum_{i=0}^{N}\|\zeta_i\overline{u}_i\|_{W^{k,p}(W_i)}
\leqslant\sum_{i=0}^{N}M_i'\|\overline{u}_i\|_{W^{k,p}(W_i)}
\notag\\
&\leqslant\Big(M_0'+C'\sum_{i=1}^{N}M_i'\Big)\|u\|_{W^{k,p}(U)}
= C''\|u\|_{W^{k,p}(U)},
\label{eq:ch05-统一-smooth估计}
\end{align}
其中$C''\triangleq M_0'+C'\sum_{i=1}^{N}M_i'$.由 $\sum\limits_{i=0}^{N}\zeta_i=1$ 于 $\overline{U}$ 以及 $\overline{u}_i=u$ 于 $W_i\cap U$，可推出 $E_0u=u$ 于 $U$。事实上，对几乎处处的 $x\in U$，因为 $\operatorname{spt}\zeta_i\subseteq W_i$，所以
\begin{align*}
\zeta_i(x)\overline{u}_i(x)=
\begin{cases}
\zeta_i(x)u(x), & x\in W_i\cap U,\\
0=\zeta_i(x)u(x), & x\in U\setminus W_i,
\end{cases}
\end{align*}
其中第一行用了 $\overline{u}_i=u$ 于 $W_i\cap U$，第二行用了 $\zeta_i=0$ 于 $U\setminus W_i$。当 $i=0$ 时 $\overline{u}_0=u$ 于 $U$，上述结论同样适用。于是
\begin{align*}
E_0u(x)
=\sum\limits_{i=0}^{N}\zeta_i(x)\overline{u}_i(x)
=\sum\limits_{i=0}^{N}\zeta_i(x)u(x)=u(x)\sum\limits_{i=0}^{N}\zeta_i(x)
=u(x).
\end{align*}
故
\begin{align}
E_0u=u\quad\text{于 }U.
\label{eq:ch05-统一-等于u}
\end{align}
由 $\operatorname{spt}\zeta_i\subseteq W_i$，有
\begin{align}
\operatorname{spt}E_0u\subseteq\bigcup_{i=0}^{N}W_i\subset\subset V.
\label{eq:ch05-统一-支集}
\end{align}

\noindent$\mathbf{Step}\bm{9}$.{\heiti $1\leqslant p<\infty$：用 $C^{\infty}(\overline{U})$ 逼近定义 $E_k$.}

设 $u\in W^{k,p}(U)$。由\hyperref[theorem:ch05-直到边界光滑的函数整体逼近定理]{直到边界光滑的函数整体逼近定理}，存在 $u_m\in C^{\infty}(\overline{U})$，使得
\begin{align*}
\|u_m-u\|_{W^{k,p}(U)}\longrightarrow0\quad(m\to\infty).
\end{align*}
由\eqref{eq:ch05-统一-smooth估计}式得
\begin{align*}
\|E_0u_m-E_0u_l\|_{W^{k,p}(\mathbb{R}^{n})}
\leqslant C''\|u_m-u_l\|_{W^{k,p}(U)}\longrightarrow0\quad(m,l\to\infty),
\end{align*}
故 $\{E_0u_m\}$ 是 $W^{k,p}(\mathbb{R}^{n})$ 中的 Cauchy 列。由\hyperref[theorem:ch05-WkpBanach]{Sobolev空间都是Banach空间}知，存在 $v\in W^{k,p}(\mathbb{R}^{n})$ 使 $E_0u_m\to v$。定义
\begin{align*}
E_ku\triangleq v.
\end{align*}
若 $v_m\in C^{\infty}(\overline{U})$ 也满足 $v_m\to u$ 于 $W^{k,p}(U)$，则由
\begin{align*}
\|E_0u_m-E_0v_m\|_{W^{k,p}(\mathbb{R}^{n})}
\leqslant C''\|u_m-v_m\|_{W^{k,p}(U)}\longrightarrow0
\end{align*}
可得极限的唯一性，故$E_ku$ 与逼近序列选取无关。

由 $E_0$ 的线性性与极限的线性性知，$E_k$ 线性，且
\begin{align*}
\|E_ku\|_{W^{k,p}(\mathbb{R}^{n})}
=\lim_{m\to\infty}\|E_0u_m\|_{W^{k,p}(\mathbb{R}^{n})}
\leqslant C''\|u\|_{W^{k,p}(U)},
\end{align*}
故 $E_k$ 有界。

又由 \eqref{eq:ch05-统一-等于u}式，$E_0u_m=u_m$ 于 $U$；而 $u_m\to u$ 于 $W^{k,p}(U)$ 蕴含 $u_m\to u$ 于 $L^{p}(U)$，故
\begin{align*}
E_0u_m|_{U}=u_m\longrightarrow u\quad\text{于 }L^{p}(U).
\end{align*}
另一方面，$E_0u_m\to E_ku$ 于 $W^{k,p}(\mathbb{R}^{n})$，故 $E_0u_m|_{U}\to E_ku|_{U}$ 于 $L^{p}(U)$。由 $L^{p}(U)$ 中极限的唯一性知
\begin{align*}
E_ku|_U=u\quad\text{于 }L^{p}(U),
\end{align*}
即 $E_ku=u$ a.e. 于 $U$。

最后证支集。由\eqref{eq:ch05-统一-支集}式，每个 $E_0u_m$ 的支集含于
\begin{align*}
K_*\triangleq\bigcup_{i=0}^{N}\overline{W_i}\subset\subset V.
\end{align*}
故对每个 $m$，$E_0u_m=0$ a.e. 于 $\mathbb{R}^{n}\setminus K_*$。由 $E_0u_m\to E_ku$ 于 $W^{k,p}(\mathbb{R}^{n})$，限制在 $\mathbb{R}^{n}\setminus K_*$ 上得
\begin{align*}
E_0u_m|_{\mathbb{R}^{n}\setminus K_*}\longrightarrow E_ku|_{\mathbb{R}^{n}\setminus K_*}\quad\text{于 }L^{p}(\mathbb{R}^{n}\setminus K_*).
\end{align*}
但左端恒为零，故 $E_ku|_{\mathbb{R}^{n}\setminus K_*}=0$，即 $E_ku=0$ a.e. 于 $\mathbb{R}^{n}\setminus K_*$。因此
\begin{align*}
\operatorname{spt}E_ku\subseteq K_*\subset\subset V.
\end{align*}
故 $1\leqslant p<\infty$ 时定理成立。

\noindent$\mathbf{Step}\bm{10}$.{\heiti $p=\infty$：用 $W^{k,\infty}$ 的 $C^{k-1}$ 代表元与第 $k$ 阶反射.}

设 $u\in W^{k,\infty}(U)$。沿用 $\mathbf{Step}\bm{3}$--$\mathbf{8}$ 中固定下来的 $W_i,\Phi_i,\Psi_i,\zeta_i$（$i=1,\cdots,N$）及 $W_0$，且$B_i$ 与 $B_i^{+}$ 是第 $i$ 个边界点对应的拉直球与上半球。对 $i=1,\cdots,N$，令
\begin{align*}
u_i'\triangleq u\circ\Psi_i.
\end{align*}
由\hyperref[theorem:ch05-弱导数性质]{弱导数性质\ref{theorem:ch05-弱导数性质-3}}，$u|_{W_i\cap U}\in W^{k,\infty}(W_i\cap U)$。由\cref{lemma:ch05-复合Wkp}（取 $p=\infty$），存在只依赖于 $\Psi,n,k,p$ 的常数$C_2>0$,使得
\begin{align}\label{qeNHJJHSAJ}
u_i'\in W^{k,\infty}(B_i^{+}),\quad
\|u_i'\|_{W^{k,\infty}(B_i^{+})}\leqslant C_2\|u\|_{W^{k,\infty}(W_i\cap U)}\leqslant C_2\|u\|_{W^{k,\infty}(U)}.
\end{align}
由\cref{lemma:ch05-Wkinfty代表元}，$u_i'$ 有代表元 $\widetilde{u_i'}\in C^{k-1}(\overline{B_i^{+}})$，其直到 $k-1$ 阶的导数在$B_i^+$上Lipschitz连续。按\eqref{eq:ch05-统一-反射}式用 $\widetilde{u_i'}$ 定义 $\overline{u_i'}$。由\cref{lemma:ch05-高阶反射}和\eqref{eq:ch05-统一-反射估计}\eqref{qeNHJJHSAJ}式得
\begin{align}\label{eqrNJJKKANZZZZ}
\overline{u_i'}\in W^{k,\infty}(B_i),\qquad
\left\| \overline{u_{i}^{\prime}} \right\| _{W^{k,\infty}\left( B_i \right)}\leqslant C_0\left\| u_{i}^{\prime} \right\| _{W^{k,\infty}(B_{i}^{+})}\leqslant C_0C_2\left\| u \right\| _{W^{k,\infty}\left( U \right)}.
\end{align}
再令
\begin{align*}
\overline{u}_i\triangleq\overline{u_i'}\circ\Phi_i\in W^{k,\infty}(W_i),
\end{align*}
则再由\eqref{eqrNJJKKANZZZZ}式和\cref{lemma:ch05-复合Wkp}（取 $p=\infty$），存在只依赖于 $\Phi,n,k,p$ 的常数$C_3>0$,使得
\begin{align}
\left\| \overline{u}_i \right\| _{W^{k,\infty}\left( W_i \right)}\leqslant C_3\left\| \overline{u_{i}^{\prime}} \right\| _{W^{k,\infty}\left( B_i \right)}\leqslant C_0C_2C_3\left\| u \right\| _{W^{k,\infty}\left( U \right)},\label{eqnjkNJNS}
\end{align}
且若 $x\in W_i\cap U$，则 $\Phi_i(x)\in B_i^{+}$，故
\begin{align*}
\overline{u}_i(x)=\overline{u_i'}(\Phi_i(x))=u_i'(\Phi_i(x))=u(\Psi_i(\Phi_i(x)))=u(x),
\end{align*}
即 $\overline{u}_i|_{W_i\cap U}=u$。
特别地,令 $\overline{u}_0\triangleq u$。由\hyperref[theorem:ch05-弱导数性质]{Leibniz公式}知 $\zeta_i\overline{u}_i\in W^{k,\infty}(W_i)$,且对 $|\alpha|\leqslant k$，有
\begin{align*}
D^{\alpha}(\zeta_i\overline{u}_i)=\sum_{\beta\leqslant\alpha}\binom{\alpha}{\beta}D^{\beta}\zeta_i\,D^{\alpha-\beta}\overline{u}_i.
\end{align*}，
从而对每个 $i=0,1,\cdots,N$ 及 $|\alpha|\leqslant k$，有
\begin{align*}
\|D^{\alpha}(\zeta_i\overline{u}_i)\|_{L^{\infty}(W_i)}
\leqslant M_i\sum_{|\gamma|\leqslant k}\|D^{\gamma}\overline{u}_i\|_{L^{\infty}(W_i)},
\end{align*}
于是
\begin{align}
\|\zeta_i\overline{u}_i\|_{W^{k,\infty}(W_i)}\leqslant M_i\|\overline{u}_i\|_{W^{k,\infty}(W_i)}.\label{eqnjkNJNS1}
\end{align}
令
\begin{align*}
E_ku\triangleq\overline{u}\triangleq\sum_{i=0}^{N}\zeta_i\overline{u}_i.
\end{align*}

下面类似$\mathbf{Step8}$证明 $\overline{u}\in W^{k,\infty}(\mathbb{R}^n)$。对每个 $i=0,1,\cdots,N$，记$v_i\triangleq\zeta_i\overline{u}_i,$则$v_i\in W^{k,\infty}(W_i)$,且由 $\zeta_i\in C_c^\infty(W_i)$ 知
\begin{align*}
\operatorname{spt}v_i\subseteq\operatorname{spt}\zeta_i\subset\subset W_i.
\end{align*}
取 $\chi_i\in C_c^\infty(W_i)$，使得 $\chi_i=1$ 于包含 $\operatorname{spt}v_i$的一个开集内。对任意 $\varphi\in C_c^\infty(\mathbb{R}^n)$，有 $\chi_i\varphi\in C_c^\infty(W_i)$，且由 由\hyperref[theorem:ch05-弱导数性质]{Leibniz公式}可得,在$\operatorname{spt}v_i$内都有
\begin{align*}
D^{\alpha}(\chi _i\varphi )=\sum_{\beta \leqslant \alpha}{\binom{\alpha}{\beta}D^{\beta}\chi _i\,D^{\alpha -\beta}\varphi}=\chi _iD^{\alpha}\varphi +\sum_{|\beta |\geqslant 1,\beta \leqslant \alpha}{\binom{\alpha}{\beta}D^{\beta}\chi _i\,D^{\alpha -\beta}\varphi}=D^{\alpha}\varphi .
\end{align*}
而在$\operatorname{spt}v_i$外都有$v_i=0$,故
\begin{align*}
v_iD^\alpha\varphi(x)=v_iD^\alpha(\chi_i\varphi)(x),\quad x\in \mathbb{R}^n.
\end{align*}
于是由 $v_i\in W^{k,\infty}(W_i)$ 的弱导数定义，得
\begin{align*}
\int_{\mathbb{R}^n}v_iD^\alpha\varphi\,\mathrm{d}x
=\int_{W_i}v_iD^\alpha(\chi_i\varphi)\,\mathrm{d}x=(-1)^{|\alpha|}\int_{W_i}D^\alpha v_i\,\chi_i\varphi\,\mathrm{d}x=(-1)^{|\alpha|}\int_{W_i}D^\alpha v_i\varphi\,\mathrm{d}x=(-1)^{|\alpha|}\int_{\mathbb{R}^n}\widetilde{D^\alpha v_i}\,\varphi\,\mathrm{d}x,
\end{align*}
其中 $\widetilde{D^\alpha v_i}$ 表示 $D^\alpha v_i$ 到 $\mathbb{R}^n$ 的零延拓。这说明 $\widetilde{D^\alpha v_i}$ 是 $v_i$ 的 $\alpha$ 阶弱导数。由于 $D^\alpha v_i\in L^\infty(W_i)$，其零延拓属于 $L^\infty(\mathbb{R}^n)$。因此 $v_i\in W^{k,\infty}(\mathbb{R}^n)$，且
\begin{align*}
\|v_i\|_{W^{k,\infty}(\mathbb{R}^n)}=\|\zeta_i\overline{u}_i\|_{W^{k,\infty}(W_i)}.
\end{align*}
由于
\begin{align*}
\overline{u}=\sum\limits_{i=0}^{N}v_i
\end{align*}
是有限和，而 $W^{k,\infty}(\mathbb{R}^n)$ 是线性空间，故
\begin{align*}
\overline{u}\in W^{k,\infty}(\mathbb{R}^n).
\end{align*}
再由\eqref{eqnjkNJNS}\eqref{eqnjkNJNS1}式可得
\begin{align}
\|\overline{u}\|_{W^{k,\infty}(\mathbb{R}^{n})}
&\leqslant\sum_{i=0}^{N}\|\zeta_i\overline{u}_i\|_{W^{k,\infty}(W_i)}
\leqslant \left( M_0+NC_0C_2C_3\underset{i=1,\cdots ,N}{\max}\left\{ M_i \right\} \right) \left\| u \right\| _{W^{k,\infty}\left( U \right)}=C\left\| u \right\| _{W^{k,\infty}\left( U \right)},\label{HJKNAJAAAA}
\end{align}
其中$C\triangleq M_0+NC_0C_2C_3\underset{i=1,\cdots ,N}{\max}\left\{ M_i \right\}$.

接下来证明 $\operatorname{spt}\overline{u}\subseteq\bigcup\limits_{i=0}^{N}W_i\subset\subset V$ 以及 $\overline{u}|_U=u$。
对每个 $i=0,1,\cdots,N$，由 $\operatorname{spt}\zeta_i\subseteq W_i$ 得
\begin{align*}
\operatorname{spt}(\zeta_i\overline{u}_i)
\subseteq\operatorname{spt}\zeta_i
\subseteq W_i.
\end{align*}
因此
\begin{align*}
\operatorname{spt}\overline{u}
=\operatorname{spt}\left(\sum\limits_{i=0}^{N}\zeta_i\overline{u}_i\right)
\subseteq\bigcup\limits_{i=0}^{N}\operatorname{spt}(\zeta_i\overline{u}_i)
\subseteq\bigcup\limits_{i=0}^{N}W_i.
\end{align*}
由于每个 $W_i\subset\subset V$，而 $i=0,1,\cdots,N$ 只有有限个，故
\begin{align*}
\overline{\bigcup\limits_{i=0}^{N}W_i}
=\bigcup\limits_{i=0}^{N}\overline{W_i}\subset V.
\end{align*}
于是 $\bigcup\limits_{i=0}^{N}W_i\subset\subset V$，从而
\begin{align*}
\operatorname{spt}\overline{u}\subseteq\bigcup\limits_{i=0}^{N}W_i\subset\subset V.
\end{align*}
再证 $\overline{u}|_U=u$。对几乎处处的 $x\in U$，由 $\operatorname{spt}\zeta_i\subseteq W_i$，有
\begin{align*}
\zeta_i(x)\overline{u}_i(x)=
\begin{cases}
\zeta_i(x)u(x), & x\in W_i\cap U,\\
0=\zeta_i(x)u(x), & x\in U\setminus W_i,
\end{cases}
\end{align*}
其中第一行用了 $\overline{u}_i=u$ 于 $W_i\cap U$，第二行用了 $\zeta_i=0$ 于 $U\setminus W_i$。当 $i=0$ 时 $\overline{u}_0=u$ 于 $U$，上述结论同样成立。于是
\begin{align*}
\overline{u}(x)
=\sum\limits_{i=0}^{N}\zeta_i(x)\overline{u}_i(x)=\sum\limits_{i=0}^{N}\zeta_i(x)u(x)=u(x)\sum\limits_{i=0}^{N}\zeta_i(x)=u(x),
\end{align*}
其中最后一步用了 $\sum\limits_{i=0}^{N}\zeta_i=1$ 于 $\overline{U}\supseteq U$。故
\begin{align*}
\overline{u}|_U=u.
\end{align*}

最后，由反射公式与单位分解的线性性，$u\mapsto E_ku$ 线性；由\eqref{HJKNAJAAAA}式知$E_k$ 有界。故 $p=\infty$ 时定理成立。
\end{proof}

\begin{theorem}[$W^{1,\infty}$的刻画]\label{theorem:ch05-W1infty刻画}
设 $U$是$\mathbb{R}^n$中的有界开集，$\partial U$ 为 $C^{1}$ 类。则 $u:U\to\mathbb{R}$ Lipschitz 连续当且仅当 $u\in W^{1,\infty}(U)$。
\end{theorem}
\begin{remark}
若 $n<p<\infty$，则每个函数 $u\in W^{1,p}$ 都属于 $C^{0,1-\frac{n}{p}}$；但一个以小于 1 的指数H\"older连续的函数未必属于任何Sobolev空间 $W^{1,p}$。
\end{remark}
\begin{proof}

1. 先设 $U=\mathbb{R}^{n}$ 且 $u$ 有紧支集。设 $u\in W^{1,\infty}(\mathbb{R}^{n})$。则 $u_{\varepsilon}\triangleq\eta_{\varepsilon}*u$（$\eta_{\varepsilon}$ 为通常磨光子）光滑，且满足
\begin{align*}
\begin{cases}
u_{\varepsilon}\to u\ \text{一致地，当}\ \varepsilon\to 0,\\
\|Du_{\varepsilon}\|_{L^{\infty}(\mathbb{R}^{n})}\leqslant\|Du\|_{L^{\infty}(\mathbb{R}^{n})}.
\end{cases}
\end{align*}
任取两点 $x,y\in\mathbb{R}^{n}$，$x\neq y$。我们有
\begin{align*}
u_{\varepsilon}(x)-u_{\varepsilon}(y)
=\int_{0}^{1}\frac{\mathrm{d}}{\mathrm{d}t}u_{\varepsilon}(tx+(1-t)y)\,\mathrm{d}t=\int_{0}^{1}Du_{\varepsilon}(tx+(1-t)y)\,\mathrm{d}t\cdot(x-y),
\end{align*}
因而
\begin{align*}
|u_{\varepsilon}(x)-u_{\varepsilon}(y)|
\leqslant\|Du_{\varepsilon}\|_{L^{\infty}(\mathbb{R}^{n})}|x-y|
\leqslant\|Du\|_{L^{\infty}(\mathbb{R}^{n})}|x-y|.
\end{align*}
令 $\varepsilon\to 0$，得
\begin{align*}
|u(x)-u(y)|\leqslant\|Du\|_{L^{\infty}(\mathbb{R}^{n})}|x-y|.
\end{align*}
故 $u$ Lipschitz 连续。

2. 现设 $u:\mathbb{R}^{n}\to\mathbb{R}$ 是 Lipschitz 连续的，记其 Lipschitz 常数为 $\operatorname{Lip}(u)$。固定 $i\in\{1,\cdots,n\}$。对任意 $h\neq 0$，定义差商
\begin{align*}
D_{i}^{-h}u(x)\triangleq\frac{u(x)-u(x-he_{i})}{h}.
\end{align*}
由 Lipschitz 条件，有
\begin{align*}
|D_{i}^{-h}u(x)|\leqslant\operatorname{Lip}(u),\quad \forall x\in\mathbb{R}^{n},
\end{align*}
故
\begin{align*}
\|D_{i}^{-h}u\|_{L^{\infty}(\mathbb{R}^{n})}\leqslant\operatorname{Lip}(u).
\end{align*}
因此差商族 $\{D_{i}^{-h}u\}_{h\neq 0}$ 在 $L^{\infty}(\mathbb{R}^{n})$ 中一致有界，从而在 $L^{2}_{\mathrm{loc}}(\mathbb{R}^{n})$ 中也一致有界。由\cref{Functional Analysis-proposition:Lp对偶与自反性}知$L^{2}_{\mathrm{loc}}(\mathbb{R}^{n})$ 是自反的，再根据\cref{Functional Analysis-proposition:自反空间中弱列紧性等价与有界}, 存在子列$h_{k}\to 0$ 及函数 $v_{i}\in L^{2}_{\mathrm{loc}}(\mathbb{R}^{n})$，使得对每个紧集 $K\subset\mathbb{R}^{n}$，有
\begin{align*}
D_{i}^{-h_{k}}u\rightharpoonup v_{i}\quad\text{在}\ L^{2}(K)\ \text{中弱收敛}.
\end{align*}
对任意 $\phi\in C_c^\infty(\mathbb{R}^n)$，由弱收敛 $D_i^{-h_k}u\rightharpoonup v_i$ 于 $L^2_{\mathrm{loc}}(\mathbb{R}^n)$，有
\begin{align*}
\int_{\mathbb{R}^n}\phi D_i^{-h_k}u\,\mathrm{d}x
\to \int_{\mathbb{R}^n}v_i\,\phi\,\mathrm{d}x.
\end{align*}
另一方面，由 $\|D_i^{-h}u\|_{L^\infty(\mathbb{R}^n)}\leqslant \operatorname{Lip}(u)$ 对一切 $h\neq 0$ 成立，得
\begin{align*}
\left|\int_{\mathbb{R}^n}\phi D_i^{-h_k}u\,\mathrm{d}x\right|
\leqslant \|D_i^{-h_k}u\|_{L^\infty(\mathbb{R}^n)}\int_{\mathbb{R}^n}|\phi|\,\mathrm{d}x
\leqslant \operatorname{Lip}(u)\int_{\mathbb{R}^n}|\phi|\,\mathrm{d}x.
\end{align*}
令 $k\to\infty$，即得
\begin{align*}
\left|\int_{\mathbb{R}^n}v_i\,\phi\,\mathrm{d}x\right|
\leqslant \operatorname{Lip}(u)\int_{\mathbb{R}^n}|\phi|\,\mathrm{d}x.
\end{align*}
由于\cref{Real Analysis-theorem:L无穷空间范数对偶刻画}以及$C_c^\infty(\mathbb{R}^n)$ 在 $L^1(\mathbb{R}^n)$ 中稠密，有
\begin{align*}
\|v_i\|_{L^\infty(\mathbb{R}^n)}
=\sup_{\substack{\phi\in C_c^\infty(\mathbb{R}^n)\\ \|\phi\|_{L^1(\mathbb{R}^n)}= 1}}
\left|\int_{\mathbb{R}^n}v_i\,\phi\,\mathrm{d}x\right|.
\end{align*}
因此综上可得
\begin{align*}
\|v_i\|_{L^\infty(\mathbb{R}^n)}\leqslant \operatorname{Lip}(u).
\end{align*}

任取 $\phi\in C^{\infty}_{c}(\mathbb{R}^{n})$。由\hyperref[theorem:ch05-差商的分部积分公式]{差商的分部积分公式}得
\begin{align*}
\int_{\mathbb{R}^{n}}u(x)D_{i}^{h}\phi(x)\,\mathrm{d}x
=-\int_{\mathbb{R}^{n}}D_{i}^{-h}u(x)\phi(x)\,\mathrm{d}x,
\end{align*}
因为 $\phi$ 光滑且具有紧支集，所以当 $h_{k}\to 0$ 时，$D_{i}^{h_{k}}\phi\to\phi_{x_{i}}$ 在 $L^{\infty}(\mathbb{R}^{n})$ 中成立。于是由\hyperref[Real Analysis-theorem:控制收敛定理]{控制收敛定理}和弱收敛定义得
\begin{align*}
\int_{\mathbb{R}^{n}}u\phi_{x_{i}}\,\mathrm{d}x
=\lim_{k\to\infty}\int_{\mathbb{R}^{n}}uD_{i}^{h_{k}}\phi\,\mathrm{d}x=-\lim_{k\to\infty}\int_{\mathbb{R}^{n}}D_{i}^{-h_{k}}u\,\phi\,\mathrm{d}x
=-\int_{\mathbb{R}^{n}}v_{i}\phi\,\mathrm{d}x.
\end{align*}
上式对一切 $\phi\in C^{\infty}_{c}(\mathbb{R}^{n})$ 成立，故 $v_{i}$ 是 $u$ 关于 $x_{i}$ 的弱导数，即 $v_{i}=u_{x_{i}}$ 在弱意义下成立。由于 $v_{i}\in L^{\infty}(\mathbb{R}^{n})$，所以 $u\in W^{1,\infty}(\mathbb{R}^{n})$。

3. 在 $U$ 有界、$\partial U$ 为 $C^{1}$ 类的一般情形，由\hyperref[theorem:ch05-扩张定理]{扩张定理}，取 $p=\infty$，存在有界线性扩张算子
\begin{align*}
E:W^{1,\infty}(U)\to W^{1,\infty}(\mathbb{R}^{n}),
\end{align*}
使得对每个 $u\in W^{1,\infty}(U)$，有 $Eu=u$ a.e. 于 $U$，且 $Eu$ 在 $\mathbb{R}^{n}$ 中具有紧支集，并且
\begin{align*}
\|Eu\|_{W^{1,\infty}(\mathbb{R}^{n})}\leqslant C\|u\|_{W^{1,\infty}(U)}.
\end{align*}

若 $u\in W^{1,\infty}(U)$，则 $Eu\in W^{1,\infty}(\mathbb{R}^{n})$ 且具有紧支集。由第 1 步的结论，$Eu$ 是 Lipschitz 连续的。因此 $u=Eu|_{U}$ 也是 Lipschitz 连续的。

反之，若 $u:U\to\mathbb{R}$ 是 Lipschitz 连续的，由 \hyperref[theorem:实值-Kirszbraun扩张定理]{Kirszbraun扩张定理}，$u$ 可以扩张为 $\mathbb{R}^{n}$ 上的 Lipschitz 函数 $\overline{u}$，且 $\operatorname{Lip}(\overline{u})=\operatorname{Lip}(u)$。由第 2 步知 $\overline{u}\in W^{1,\infty}(\mathbb{R}^{n})$。限制在 $U$ 上得 $u=\overline{u}|_{U}\in W^{1,\infty}(U)$。

综上，$u$ Lipschitz 连续当且仅当 $u\in W^{1,\infty}(U)$。
\end{proof}



\end{document}
